\chapter{Differentiability, orthogonality, and the variational equation}
\index{differentiability!non-Archimedean volume}
\index{differentiability!of the Monge--Ampere energy}
\index{orthogonality!BFJ}
\index{variational method!BFJ}
\index{Euler--Lagrange equation!non-Archimedean}
\label{chap:bfj-differentiability}

There are two different differentiability statements in this chapter, and the
whole architecture of the variational method depends on not confusing them.
Along a segment contained in $\mathcal E^1$, energy differentiates directly
to the Monge--Amp\`ere measure; this is a theorem about a concave function of
one real variable, and it costs nothing beyond the energy calculus of
\Cref{chap:bfj-energy}.  An arbitrary continuous perturbation does not
preserve psh positivity, so it must first be projected by the psh envelope.
Differentiability after that projection is a much harder statement, and it
turns out to be \emph{equivalent} to orthogonality of the envelope.  This
second statement is the one that turns a variational maximizer into a
solution of the Monge--Amp\`ere equation.

The two statements are separated here by an intermediate object which is the
real organizing principle of the chapter: the \emph{superdifferential} of the
concave functional $E_\omega\circ P_\omega$ on $C^0(X^{\An})$.  Soft
arguments---concavity, monotonicity, translation equivariance, and the
Hahn--Banach theorem---show that this superdifferential is a nonempty,
weak-$*$ compact, convex set of positive measures of total mass $V$
carried by the contact set of the obstacle, using envelope continuity for
that last assertion.  Orthogonality is precisely the
assertion that the Monge--Amp\`ere measure of the envelope is one of them,
and Gateaux differentiability is precisely the assertion that there is only
one.  Once this is said, the logical order of the chapter is forced, and the
only deep input is orthogonality itself.

\paragraph{Standing hypotheses.}
We retain the original BFJ geometric hypotheses fixed at the beginning of
\Cref{chap:bfj-model-forms} and used throughout
\Cref{chap:bfj-psh,chap:bfj-compactness,chap:bfj-envelopes,chap:bfj-ma,chap:bfj-capacity,chap:bfj-energy}:
$K$ is complete and discretely valued with residue characteristic zero and
$v(\pi)=1$; $X$ is smooth, connected and projective of dimension $n$
over $K$; and $\omega$ is a closed model $(1,1)$-form which is
semipositive and whose de Rham class $\{\omega\}$ is ample.  As before,
\[
                         V=(\{\omega\}^n)>0 .
\]
\Cref{chap:bfj-envelopes} writes $\theta$ for the reference form; the two
letters denote the same kind of object, and every result quoted from that
chapter is applied with $\theta=\omega$.  The bounded and finite-energy
Monge--Amp\`ere measures used in the variational statements have total mass
$V$, never a tacit probability normalization; see
\eqref{eq:appJ-cl-total-mass} and \eqref{eq:appJ-model-ma}.

The extra hypothesis needed in the original orthogonality proof, and its
later removal, are stated separately in
\Cref{thm:orthogonality-provenance}.

\paragraph{What the reader should learn.}
By the end of the chapter the reader should be able to: compute the two
one-sided derivatives of $E_\omega$ along a segment in $\mathcal E^1$ and
know why they exist; state the two-sided energy sandwich for an arbitrary
pair of finite-energy potentials, with no ordering hypothesis; recognize
orthogonality in each of its five equivalent forms, including the Legendre
form; reproduce the two-line derivation of the derivative of
$E_\omega\circ P_\omega$ from orthogonality, and the converse derivation of
orthogonality from that derivative; write down the Euler--Lagrange equation
of the variational functional and verify it on a graph; and---most
important---say, at every single step, whether continuity of envelopes has
been used, and if so in which of its forms.

\paragraph{Prerequisites.}
The energy calculus of \Cref{chap:bfj-energy}, in particular the cocycle
\eqref{eq:energy-cocycle}, the first variation
\eqref{eq:energy-first-variation}, the extension
\Cref{thm:extended-energy-basic} and the monotone continuity
\Cref{prop:e1-monotone-continuity}; the bounded Monge--Amp\`ere operator
\Cref{thm:bounded-ma-extension} and the locality statements of
\Cref{chap:bfj-capacity}; the envelope calculus
\Cref{prop:ch23-formal-envelope-calculus} and the continuity theorem
\Cref{thm:bfj-envelope-regularity} of \Cref{chap:bfj-envelopes}, together
with the model-level orthogonality results
\Cref{thm:ch23-model-orthogonality-curve} and
\Cref{thm:ch23-convex-envelope-orthogonality} proved there.  From
\Cref{chap:topological-tools} we use the portmanteau theorem
\Cref{prop:portmanteau}, the joint convergence lemma
\Cref{lem:joint-weak-uniform}, the monotone integral statement
\Cref{cor:usc-net-monotone-integral}, the Riesz interface
\Cref{thm:riesz-interface}, and the concave-conjugate formalism of
\Cref{subsec:legendre-finite}, in particular
\Cref{prop:concave-biconjugation} and \eqref{eq:dual-energy-abstract}.
Convex-analytic facts about the real Monge--Amp\`ere operator are in
\Cref{sec:appD-real-ma}.

\section{Orientation}
\label{sec:ch27-orientation}

\subsection{Four functionals, and which variable moves}
\label{sec:ch27-four-functionals}

Four different functions of a real variable occur in this chapter, and they
are constantly confused in informal accounts.  We name them once.

\begin{enumerate}
\item
The \emph{energy along a psh segment}
\[
 h(t)=E_\omega\bigl((1-t)\varphi+t\psi\bigr),
 \qquad \varphi,\psi\in\mathcal E^1(X,\omega).
\]
Here the moving object is already psh at every time.  The function $h$ is
concave, and for bounded potentials it is a \emph{polynomial} of degree at
most $n+1$.  Its derivative is \Cref{thm:e1-energy-differentiability}.
\item
The \emph{energy of a projected continuous perturbation}
\[
 F(f+tg)=E_\omega\bigl(P_\omega(f+tg)\bigr),
 \qquad f,g\in C^0(X^{\An}).
\]
Here the moving object $f+tg$ is not psh, and the projection $P_\omega$
intervenes.  The function $t\mapsto F(f+tg)$ is concave and
$V\lVert g\rVert_\infty$-Lipschitz, but it need not be a polynomial, and
its differentiability is \Cref{thm:energy-envelope-differentiability}.
\item
The \emph{linearized functional}
\[
 q(t)=\int_{X^{\An}}P_\omega(f+tg)\,\rd\mu,
 \qquad \mu=\MA_\omega(P_\omega f),
\]
in which the energy has been replaced by integration against a fixed
measure.  It is concave for the same reason as
item~(2), and the assertion that its derivative
agrees with that of $F$ is \eqref{eq:linearized-envelope-energy}.
\item
The \emph{variational functional}
$F_\mu(\varphi)=E_\omega(\varphi)-\int\varphi\,\rd\mu$ of
\eqref{eq:bfj-variational-functional}, whose stationarity is the
Euler--Lagrange equation \eqref{eq:euler-lagrange-equation}.
\end{enumerate}

The passage from item~(1) to
item~(2) is not a technical refinement.  A segment
between two psh functions stays in the convex psh class, and its energy
derivative is a mixed-product integral, an intersection number at model
level.  A segment between two \emph{obstacles} moves the
constraint set itself, and the derivative of the value of the associated
obstacle problem is a first-order sensitivity statement about a
constrained optimization.  Orthogonality is exactly the complementary
slackness that makes that sensitivity computable.

\subsection{The dependency order}
\label{sec:ch27-dependency}

\index{dependency!differentiability and orthogonality}
This is the most delicate point in the chapter, and we state it before any
proof.  Differentiability of energy is used to derive orthogonality;
orthogonality is used to prove differentiability of $E_\omega\circ
P_\omega$; differentiability of $E_\omega\circ P_\omega$ is used in the
existence proof of \Cref{chap:bfj-finite-energy-measures}; and continuity of
envelopes enters as a hypothesis in some routes and as a conclusion in
others.  Written like that the situation looks circular.  It is not, because
the three occurrences of the word ``differentiability'' refer to three
different statements.  The following list is the order actually adopted
here, and no step uses a later one.

\begin{enumerate}
\item[\textup{(D1)}]
\emph{Skeletal equicontinuity, compactness, and the upper-envelope
calculus.}
\Cref{thm:bfj-skeletal-equicontinuity}, \Cref{thm:bfj-psh-compactness},
\Cref{thm:ch23-upper-envelope}.  Input: residue characteristic zero only
through cofinality of SNC models \Cref{prop:snc-cofinal}.  No envelope
continuity, no measures, no energy.
\item[\textup{(D2)}]
\emph{Continuity of envelopes of continuous obstacles.}
\Cref{thm:bfj-envelope-regularity}.  Input: asymptotic multiplier ideals
with subadditivity and uniform global generation, that is, residue
characteristic zero in an essential way.  \Cref{sec:ch23-imported} verifies
that its proof uses nothing from
\Cref{chap:bfj-ma,chap:bfj-energy,chap:bfj-differentiability}.
\item[\textup{(D3)}]
\emph{Decreasing model regularization.}
\Cref{thm:bfj-decreasing-regularization}, deduced from \textup{(D2)} in
\Cref{cor:ch23-decreasing-regularization-pointer}.
\item[\textup{(D4)}]
\emph{The bounded Monge--Amp\`ere operator and its locality.}
\Cref{thm:bounded-ma-extension}, \Cref{thm:bounded-maximum-principle},
\Cref{cor:bounded-comparison-principle}.  These use \textup{(D3)} to produce
the decreasing model approximants along which the operator is defined.
\item[\textup{(D5)}]
\emph{Energy, the class $\mathcal E^1$, and the nonpluripolar measure.}
\Cref{def:model-energy}, \Cref{prop:energy-cocycle},
\Cref{thm:extended-energy-basic}, \Cref{def:nonpluripolar-ma},
\Cref{prop:e1-integrability-tail}, \Cref{prop:e1-monotone-continuity}.  These
use \textup{(D4)}.
\item[\textup{(D6)}]
\emph{Differentiability of energy along segments in $\mathcal E^1$, and
the two-sided energy sandwich.}
\Cref{thm:e1-energy-differentiability} and
\Cref{cor:ch27-two-sided-sandwich} below.  These use \textup{(D5)} and
nothing else.  \emph{In particular they do not use orthogonality.}
\item[\textup{(D7)}]
\emph{The superdifferential of $E_\omega\circ P_\omega$.}
\Cref{prop:ch27-F-elementary}, \Cref{thm:ch27-superdifferential},
\Cref{prop:ch27-supergradients-on-contact}.  These use \textup{(D5)},
\textup{(D6)}, the envelope calculus, and---for the statement that every
supergradient is carried by the contact set---continuity of envelopes
\textup{(D2)}, because $P_\omega f$ must itself be an admissible
continuous obstacle.  \emph{They do not use orthogonality.}
\item[\textup{(D8)}]
\emph{Orthogonality.}  \Cref{def:orthogonality-property}.  This is
\emph{not} proved here for $n\geq2$.  It is available through two
historically distinct routes recorded in
\Cref{thm:orthogonality-provenance}, both of which lie outside the chain
\textup{(D1)}--\textup{(D7)} and neither of which is a formal consequence of
it.  In dimension one it is proved here under a model-level hypothesis satisfied
by strictly semistable models, \Cref{prop:ch27-curve-orthogonality} and
\Cref{cor:ch27-curve-orthogonality-global} --- see
\Cref{rem:ch27-curve-base-change} for what base change costs --- and in the
invariant toric category it is proved in every dimension for the
\emph{invariant} envelope, \Cref{thm:ch27-toric-orthogonality} with the
scope restrictions of \Cref{rem:ch27-toric-scope}; building on
\Cref{thm:ch23-model-orthogonality-curve} and
\Cref{thm:ch23-convex-envelope-orthogonality}.
\item[\textup{(D9)}]
\emph{Differentiability of $E_\omega\circ P_\omega$.}
\Cref{thm:ch27-model-direction-differentiability} for model directions and
\Cref{thm:energy-envelope-differentiability} for continuous directions.
These use \textup{(D6)}, \textup{(D7)} and \textup{(D8)}.  Conversely
\Cref{prop:differentiability-iff-orthogonality} recovers \textup{(D8)} from
\textup{(D9)}, using \textup{(D2)} once more; so \textup{(D8)} and
\textup{(D9)} are equivalent, and this equivalence is not a circularity but
the content of \Cref{cor:ch27-orthogonality-iff-differentiability}.
\item[\textup{(D10)}]
\emph{Finite-energy projected variations and the Euler--Lagrange equation.}
\Cref{cor:e1-projected-variation} and \Cref{prop:bfj-euler-lagrange}.  These
use \textup{(D9)} for continuous obstacles and \textup{(D5)}--\textup{(D6)}
for the passage to $\mathcal E^1$.
\item[\textup{(D11)}]
\emph{Existence of a maximizer.}  This is not in this chapter.  It requires
compactness and coercivity and is
\Cref{thm:bfj-supported-cy} in \Cref{chap:bfj-finite-energy-measures}.
\end{enumerate}

In diagram form, with $\Longrightarrow$ meaning ``is used in the proof
of'':
\begin{equation}\label{eq:ch27-dependency-diagram}
 \begin{aligned}
  \text{multiplier ideals}
   &\Longrightarrow \textup{(D2)}\ \text{envelope continuity},\\[2pt]
  \textup{(D2)}
   &\Longrightarrow \textup{(D3)}\Longrightarrow \textup{(D4)}
    \Longrightarrow \textup{(D5)}\Longrightarrow \textup{(D6)},\\[2pt]
  \textup{(D2)}+\textup{(D6)}
   &\Longrightarrow \textup{(D7)}\ \text{superdifferential calculus},\\[2pt]
  \left.
   \begin{gathered}
    \text{asymptotic orthogonality}\\[-2pt]
    \text{together with algebraizability}\\[-2pt]
    \text{or differentiability of NA volumes}\\[-2pt]
    \text{together with \textup{(D2)}}
   \end{gathered}
  \right\}
   &\Longrightarrow \textup{(D8)}\ \text{orthogonality},\\[2pt]
  \textup{(D6)}+\textup{(D7)}+\textup{(D8)}
   &\Longrightarrow \textup{(D9)}\\[-1mm]
   &\Longrightarrow \textup{(D10)}\Longrightarrow \textup{(D11)} .
 \end{aligned}
\end{equation}

The only arrow that could conceivably close a loop is the one from
\textup{(D9)} back to \textup{(D8)}, and it does not, because
\Cref{prop:differentiability-iff-orthogonality} is an implication between two
statements neither of which is used in the proof of the other: the forward
implication assumes orthogonality \emph{as a hypothesis} and produces a
derivative, while the backward implication assumes the derivative
\emph{as a hypothesis} and produces orthogonality.  Neither is used to prove
\textup{(D1)}--\textup{(D7)}.  This is the same separation that
\eqref{eq:ch23-dependency-chain} enforces on the envelope side.

\begin{center}
\fbox{\begin{minipage}{0.91\textwidth}
\textbf{The one forbidden shortcut.}
It is tempting to argue: $P_\omega f$ is maximal among $\omega$-psh
functions below $f$, hence perturbing it upward must fail, hence
$\MA_\omega(P_\omega f)$ is carried by the contact set.  This is false
reasoning for $n\geq2$, for the reason given in
\Cref{rem:model-slackness-heuristic} and analysed in
\Cref{rem:ch23-model-slackness-fails}: the perturbation need not remain
semipositive, and the coefficient of the model measure at a divisorial point
is a degree-$n$ polynomial in the vertex data, not an affine function of
it.  No argument of this chapter uses that shortcut.
\end{minipage}}
\end{center}

\subsection{Hypothesis signatures}
\label{sec:ch27-signatures}

\begin{center}
\small
\begin{tabular}{p{.27\textwidth}|p{.31\textwidth}|p{.32\textwidth}}
\toprule
\textbf{Statement}&\textbf{Uses}&\textbf{Does not use}\\
\midrule
\Cref{thm:e1-energy-differentiability}
 (energy along an $\mathcal E^1$ segment)
 &the energy cocycle; monotone continuity; the mass tail
  \eqref{eq:e1-mass-tail}
 &orthogonality; envelope continuity beyond \textup{(D2)}$\to$\textup{(D5)}\\
\Cref{thm:ch27-superdifferential}
 (superdifferential is a compact convex set of measures)
 &concavity, monotonicity, translation, Hahn--Banach
 &orthogonality; any intersection theory\\
\Cref{prop:ch27-supergradients-on-contact}
 (supergradients live on the contact set)
 &\textup{(D2)}: $P_\omega f\in C^0$
 &orthogonality\\
\Cref{thm:ch27-model-direction-differentiability}
 (model directions)
 &orthogonality at $f$ and at $f+tg$; the two-sided sandwich
 &density of model functions; approximation\\
\Cref{thm:energy-envelope-differentiability}
 (continuous directions)
 &the model case; density; uniform contraction
 &new geometric input\\
\Cref{prop:differentiability-iff-orthogonality}
 (converse)
 &\textup{(D2)}: $P_\omega f\in C^0$ is an admissible direction
 &intersection theory\\
\Cref{prop:bfj-euler-lagrange}
 (Euler--Lagrange)
 &\Cref{cor:e1-projected-variation}; density of model functions
 &existence of a maximizer\\
\bottomrule
\end{tabular}
\end{center}

Two rows illustrate a recurring direct use of envelope continuity.
In the contact-set argument and the converse to differentiability,
$P_\omega f$ must be fed back into a statement whose input is continuous.
The elementary identities for $F$, the uniform-limit orthogonality argument
and the conjugate identification have related direct uses.
\Cref{sec:ch27-where-continuity} records all five groups
\textup{(EC2)}--\textup{(EC6)}, separately from the upstream construction
\textup{(D2)}$\to$\textup{(D3)}$\to$\textup{(D4)}.

\subsection{What is proved here and what is imported}
\label{sec:ch27-imported}

Everything below is proved in full from
\Cref{chap:bfj-model-forms,chap:bfj-psh,chap:bfj-compactness,chap:bfj-envelopes,chap:bfj-ma,chap:bfj-capacity,chap:bfj-energy}
and from exactly one imported package, namely orthogonality in dimension
$\geq2$, which is stated with its two provenances in
\Cref{thm:orthogonality-provenance} and is never proved here.  All
statements about positive residue characteristic are quoted with their
sources and hypotheses in \Cref{sec:ch27-where-continuity}; none of them is
used in a proof.

Nothing in this chapter is used in \Cref{chap:bfj-envelopes}, which is what
makes \eqref{eq:ch23-dependency-chain} and
\eqref{eq:ch27-dependency-diagram} compatible.

\subsection{Plan}
\label{sec:ch27-plan}

\Cref{sec:ch27-directional} computes the one-sided directional derivatives
of $E_\omega$ inside the psh cone, first as an exact polynomial identity
for bounded potentials, then in $\mathcal E^1$; it ends with the two-sided
energy sandwich, which is the workhorse of the rest of the chapter.
\Cref{sec:ch27-superdifferential} introduces $F=E_\omega\circ P_\omega$
and computes its superdifferential as far as soft arguments allow.
\Cref{sec:ch27-orthogonality} defines orthogonality, gives its five
equivalent forms, records its provenance, and proves it in the two cases in
which it is provable here.  \Cref{sec:ch27-model-directions} proves
differentiability along model directions and
\Cref{sec:ch27-continuous-directions} reduces continuous directions to model
ones.  \Cref{sec:ch27-e1-variations} extends the derivative to
finite-energy obstacles, and \Cref{sec:ch27-euler-lagrange} derives the
Euler--Lagrange equation.  \Cref{sec:ch27-where-continuity} is the audit of
envelope continuity, including what is known in positive residue
characteristic.  \Cref{sec:ch27-worked} contains the computations, carried
to numbers, and \Cref{sec:ch27-dictionary} the chapter dictionary.

\section{Directional derivatives of energy inside the psh cone}
\label{sec:ch27-directional}

\subsection{One-sided derivatives of a concave function}
\label{sec:ch27-concave-calculus}

Everything in this section rests on one elementary fact, which we isolate
because it is used at least eight times below and because it is the only
reason one-sided derivatives exist at all in situations where nothing else
is known.

\begin{lemma}
\label{lem:ch27-concave-quotients}
Let $I\subset\mathbf R$ be an interval and let $h:I\to\mathbf R$ be
concave.  Fix $t_0$ in $I$.
\begin{enumerate}
\item The difference quotient
\[
 \Delta(s):=\frac{h(t_0+s)-h(t_0)}{s}
\]
is nonincreasing in $s$ on $\{s\neq0: t_0+s\in I\}$.
\item If $t_0$ is not the right endpoint of $I$, the right derivative
$h'(t_0+)=\lim_{s\downarrow0}\Delta(s)=\sup_{s>0}\Delta(s)$ exists in
$\mathbf R\cup\{+\infty\}$; if $t_0$ is interior, it is finite, and
similarly $h'(t_0-)=\inf_{s<0}\Delta(s)$ exists and
\begin{equation}\label{eq:ch27-one-sided-order}
 h'(t_0-)\ \geq\ h'(t_0+),
 \qquad
 h(t)\leq h(t_0)+h'(t_0+)(t-t_0)\ \ (t\geq t_0).
\end{equation}
\item At an interior point $t_0$, $h$ is differentiable if and only if
$h'(t_0-)=h'(t_0+)$; the exceptional interior set is countable.
On the interior, $t\mapsto h'(t+)$ is nonincreasing.
\item If $h$ is concave on $[a,b]$ and finite, then for all
$a\leq t_1<t_2\leq b$,
\begin{equation}\label{eq:ch27-endpoint-sandwich}
 h'(t_2-)\ \leq\ \frac{h(t_2)-h(t_1)}{t_2-t_1}\ \leq\ h'(t_1+).
\end{equation}
\end{enumerate}
\end{lemma}

\begin{proof}
(1) is the three-chord inequality: for $s_1<s_2$ with $t_0+s_i\in I$ and
$s_i\neq0$, concavity of $h$ on the three points
$t_0,t_0+s_1,t_0+s_2$ gives $\Delta(s_1)\geq\Delta(s_2)$.  One checks
the three sign cases separately; each is the statement that the slope of a
chord of a concave function decreases as the chord moves to the right.
Assertion (2) follows because a monotone function has one-sided limits, and
finiteness at an interior point holds because $\Delta(s)$ for $s<0$
bounds $\Delta(s')$ for $s'>0$ from above by (1).  The first inequality
in \eqref{eq:ch27-one-sided-order} is again (1); the second is the
definition of the supremum in (2).  Assertion (3) is the standard
monotonicity argument: $h'(\cdot+)$ is nonincreasing by (1), a
nonincreasing function has at most countably many jumps, and
$h'(t_0-)-h'(t_0+)$ is the size of the jump.  Assertion (4) is
\eqref{eq:ch27-one-sided-order} applied at $t_1$ and at $t_2$.
\end{proof}

\begin{lemma}
\label{lem:ch27-concave-nets}
Let $I=[a,b]$ and let $(h_\alpha)$ be a net of concave functions
$I\to\mathbf R$ converging pointwise to a finite function $h$.  Then:
\begin{enumerate}
\item $h$ is concave;
\item if the $h_\alpha$ are uniformly $L$-Lipschitz, the convergence is
uniform on $I$;
\item if $(k_\alpha)$ is a net of nonincreasing functions
$I\to[-M,M]$ converging pointwise to $k$, then for every
$t\in I$
\begin{equation}\label{eq:ch27-helly-integral}
 \int_a^tk_\alpha(s)\,\rd s\longrightarrow\int_a^tk(s)\,\rd s .
\end{equation}
\end{enumerate}
\end{lemma}

\begin{proof}
(1) Concavity is a pointwise inequality among three values, preserved by
pointwise limits.

(2) If $L=0$, each function is constant and pointwise convergence at one
point is uniform convergence.  Suppose $L>0$.  Fix $\varepsilon>0$ and
choose a finite $\varepsilon/(3L)$-net
$t_1,\dots,t_N$ of $I$.  By pointwise convergence and directedness there
is an index $\alpha_0$ beyond which $|h_\alpha(t_i)-h(t_i)|<\varepsilon/3$
for every $i$.  The limit $h$ is $L$-Lipschitz as a pointwise limit of
$L$-Lipschitz functions, so for arbitrary $t\in I$, choosing $t_i$
with $|t-t_i|\leq\varepsilon/(3L)$,
\[
 |h_\alpha(t)-h(t)|
 \leq|h_\alpha(t)-h_\alpha(t_i)|+|h_\alpha(t_i)-h(t_i)|+|h(t_i)-h(t)|
 <\varepsilon .
\]

(3) The pointwise limit $k$ of nonincreasing functions is nonincreasing
and bounded by $M$, hence Riemann integrable on $[a,t]$.  Let
$a=s_0<s_1<\cdots<s_N=t$ be a partition of mesh $\delta$.  Monotonicity
gives, for every index $\alpha$,
\[
 \sum_{j=1}^N(s_j-s_{j-1})k_\alpha(s_j)
 \ \leq\ \int_a^tk_\alpha
 \ \leq\ \sum_{j=1}^N(s_j-s_{j-1})k_\alpha(s_{j-1}),
\]
and the same for $k$.  The two extreme sums are finite linear
combinations of values, so they converge to the corresponding sums for
$k$; and those two sums differ by at most
$\delta\,(k(a)-k(t))\leq2M\delta$.  Hence
\[
 \limsup_\alpha\int_a^tk_\alpha-\liminf_\alpha\int_a^tk_\alpha\leq4M\delta,
 \qquad
 \Bigl|\limsup_\alpha\int_a^tk_\alpha-\int_a^tk\Bigr|\leq4M\delta .
\]
Let $\delta\downarrow0$.
\end{proof}

\begin{remark}[Why a net and not a sequence]
\label{rem:ch27-nets-here}
Decreasing model regularization \Cref{thm:bfj-decreasing-regularization}
provides a \emph{net}, whose index set need not be countably cofinal
(\Cref{rem:ch23-selection-scope}, \Cref{ex:no-cofinal-sequence}).  Ordinary
dominated convergence is therefore unavailable, and
\Cref{lem:ch27-concave-nets}\,(3) is the substitute: it replaces domination
by monotonicity in the \emph{integration} variable, which is exactly what
concavity of the energy supplies.  No countability of the ground field or
of the model system is assumed.  Sequences chosen for particular potentials
or uniform approximations remain available; see
\Cref{cor:psh-sequential-regularization}.
\end{remark}

\subsection{The energy polynomial on a bounded segment}
\label{sec:ch27-energy-polynomial}

We first compute everything exactly.  Let $\varphi,\psi$ be bounded
$\omega$-psh functions and put
\[
 \varphi_t:=(1-t)\varphi+t\psi=\varphi+t(\psi-\varphi),
 \qquad 0\leq t\leq1,
\]
which is bounded $\omega$-psh for every such $t$ because
$\PSH(X,\omega)$ is convex.  Write
\begin{equation}\label{eq:ch27-mixed-coefficients}
 a_k:=\int_{X^{\An}}(\psi-\varphi)\,
       \omega_\varphi^{\,n-k}\wedge\omega_\psi^{\,k},
 \qquad k=0,1,\dots,n,
\end{equation}
with $\omega_\varphi=\omega+\ddc\varphi$ as in \Cref{chap:bfj-energy}.
Each $a_k$ is a finite real number: the mixed measures are positive of
mass $V$ by \Cref{thm:bounded-ma-extension}, and $\psi-\varphi$ is a
bounded Borel function.

\begin{proposition}
\label{prop:ch27-energy-polynomial}
With the above notation, the function
$h(t)=E_\omega(\varphi_t)$ on $[0,1]$ is a polynomial of degree at most
$n+1$.  Explicitly,
\begin{equation}\label{eq:ch27-energy-derivative-binomial}
 h'(t)=\int_{X^{\An}}(\psi-\varphi)\,\MA_\omega(\varphi_t)
 =\sum_{k=0}^n\binom nk(1-t)^{\,n-k}t^{\,k}a_k ,
\end{equation}
and consequently
\begin{equation}\label{eq:ch27-energy-integral-form}
 h(t)=E_\omega(\varphi)+\sum_{k=0}^n\binom nk a_k
      \int_0^t(1-s)^{\,n-k}s^{\,k}\,\rd s .
\end{equation}
In particular
\begin{equation}\label{eq:ch27-endpoint-derivatives}
 h'(0)=a_0=\int(\psi-\varphi)\MA_\omega(\varphi),
 \qquad
 h'(1)=a_n=\int(\psi-\varphi)\MA_\omega(\psi),
\end{equation}
and the cocycle identity \eqref{eq:energy-cocycle} is recovered as
\begin{equation}\label{eq:ch27-cocycle-recovered}
 E_\omega(\psi)-E_\omega(\varphi)=h(1)-h(0)
 =\frac1{n+1}\sum_{k=0}^na_k .
\end{equation}
Finally $h$ is concave, and $h'$ is nonincreasing.
\end{proposition}

\begin{proof}
\emph{Step 1: the derivative formula.}
For model $\varphi,\psi$ the first variation
\eqref{eq:energy-first-variation} gives
$h'(t)=\int(\psi-\varphi)\MA_\omega(\varphi_t)$.  By
\Cref{prop:bounded-energy-identities} the same holds for bounded
$\omega$-psh $\varphi,\psi$, the mixed products being those of
\Cref{thm:bounded-ma-extension}.

\emph{Step 2: expansion of the measure.}
The bounded affine-product identity
\Cref{cor:bounded-ma-multilinearity} gives the following expansion.
Here $\omega_{\varphi_t}=(1-t)\omega_\varphi+t\,\omega_\psi$ is notation
for the entries of completed mixed products; at bounded regularity it is
not an assertion that each entry is a closed model form.  Thus
\[
 \MA_\omega(\varphi_t)=\omega_{\varphi_t}^{\,n}
 =\sum_{k=0}^n\binom nk(1-t)^{\,n-k}t^{\,k}\,
   \omega_\varphi^{\,n-k}\wedge\omega_\psi^{\,k}.
\]
Integrating $\psi-\varphi$ against this identity gives the second equality
in \eqref{eq:ch27-energy-derivative-binomial}.  The right-hand side is a
polynomial in $t$ of degree at most $n$; integrating from $0$ gives
\eqref{eq:ch27-energy-integral-form} and shows that $h$ is a polynomial of
degree at most $n+1$.

\emph{Step 3: the endpoints.}
Setting $t=0$ and $t=1$ in
\eqref{eq:ch27-energy-derivative-binomial} kills all terms but one and gives
\eqref{eq:ch27-endpoint-derivatives}.

\emph{Step 4: the cocycle.}
The Euler Beta integral gives
\[
 \binom nk\int_0^1(1-s)^{\,n-k}s^{\,k}\,\rd s
 =\binom nk\cdot\frac{k!\,(n-k)!}{(n+1)!}
 =\frac1{n+1}
\]
for every $k$, independently of $k$.  Substituting into
\eqref{eq:ch27-energy-integral-form} at $t=1$ yields
\eqref{eq:ch27-cocycle-recovered}, which is
\eqref{eq:energy-cocycle} after renaming the mixed products.

\emph{Step 5: concavity.}
This is \eqref{eq:energy-second-variation} for model functions, extended to
bounded potentials by \Cref{prop:bounded-energy-identities}; the sign is the
relative Hodge index inequality \eqref{eq:model-dirichlet-negativity}, that
is, \Cref{thm:ch20-mixed-hodge}.  A concave differentiable function has
nonincreasing derivative.
\end{proof}

\begin{remark}[The coefficient identity is the whole of first-order energy
theory]
\label{rem:ch27-beta-identity}
The uniformity of the Beta integral in Step~4 is worth pausing on.  It says
that the $n+1$ mixed numbers $a_0,\dots,a_n$ enter the energy
\emph{difference} with equal weight $1/(n+1)$, while they enter the
\emph{derivative at an endpoint} with weight $1$ and $0$.  Every
inequality of this chapter is a comparison between these two weightings,
made possible by the fact that $a_k\geq a_{k+1}$ for every bounded pair.
Indeed their difference is the negative Dirichlet pairing of
$\psi-\varphi$ with itself and the remaining semipositive factors, hence
nonnegative by \Cref{lem:ch26-dirichlet-coefficients}.  No ordering is needed for this
coefficient monotonicity or for the sandwich below.
\end{remark}

\begin{corollary}
\label{cor:ch27-bounded-sandwich}
For bounded $\omega$-psh $\varphi,\psi$, with no ordering assumption,
\begin{equation}\label{eq:ch27-bounded-two-sided}
 \int(\psi-\varphi)\MA_\omega(\psi)
 \ \leq\ E_\omega(\psi)-E_\omega(\varphi)
 \ \leq\ \int(\psi-\varphi)\MA_\omega(\varphi).
\end{equation}
\end{corollary}

\begin{proof}
Apply \eqref{eq:ch27-endpoint-sandwich} to the concave function $h$ of
\Cref{prop:ch27-energy-polynomial} with $t_1=0$, $t_2=1$, and use
\eqref{eq:ch27-endpoint-derivatives}.
\end{proof}

\Cref{cor:ch27-bounded-sandwich} uses the unconditional concavity form of
\eqref{eq:energy-sandwich}, explained in \Cref{rem:ch26-sandwich-scope}.
The ordered case also gives monotonicity of energy, but the variational
proof below needs the sandwich for pairs that need not be comparable.

\subsection{Mutual integrability in \texorpdfstring{$\mathcal E^1$}{E1}}
\label{sec:ch27-mutual-integrability}

Before differentiating in $\mathcal E^1$ we must know that the integrals
occurring in \eqref{eq:ch27-endpoint-derivatives} converge.  This is not
automatic: $\psi$ may be $-\infty$ on a large set, and
$\MA_\omega(\varphi)$ is only known not to charge pluripolar sets.

\begin{lemma}
\label{lem:ch27-mutual-integrability}
Let $\varphi,\psi\in\mathcal E^1(X,\omega)$ with $\varphi\leq0$ and
$\psi\leq0$.  Then $\psi$ is $\MA_\omega(\varphi)$-integrable, and
\begin{equation}\label{eq:ch27-mutual-integrability}
 0\ \geq\ \int_{X^{\An}}\psi\,\MA_\omega(\varphi)
   \ \geq\ E_\omega(\psi)+(n+1)E_\omega(\varphi)\ >\ -\infty .
\end{equation}
Consequently, for arbitrary $\varphi,\psi\in\mathcal E^1(X,\omega)$ the
integral $\int(\psi-\varphi)\MA_\omega(\varphi)$ is a well-defined real
number: the integrand is defined $\MA_\omega(\varphi)$-almost everywhere,
and both of its parts are integrable.
\end{lemma}

\begin{proof}
Write $\varphi_s:=\varphi^{\langle s\rangle}=\max\{\varphi,-s\}$ and
$\psi_s:=\psi^{\langle s\rangle}$, bounded $\omega$-psh functions with
values in $[-s,0]$, decreasing to $\varphi$ and $\psi$ as
$s\to\infty$.

\emph{Step 1: a uniform lower bound.}
Apply \eqref{eq:ch27-bounded-two-sided} to the pair
$(\varphi_s,\psi_s)$:
\[
 E_\omega(\psi_s)-E_\omega(\varphi_s)
 \ \leq\ \int(\psi_s-\varphi_s)\MA_\omega(\varphi_s),
\]
that is,
\[
 \int\psi_s\,\MA_\omega(\varphi_s)
 \ \geq\ E_\omega(\psi_s)-E_\omega(\varphi_s)
        +\int\varphi_s\,\MA_\omega(\varphi_s).
\]
Now $\psi_s\geq\psi$ gives $E_\omega(\psi_s)\geq E_\omega(\psi)$;
$\varphi_s\leq0$ gives $E_\omega(\varphi_s)\leq E_\omega(0)=0$, so
$-E_\omega(\varphi_s)\geq0$; and
\eqref{eq:e1-integrability} applied to the bounded finite-energy function
$\varphi_s$ gives
\[
 \int\varphi_s\,\MA_\omega(\varphi_s)
 =-\int(-\varphi_s)\MA_\omega(\varphi_s)
 \geq(n+1)E_\omega(\varphi_s)\geq(n+1)E_\omega(\varphi),
\]
the last step because $E_\omega(\varphi_s)\geq E_\omega(\varphi)$ and
$n+1>0$.  Combining,
\begin{equation}\label{eq:ch27-uniform-lower-bound}
 \int\psi_s\,\MA_\omega(\varphi_s)
 \ \geq\ E_\omega(\psi)+(n+1)E_\omega(\varphi)
 \qquad\text{for every }s>0 .
\end{equation}

\emph{Step 2: an upper bound in the limit.}
Fix $s_0>0$.  For $s\geq s_0$ we have $\psi_s\leq\psi_{s_0}$, so
\[
 \int\psi_s\,\MA_\omega(\varphi_s)
 \ \leq\ \int\psi_{s_0}\,\MA_\omega(\varphi_s).
\]
The function $\psi_{s_0}$ is bounded and upper semicontinuous, and
$\varphi_s\searrow\varphi$ in $\mathcal E^1$, so
$\MA_\omega(\varphi_s)\to\MA_\omega(\varphi)$ weakly by
\Cref{prop:e1-monotone-continuity}.  All these measures have the same mass
$V$; dividing by $V$ and applying the portmanteau inequality
\Cref{prop:portmanteau} to the upper semicontinuous function
$\psi_{s_0}$,
\[
 \limsup_{s\to\infty}\int\psi_{s_0}\,\MA_\omega(\varphi_s)
 \ \leq\ \int\psi_{s_0}\,\MA_\omega(\varphi).
\]
Hence
$\limsup_s\int\psi_s\MA_\omega(\varphi_s)
 \leq\int\psi_{s_0}\MA_\omega(\varphi)$ for every $s_0$.  Letting
$s_0\to\infty$ and using \Cref{cor:usc-net-monotone-integral} for the
fixed measure $\MA_\omega(\varphi)$ and the decreasing family
$\psi_{s_0}\searrow\psi$,
\begin{equation}\label{eq:ch27-upper-limsup}
 \limsup_{s\to\infty}\int\psi_s\,\MA_\omega(\varphi_s)
 \ \leq\ \int\psi\,\MA_\omega(\varphi),
\end{equation}
the right-hand side being well defined in $[-\infty,0]$ because
$\psi\leq0$.

\emph{Step 3: conclusion.}
Combining \eqref{eq:ch27-uniform-lower-bound} and
\eqref{eq:ch27-upper-limsup},
\[
 \int\psi\,\MA_\omega(\varphi)\ \geq\ E_\omega(\psi)+(n+1)E_\omega(\varphi),
\]
which is finite because $\varphi,\psi\in\mathcal E^1$.  Since
$\psi\leq0$, the integral is at most $0$; this is
\eqref{eq:ch27-mutual-integrability}, and it says exactly that
$\psi\in L^1(\MA_\omega(\varphi))$.

For the last assertion, translate $\varphi$ and $\psi$ by constants so
that both are $\leq0$; the set $\{\psi=-\infty\}$ is pluripolar, hence
$\MA_\omega(\varphi)$-null by \Cref{thm:e1-mixed-products}, and likewise
for $\{\varphi=-\infty\}$.  Thus $\psi-\varphi$ is defined
$\MA_\omega(\varphi)$-almost everywhere, and both $\psi$ and
$\varphi$ are integrable by the case just proved (for $\varphi$, take
$\psi=\varphi$).  Adding constants changes the integral by a finite
multiple of $V$.
\end{proof}

\begin{remark}[The estimate is quantitative and it is used again]
\label{rem:ch27-mutual-quantitative}
Inequality \eqref{eq:ch27-mutual-integrability} is not merely a finiteness
statement: it bounds $\int(-\psi)\MA_\omega(\varphi)$ by
$-E_\omega(\psi)-(n+1)E_\omega(\varphi)$, a quantity controlled on the
compact sets $\mathcal K_C$ of \eqref{eq:normalized-energy-superlevel}.
This uniformity is what makes the coercivity estimates of
\Cref{chap:bfj-finite-energy-measures} possible, and it is the exact
non-Archimedean analogue of the complex statement in
\cite[\S2]{BBGZ13}.  Note also the asymmetry of the two sides: the bound
degrades linearly in the energy of the \emph{measure's} potential with the
factor $n+1$, and only linearly with factor $1$ in the energy of the
integrand.
\end{remark}

\subsection{Differentiability along a segment in
\texorpdfstring{$\mathcal E^1$}{E1}}
\label{sec:ch27-e1-segment}

We can now prove the first of the two differentiability theorems, including
the identification of its finite signed integral at a singular endpoint.

\begin{theorem}\label{thm:e1-energy-differentiability}
Let $\varphi,\psi\in\mathcal E^1(X,\omega)$. The function
\[
              h(t)=E_\omega((1-t)\varphi+t\psi),
              \qquad 0\leq t\leq1,
\]
is a polynomial of degree at most $n+1$, differentiable with one-sided
derivatives at the endpoints, and
\begin{equation}\label{eq:e1-energy-derivative}
 h'(0+)=\int(\psi-\varphi)\MA_\omega(\varphi).
\end{equation}
At every interior $t$, its derivative is
$\int(\psi-\varphi)\MA_\omega((1-t)\varphi+t\psi)$.
These are finite signed integrals, agreeing with the prescribed truncation
limits; no pointwise difference is taken on a common pole set.
\end{theorem}

\begin{proof}
After adding constants we may and do assume $\varphi\leq0$ and
$\psi\leq0$; both sides of \eqref{eq:e1-energy-derivative} then change by
the same finite amount under a further translation, because
$E_\omega(\cdot+c)=E_\omega(\cdot)+cV$ and
$\MA_\omega(\varphi)$ has mass $V$.

\emph{Step 0: $h$ is finite and concave.}
For $0\leq t\leq1$, concavity of $E_\omega$
(\Cref{thm:extended-energy-basic}) gives
\[
 h(t)=E_\omega((1-t)\varphi+t\psi)
 \geq(1-t)E_\omega(\varphi)+tE_\omega(\psi)>-\infty ,
\]
so $(1-t)\varphi+t\psi\in\mathcal E^1$; and $h(t)\leq E_\omega(0)=0$
by monotonicity.  Concavity of $h$ is concavity of $E_\omega$ along the
segment.  By \Cref{lem:ch27-concave-quotients} the one-sided derivatives
$h'(0+)$ and $h'(1-)$ exist in $\mathbf R\cup\{\pm\infty\}$ and
$h$ is differentiable off a countable set.

\emph{Step 1: truncation produces polynomials.}
Put $\varphi_s=\max\{\varphi,-s\}$, $\psi_s=\max\{\psi,-s\}$, and
\[
 h_s(t):=E_\omega\bigl((1-t)\varphi_s+t\psi_s\bigr).
\]
By \Cref{prop:ch27-energy-polynomial}, each $h_s$ is a polynomial of
degree at most $n+1$ on $[0,1]$.  As $s$ increases,
$(1-t)\varphi_s+t\psi_s$ decreases to $(1-t)\varphi+t\psi$, so
\Cref{thm:extended-energy-basic} gives
\[
 h_s(t)\ \searrow\ h(t)
 \qquad\text{for every }t\in[0,1].
\]

\emph{Step 2: coefficientwise convergence.}
Choose $n+2$ distinct points $t_0,\dots,t_{n+1}$ in $[0,1]$.  A
polynomial of degree at most $n+1$ is determined by its values at these
points through Lagrange interpolation, and the map from values to
coefficients is a fixed linear isomorphism
$\mathbf R^{n+2}\to\mathbf R^{n+2}$, hence continuous.  Since
$h_s(t_i)\to h(t_i)$ for each $i$, the coefficient vectors of $h_s$
converge; therefore $h$ is itself a polynomial of degree at most $n+1$
and $h_s\to h$ in the topology of coefficientwise convergence.  In
particular
\begin{equation}\label{eq:ch27-derivative-limit}
 h'(t)=\lim_{s\to\infty}h_s'(t)
 \qquad\text{for every }t\in[0,1],
\end{equation}
and $h$ is differentiable on $[0,1]$, with one-sided derivatives at the
endpoints, as asserted.  (This is where the polynomial structure of
\Cref{prop:ch27-energy-polynomial} does real work: for a general decreasing
limit of concave functions no convergence of derivatives may be expected.)

\emph{Step 3: identification of $h'(0)$, the diagonal term.}
By \eqref{eq:ch27-endpoint-derivatives},
\[
 h_s'(0)=\int(\psi_s-\varphi_s)\MA_\omega(\varphi_s)
 =\int\psi_s\,\MA_\omega(\varphi_s)-\int\varphi_s\,\MA_\omega(\varphi_s).
\]
We treat the two terms separately.  Write
\[
 \nu_s:=\one_{\{\varphi>-s\}}\MA_\omega(\varphi_s),
\]
a positive measure.  By \Cref{def:nonpluripolar-ma} the family
$(\nu_s)$ is increasing and $\nu_s(B)\uparrow\MA_\omega(\varphi)(B)$ for
every Borel $B$; write $\nu=\MA_\omega(\varphi)$, so $\nu_s\leq\nu$.
On $\{\varphi>-s\}$ we have $\varphi_s=\varphi$, so
\[
 \int\varphi_s\,\MA_\omega(\varphi_s)
 =\int\varphi\,\rd\nu_s
   -s\,\MA_\omega(\varphi_s)\bigl(\{\varphi\leq-s\}\bigr).
\]
The second term tends to $0$ by the mass tail \eqref{eq:e1-mass-tail},
which says $\MA_\omega(\varphi_s)(\{\varphi\leq-s\})=o(s^{-1})$.  For the
first term, choose an increasing sequence $s_1<s_2<\cdots\to\infty$ and
decompose
$\nu=\nu_{s_1}+\sum_{j\geq1}(\nu_{s_{j+1}}-\nu_{s_j})$ as a countable sum
of positive measures; countable additivity of the integral of the
nonnegative Borel function $-\varphi$ gives
$\int(-\varphi)\rd\nu_{s_j}\uparrow\int(-\varphi)\rd\nu$.  As the family
is monotone in $s$, the limit along the sequence is the limit along $s$.
Hence
\begin{equation}\label{eq:ch27-diagonal-limit}
 \int\varphi_s\,\MA_\omega(\varphi_s)
 \ \longrightarrow\ \int\varphi\,\MA_\omega(\varphi),
\end{equation}
a finite number by \Cref{lem:ch27-mutual-integrability}.

\emph{Step 4: identification of $h'(0)$, the mixed term.}
Similarly,
\[
 \int\psi_s\,\MA_\omega(\varphi_s)
 =\int\psi_s\,\rd\nu_s
  +\int_{\{\varphi\leq-s\}}\psi_s\,\MA_\omega(\varphi_s),
\]
and the second term is bounded in absolute value by
$s\,\MA_\omega(\varphi_s)(\{\varphi\leq-s\})\to0$, again by
\eqref{eq:e1-mass-tail}.  For the first term we sandwich.  Since
$\psi\leq\psi_s\leq0$ and $\nu_s\leq\nu$,
\[
 \int\psi_s\,\rd\nu_s\ \geq\ \int\psi\,\rd\nu_s\ \geq\ \int\psi\,\rd\nu ,
\]
the second inequality because $\psi\leq0$ and $\nu-\nu_s\geq0$.  In the
other direction, for $s\geq s_0$ we have $\psi_s\leq\psi_{s_0}\leq0$,
hence
\[
 \int\psi_s\,\rd\nu_s\ \leq\ \int\psi_{s_0}\,\rd\nu_s
 \ \xrightarrow[\ s\to\infty\ ]{}\ \int\psi_{s_0}\,\rd\nu ,
\]
the convergence because $\psi_{s_0}$ is a bounded Borel function and
$\nu_s\uparrow\nu$ setwise, so that the same countable decomposition as in
Step~3 applies to $\psi_{s_0}+s_0\geq0$.  Letting $s_0\to\infty$ and
using \Cref{cor:usc-net-monotone-integral} gives
$\limsup_s\int\psi_s\,\rd\nu_s\leq\int\psi\,\rd\nu$.  Combining the two
directions,
\begin{equation}\label{eq:ch27-mixed-limit}
 \int\psi_s\,\MA_\omega(\varphi_s)
 \ \longrightarrow\ \int\psi\,\MA_\omega(\varphi),
\end{equation}
a finite number by \Cref{lem:ch27-mutual-integrability}.

\emph{Step 5: conclusion.}
By \eqref{eq:ch27-derivative-limit} at $t=0$, together with
\eqref{eq:ch27-diagonal-limit} and \eqref{eq:ch27-mixed-limit},
\[
 h'(0+)=\lim_{s\to\infty}h_s'(0)
 =\int\psi\,\MA_\omega(\varphi)-\int\varphi\,\MA_\omega(\varphi)
 =\int(\psi-\varphi)\MA_\omega(\varphi),
\]
where the last equality is legitimate because both integrals are finite and
$\psi-\varphi$ is defined $\MA_\omega(\varphi)$-almost everywhere.  This
proves \eqref{eq:e1-energy-derivative}.  At an interior time $t_0$,
apply the same endpoint formula to the segment from
$\varphi_{t_0}=(1-t_0)\varphi+t_0\psi$ to $\psi$.
Its parameter is $t=t_0+r(1-t_0)$ and its increment is
$(1-t_0)(\psi-\varphi)$ almost everywhere for the endpoint measure.
Dividing by $1-t_0$ gives the stated formula at $t_0$.
By construction the endpoint value is the
limit of the truncated integrals, so the parenthetical clause of the
statement holds; and no pointwise difference has been formed on
$\{\varphi=\psi=-\infty\}$, which is $\MA_\omega(\varphi)$-null.
\end{proof}

\begin{corollary}
\label{cor:ch27-far-endpoint}
In the situation of \Cref{thm:e1-energy-differentiability},
\begin{equation}\label{eq:ch27-far-endpoint}
 h'(1-)=\int(\psi-\varphi)\MA_\omega(\psi).
\end{equation}
\end{corollary}

\begin{proof}
Apply \Cref{thm:e1-energy-differentiability} to the pair $(\psi,\varphi)$
in that order: the resulting function is $k(r)=E_\omega((1-r)\psi+r\varphi)
=h(1-r)$, and $k'(0+)=\int(\varphi-\psi)\MA_\omega(\psi)$.  Since
$k'(0+)=-h'(1-)$, formula \eqref{eq:ch27-far-endpoint} follows.
\end{proof}

\begin{corollary}
\label{cor:ch27-two-sided-sandwich}
For all $u,v\in\mathcal E^1(X,\omega)$, with no ordering assumption,
\begin{equation}\label{eq:ch27-two-sided-sandwich}
 \int(v-u)\,\MA_\omega(v)
 \ \leq\ E_\omega(v)-E_\omega(u)
 \ \leq\ \int(v-u)\,\MA_\omega(u),
\end{equation}
all three quantities being finite.  If equality holds on either side, then
the function $t\mapsto E_\omega((1-t)u+tv)$ is affine on $[0,1]$.
\end{corollary}

\begin{proof}
Finiteness is \Cref{lem:ch27-mutual-integrability}.  The inequalities are
\eqref{eq:ch27-endpoint-sandwich} applied to the concave function
$h(t)=E_\omega((1-t)u+tv)$ with $t_1=0$, $t_2=1$, combined with
\eqref{eq:e1-energy-derivative} and \eqref{eq:ch27-far-endpoint}.

For the last assertion, equality on the right means
$h(1)-h(0)=h'(0+)$.  For a concave function this forces $h$ to be
affine on $[0,1]$: indeed \eqref{eq:ch27-one-sided-order} gives
$h(t)\leq h(0)+h'(0+)t$ for every $t\in[0,1]$, while concavity and the
endpoint values give
$h(t)\geq(1-t)h(0)+th(1)=h(0)+t\,(h(1)-h(0))=h(0)+h'(0+)t$.  Equality on
the left is the same computation read from the endpoint $t=1$.
\end{proof}

\begin{remark}[Equality does not immediately give a constant difference]
\label{rem:ch27-equality-case}
One would like to conclude from equality in
\eqref{eq:ch27-two-sided-sandwich} that $v-u$ is constant.  If $u,v$
are bounded, affineness of $h$ makes its second variation vanish.  For
model endpoints this is the Dirichlet form
\[
 \int(v-u)\,\ddc(v-u)\wedge
             \omega_{u+t(v-u)}^{\,n-1}.
\]
Its weights are relatively nef; the equality clause in
\Cref{thm:ch20-mixed-hodge} requires relatively ample weights.
Consequently that clause cannot simply be applied here.

For singular endpoints, let $h_s(t)=E_\omega((1-t)u_s+tv_s)$.
The interpolation argument gives $h_s''\to h''=0$, but does not say that
each $h_s$ is affine or each truncated Dirichlet form vanishes.
Additional rigidity is needed before passing from this limiting statement
to a constant difference.  We therefore do not assert that conclusion
here.  The continuous Calabi--Yau uniqueness theorem is used with its own
source in \Cref{cor:bfj-original-continuous-solution}.
\end{remark}

\begin{remark}[Which inequality is used where]
\label{rem:ch27-which-inequality}
The right-hand inequality in \eqref{eq:ch27-two-sided-sandwich} is the
supporting-hyperplane inequality: it says that
$\MA_\omega(u)$ is a supergradient of $E_\omega$ at $u$, and it is the
one used to bound the energy increment of a perturbed envelope
\emph{from above}.  The left-hand inequality is the same statement read at
the other endpoint, and it is used to bound the increment \emph{from below}.
Both are needed in \Cref{thm:ch27-model-direction-differentiability}, and it
is precisely the availability of both that makes a two-sided derivative
accessible.  A proof that uses only one of them cannot produce
differentiability; it produces a sandwich
$D^+\leq\int g\,\rd\mu\leq D^-$ which is compatible with a genuine corner,
as \Cref{ex:ch27-corner-without-orthogonality} shows.
\end{remark}

\begin{remark}[No continuous direction is allowed here]
\label{rem:ch27-no-arbitrary-direction}
\Cref{thm:e1-energy-differentiability} does not allow replacing
$\psi-\varphi$ by an arbitrary $f\in C^0(X^{\An})$: the function
$\varphi+tf$ need not be psh. Envelopes provide an admissible replacement, whose projected-energy
derivative requires the orthogonality input studied next.
\end{remark}

\begin{example}[Why the finite-energy endpoint is essential]
\label{ex:ch27-outside-e1}
Let $\varphi\in\PSH(X,\omega)\setminus\mathcal E^1$ and let $\psi=0$.
For
\[
                 h(t):=E_\omega((1-t)\varphi),\qquad0\leq t\leq1,
\]
one has $h(0)=E_\omega(\varphi)=-\infty$.  Thus the endpoint quotient
\[
                    \frac{h(t)-h(0)}{t}
\]
is not a finite real quotient.  If $h(t)$ is finite, it is $+\infty$
in extended-real arithmetic; if $h(t)=-\infty$, it is undefined.
Neither case gives the finite signed variation of
\Cref{thm:e1-energy-differentiability}, under the convention of
\Cref{sec:ch26-extended-real}.  No assertion is made about whether
$(1-t)\varphi$ belongs to $\mathcal E^1$; that depends on the
singularity, and finite energy is not obtained by a formal homogeneity
argument.  Likewise the nonpluripolar measure of $\varphi$ may have mass
strictly below $V$, so the integral appearing in
\eqref{eq:e1-energy-derivative} is not a substitute for the missing
quotient.  This is the reason every statement below carries the hypothesis
$\varphi\in\mathcal E^1$ explicitly, and it is why
\Cref{lem:ch27-mutual-integrability} is proved before, not after,
\Cref{thm:e1-energy-differentiability}.
\end{example}

\section{The envelope energy and its superdifferential}
\label{sec:ch27-superdifferential}

\subsection{The functional \texorpdfstring{$F=E_\omega\circ P_\omega$}{F}}
\label{sec:ch27-F-elementary}

\begin{definition}[Envelope energy]
\label{def:ch27-envelope-energy}
\index{envelope!energy of}
For $f\in C^0(X^{\An})$ set
\begin{equation}\label{eq:ch27-F-definition}
 F(f):=E_\omega\bigl(P_\omega(f)\bigr).
\end{equation}
This is well defined and finite: $P_\omega(f)\in\PSH(X,\omega)$ by
\Cref{cor:ch23-continuous-obstacle-psh}, it is bounded by
\Cref{cor:ch23-continuous-finite}, and $E_\omega$ is finite on bounded
$\omega$-psh functions by \Cref{prop:bounded-energy-identities}.
\end{definition}

Note that \Cref{def:ch27-envelope-energy} does \emph{not} require
$P_\omega(f)$ to be continuous.  Boundedness suffices, and boundedness
comes from the crude two-sided estimate \eqref{eq:ch23-envelope-two-sided},
which uses only ampleness of $\{\omega\}$.  This is worth recording,
because it is the first of several places where one might carelessly invoke
\Cref{thm:bfj-envelope-regularity} and does not need to.

\begin{proposition}
\label{prop:ch27-F-elementary}
For all $f,g\in C^0(X^{\An})$ and $c\in\mathbf R$:
\begin{enumerate}
\item $F(f+c)=F(f)+cV$;
\item if $f\leq g$, then $F(f)\leq F(g)$;
\item $|F(f)-F(g)|\leq V\lVert f-g\rVert_\infty$;
\item $F$ is concave on $C^0(X^{\An})$;
\item $F(P_\omega f)=F(f)$; and if $f$ is $\omega$-psh and continuous,
      then $F(f)=E_\omega(f)$;
\item $F(f)=\sup\{E_\omega(u): u\in\PSH(X,\omega)\cap C^0(X^{\An}),\
      u\leq f\}$, and the supremum is attained at $u=P_\omega(f)$;
\item $|E_\omega(u)-E_\omega(v)|\leq V\lVert u-v\rVert_\infty$ for
      bounded $\omega$-psh $u,v$.
\end{enumerate}
\end{proposition}

\begin{proof}
(1) Combine \eqref{eq:ch23-translation}, $P_\omega(f+c)=P_\omega(f)+c$,
with the translation rule \eqref{eq:bfj-energy-translation} for energy.

(2) Monotonicity of $P_\omega$ is part (1) of
\Cref{prop:ch23-formal-envelope-calculus}, and monotonicity of $E_\omega$
is \Cref{cor:energy-monotone-concave}.

(3) Put $\varepsilon=\lVert f-g\rVert_\infty$.  Then $f\leq g+\varepsilon$,
so by (2) and (1), $F(f)\leq F(g)+\varepsilon V$; exchange $f$ and
$g$.

(4) For $0\leq t\leq1$, joint concavity \eqref{eq:ch23-joint-concavity} of
the envelope, taken with equal reference forms, gives
\[
 P_\omega\bigl(tf+(1-t)g\bigr)\ \geq\ tP_\omega(f)+(1-t)P_\omega(g).
\]
Applying (2)-style monotonicity of $E_\omega$ and then its concavity
(\Cref{cor:energy-monotone-concave}),
\[
 F\bigl(tf+(1-t)g\bigr)
 \geq E_\omega\bigl(tP_\omega f+(1-t)P_\omega g\bigr)
 \geq tF(f)+(1-t)F(g).
\]

(5) Envelope continuity makes $P_\omega(f)$ an admissible argument of
$F$.  It is $\omega$-psh, so
$P_\omega(P_\omega f)=P_\omega f$ by part (6) of
\Cref{prop:ch23-formal-envelope-calculus}; apply $E_\omega$.  The second
assertion is the same statement for $f$ itself.

(6) Any such $u$ satisfies $u\leq P_\omega(f)$ by maximality
(\Cref{cor:ch23-continuous-obstacle-psh}), so
$E_\omega(u)\leq E_\omega(P_\omega f)=F(f)$.  The competitor
$u=P_\omega(f)$ is admissible because it is continuous, by
\Cref{thm:bfj-envelope-regularity}.  Parts (5) and (6) use that theorem
directly, as do \Cref{prop:ch27-supergradients-on-contact} and
\Cref{prop:ch27-conjugate-identification} later in this section.

(7) If $\varepsilon=\lVert u-v\rVert_\infty$ then $u\leq v+\varepsilon$,
so $E_\omega(u)\leq E_\omega(v+\varepsilon)=E_\omega(v)+\varepsilon V$ by
\Cref{cor:energy-monotone-concave} and
\eqref{eq:bfj-energy-translation}; exchange $u$ and $v$.
\end{proof}

\begin{remark}[The three constants]
\label{rem:ch27-three-constants}
Properties (1) and (3) of \Cref{prop:ch27-F-elementary} both involve the
number $V$, and they are not independent: (1) plus (2) imply (3).  The
reader who normalizes to probability measures will see the constant $1$ in
all three places; the ledger \eqref{eq:appJ-energy-translation} and
\eqref{eq:appJ-prob-ma} records the conversion.  Every estimate below
carries the factor $V$ explicitly, so that the mass test of
\eqref{eq:appJ-cl-total-mass} can be applied to it.
\end{remark}

\subsection{Supergradients}
\label{sec:ch27-supergradients}

We now record what concavity alone gives.  Recall from
\Cref{subsec:legendre-finite} that for a concave function on a vector space
the correct dual object is the set of \emph{supergradients}.

\begin{definition}[Superdifferential]
\label{def:ch27-superdifferential}
\index{superdifferential!of the envelope energy}
Let $f\in C^0(X^{\An})$.  A continuous linear functional $\ell$ on the
Banach space $C^0(X^{\An})$ is a \emph{supergradient} of $F$ at $f$ if
\begin{equation}\label{eq:ch27-supergradient-inequality}
 F(g)\ \leq\ F(f)+\ell(g-f)
 \qquad\text{for every }g\in C^0(X^{\An}).
\end{equation}
The set of supergradients is the \emph{superdifferential} $\partial F(f)$.
We also write
\begin{equation}\label{eq:ch27-one-sided-directional}
 D^+F(f;g):=\lim_{t\downarrow0}\frac{F(f+tg)-F(f)}{t},
 \qquad
 D^-F(f;g):=\lim_{t\uparrow0}\frac{F(f+tg)-F(f)}{t},
\end{equation}
which exist in $\mathbf R$ by \Cref{lem:ch27-concave-quotients} applied to
the concave, $V\lVert g\rVert_\infty$-Lipschitz function
$t\mapsto F(f+tg)$, and satisfy $D^+F(f;g)\leq D^-F(f;g)$ and
$D^-F(f;g)=-D^+F(f;-g)$.
\end{definition}

\begin{theorem}
\label{thm:ch27-superdifferential}
Let $f\in C^0(X^{\An})$.  Then:
\begin{enumerate}
\item $\partial F(f)$ is nonempty, convex, and weak-$*$ compact;
\item every $\ell\in\partial F(f)$ is represented by a positive Radon
measure $\mu_\ell$ on $X^{\An}$ with
$\mu_\ell(X^{\An})=V$;
\item for every $g\in C^0(X^{\An})$,
\begin{equation}\label{eq:ch27-minmax-formula}
 D^+F(f;g)=\min_{\ell\in\partial F(f)}\ell(g),
 \qquad
 D^-F(f;g)=\max_{\ell\in\partial F(f)}\ell(g),
\end{equation}
both extrema being attained;
\item $F$ has a two-sided directional derivative at $f$ in the direction
$g$ if and only if all supergradients agree on $g$; and $F$ is Gateaux
differentiable at $f$ if and only if $\partial F(f)$ is a singleton.
\end{enumerate}
\end{theorem}

\begin{proof}
Write $G(g):=D^+F(f;g)$.

\emph{Step 1: $G$ is superlinear.}
Positive homogeneity, $G(\lambda g)=\lambda G(g)$ for $\lambda>0$, is
immediate from \eqref{eq:ch27-one-sided-directional}.  Concavity of $F$
gives, for $g_1,g_2\in C^0$ and $t>0$,
\[
 \frac{F(f+t(g_1+g_2))-F(f)}{t}
 \geq\frac{F(f+2tg_1)-F(f)}{2t}+\frac{F(f+2tg_2)-F(f)}{2t},
\]
because $f+t(g_1+g_2)=\tfrac12(f+2tg_1)+\tfrac12(f+2tg_2)$.  Letting
$t\downarrow0$ gives $G(g_1+g_2)\geq G(g_1)+G(g_2)$.  Moreover
$|G(g)|\leq V\lVert g\rVert_\infty$ by
\Cref{prop:ch27-F-elementary}\,(3).  Finally, concavity gives
\begin{equation}\label{eq:ch27-G-dominates}
 F(f+g)-F(f)\ \leq\ G(g)
 \qquad\text{for every }g,
\end{equation}
since the difference quotient at $t=1$ is at most its limit at
$t\downarrow0$ by \Cref{lem:ch27-concave-quotients}\,(1).

\emph{Step 2: existence of supergradients, and \eqref{eq:ch27-minmax-formula}.}
The function $p:=-G$ is sublinear and $p(g)\leq V\lVert g\rVert_\infty$.
Fix $g_0\in C^0$ and define a linear functional on $\mathbf Rg_0$ by
$\lambda(sg_0):=-sG(g_0)$.  For $s>0$ this equals $p(sg_0)$; for
$s<0$ it equals $|s|G(g_0)$, while
$p(sg_0)=|s|p(-g_0)=-|s|G(-g_0)$, and superadditivity gives
$G(g_0)+G(-g_0)\leq G(0)=0$, so $\lambda\leq p$ on $\mathbf Rg_0$.  By
the Hahn--Banach theorem $\lambda$ extends to a linear functional
$\lambda\leq p$ on $C^0(X^{\An})$; then $\ell:=-\lambda$ satisfies
$\ell\geq G$ everywhere and $\ell(g_0)=G(g_0)$.  From $\ell\geq G$ and
$|G(g)|\leq V\lVert g\rVert_\infty$ we get
$\ell(g)\geq-V\lVert g\rVert_\infty$ and, applying this to $-g$,
$\ell(g)\leq V\lVert g\rVert_\infty$; so $\ell$ is continuous of norm at
most $V$.  By \eqref{eq:ch27-G-dominates}, $\ell\geq G$ implies
\eqref{eq:ch27-supergradient-inequality}, so $\ell\in\partial F(f)$.  This
proves nonemptiness, and it proves that the minimum in
\eqref{eq:ch27-minmax-formula} is attained at $g_0$, since conversely
every $\ell\in\partial F(f)$ satisfies
$\ell(g)\geq D^+F(f;g)$: divide
\eqref{eq:ch27-supergradient-inequality}, applied to $g$ replaced by
$f+tg$, by $t>0$.  The formula for $D^-$ follows from
$D^-F(f;g)=-D^+F(f;-g)$.

\emph{Step 3: positivity and mass.}
Let $\ell\in\partial F(f)$.  If $h\in C^0$ with $h\geq0$, then
$f+th\geq f$ for $t>0$, so $F(f+th)\geq F(f)$ by
\Cref{prop:ch27-F-elementary}\,(2), while
\eqref{eq:ch27-supergradient-inequality} gives
$F(f+th)\leq F(f)+t\ell(h)$.  Hence $\ell(h)\geq0$.  A positive
continuous linear functional on $C^0$ of a compact Hausdorff space is
represented by a positive Radon measure $\mu_\ell$, by the Riesz
representation theorem \Cref{thm:riesz-interface}.  For the mass, apply
\eqref{eq:ch27-supergradient-inequality} with $g=f+c$ and use
\Cref{prop:ch27-F-elementary}\,(1):
$cV\leq c\,\ell(1)$ for every $c\in\mathbf R$, whence
$\ell(1)=V$, that is, $\mu_\ell(X^{\An})=V$.

\emph{Step 4: convexity and compactness.}
Convexity is clear from \eqref{eq:ch27-supergradient-inequality}.  By
Step~3, $\partial F(f)$ is contained in the set of positive measures of
mass $V$, which is weak-$*$ compact: dividing by $V$ identifies it
with the space of Radon probability measures on the compact Hausdorff space
$X^{\An}$, compact by \Cref{prop:measure-compactness}.  It is weak-$*$ closed, because
\eqref{eq:ch27-supergradient-inequality} is a family of closed conditions
indexed by $g$.  A closed subset of a compact set is compact.

\emph{Step 5: differentiability.}
By \eqref{eq:ch27-minmax-formula}, $D^+F(f;g)=D^-F(f;g)$ if and only if
the minimum and maximum of $\ell\mapsto\ell(g)$ over $\partial F(f)$
coincide, that is, if and only if all supergradients agree on $g$.  If
this holds for every $g$, then $\partial F(f)$ is a singleton, because
two distinct measures are separated by some continuous function; and then
$g\mapsto D^+F(f;g)=\ell(g)$ is linear, which is Gateaux differentiability.
\end{proof}

\begin{example}[A one-sided sandwich is not differentiability]
\label{ex:ch27-corner-without-orthogonality}
The formalism of \Cref{thm:ch27-superdifferential} is not vacuous, and the
following two-dimensional caricature shows exactly what the arguments of
\Cref{sec:ch27-model-directions} must avoid.  On $\mathbf R^2$ let
\[
 \Phi(x,y):=\min\{x,y\},
\]
a concave, $1$-Lipschitz, monotone function which satisfies
$\Phi(x+c,y+c)=\Phi(x,y)+c$.  At the origin,
\[
 D^+\Phi(0;(g_1,g_2))=\min\{g_1,g_2\},
 \qquad
 D^-\Phi(0;(g_1,g_2))=\max\{g_1,g_2\},
\]
and $\partial\Phi(0)$ is the segment
$\{(\lambda,1-\lambda):0\leq\lambda\leq1\}$, a compact convex set of
``measures of mass one'' on the two-point space.  The midpoint
$\mu=(\tfrac12,\tfrac12)$ satisfies
\[
 D^+\Phi(0;g)\ \leq\ \langle\mu,g\rangle\ \leq\ D^-\Phi(0;g)
 \qquad\text{for every }g,
\]
and yet $\Phi$ is not differentiable at $0$ in any direction with
$g_1\neq g_2$.  Thus a two-sided \emph{sandwich} of a candidate derivative
between the two one-sided derivatives carries no information at all: it
holds automatically for every element of the superdifferential.  What is
needed is the identification of \emph{one} one-sided derivative with the
candidate, from which the other follows by the sign change
$g\mapsto-g$; or, as in
\Cref{thm:ch27-model-direction-differentiability}, a genuine two-sided
estimate of the difference quotient itself.  This is the content of
\Cref{rem:ch27-which-inequality}.
\end{example}

\subsection{Supergradients live on the contact set}
\label{sec:ch27-contact-support}

The next proposition uses continuity of envelopes directly, as did
parts (5)--(6) of \Cref{prop:ch27-F-elementary}.  Here its role is to know that
$P_\omega f$ may itself be substituted for $g$ in
\eqref{eq:ch27-supergradient-inequality}.

\begin{proposition}
\label{prop:ch27-supergradients-on-contact}
Let $f\in C^0(X^{\An})$ and put
\begin{equation}\label{eq:ch27-contact-set}
 C_f:=\{x\in X^{\An}: P_\omega(f)(x)=f(x)\},
\end{equation}
a closed subset of $X^{\An}$.  Then every $\ell\in\partial F(f)$
satisfies
\begin{equation}\label{eq:ch27-supergradient-contact}
 \int_{X^{\An}}\bigl(f-P_\omega f\bigr)\,\rd\mu_\ell=0,
 \qquad\text{equivalently}\qquad
 \mu_\ell\bigl(X^{\An}\setminus C_f\bigr)=0 .
\end{equation}
\end{proposition}

\begin{proof}
By \Cref{thm:bfj-envelope-regularity} the function $P_\omega f$ is
continuous, so $C_f$ is closed and $P_\omega f$ is an admissible
argument of $F$.  By \Cref{prop:ch27-F-elementary}\,(5),
$F(P_\omega f)=F(f)$.  Substituting $g=P_\omega f$ into
\eqref{eq:ch27-supergradient-inequality},
\[
 F(f)=F(P_\omega f)\ \leq\ F(f)+\ell\bigl(P_\omega f-f\bigr),
\]
so $\ell(f-P_\omega f)\leq0$.  On the other hand $f-P_\omega f\geq0$ and
$\mu_\ell\geq0$ by \Cref{thm:ch27-superdifferential}\,(2), so
$\ell(f-P_\omega f)\geq0$.  Hence the integral in
\eqref{eq:ch27-supergradient-contact} vanishes.  Since the integrand is
continuous and nonnegative and vanishes exactly on $C_f$, vanishing of the
integral says $\mu_\ell(\{f-P_\omega f>\varepsilon\})=0$ for every
$\varepsilon>0$, whence $\mu_\ell(X^{\An}\setminus C_f)=0$.
\end{proof}

\begin{center}
\fbox{\begin{minipage}{0.91\textwidth}
\textbf{What is, and what is not, proved by soft arguments.}
Without any orthogonality hypothesis we now know that
$\partial F(f)$ is a nonempty compact convex set of positive measures of
mass $V$, all of them carried by the contact set $C_f$.  What we do
\emph{not} know is that $\MA_\omega(P_\omega f)$ is one of them, nor that
the set is a singleton.  Orthogonality supplies the first, and
\Cref{thm:ch27-model-direction-differentiability} deduces the second.  The
gap between ``some measure carried by the contact set'' and ``the
Monge--Amp\`ere measure of the envelope'' is the entire difficulty.
\end{minipage}}
\end{center}

\subsection{The concave conjugate, and the link with dual energy}
\label{sec:ch27-conjugate}

We close the soft part of the theory by identifying the concave conjugate of
$F$ with the dual energy of \Cref{chap:bfj-finite-energy-measures}.  In
the notation of \eqref{eq:dual-energy-abstract}, put
\begin{equation}\label{eq:ch27-F-conjugate}
 F^\vee(\mu):=\sup_{f\in C^0(X^{\An})}
   \Bigl(F(f)-\int_{X^{\An}}f\,\rd\mu\Bigr)
 \in[0,+\infty]
\end{equation}
for a positive Radon measure $\mu$ on $X^{\An}$.

\begin{proposition}
\label{prop:ch27-conjugate-identification}
Let $\mu$ be a positive Radon measure on $X^{\An}$.
\begin{enumerate}
\item $F^\vee(\mu)<+\infty$ forces $\mu(X^{\An})=V$.
\item If $\mu(X^{\An})=V$, then
\begin{equation}\label{eq:ch27-conjugates-agree}
 F^\vee(\mu)=E_\omega^\vee(\mu),
\end{equation}
with $E^\vee_\omega$ as in \eqref{eq:dual-energy-measure}.
\item For $f\in C^0(X^{\An})$ and $\mu$ of mass $V$, one has
$\mu\in\partial F(f)$ if and only if the Fenchel equality
\begin{equation}\label{eq:ch27-fenchel-equality}
 F(f)-\int f\,\rd\mu=E^\vee_\omega(\mu)
\end{equation}
holds.
\end{enumerate}
\end{proposition}

\begin{proof}
(1) is \Cref{lem:dual-mass} applied to the concave functional $F$ on
$\mathcal H=C^0(X^{\An})$, which satisfies the translation identity
\eqref{eq:energy-translation} with $M=V$ by
\Cref{prop:ch27-F-elementary}\,(1) and is finite at $f=0$.

(2) For the inequality $F^\vee\leq E^\vee_\omega$, let $f\in C^0$ and put
$u=P_\omega f$.  Since $u\leq f$ and $\mu\geq0$,
\[
 F(f)-\int f\,\rd\mu\ \leq\ E_\omega(u)-\int u\,\rd\mu .
\]
By \Cref{thm:bfj-envelope-regularity} and the approximation criterion
\Cref{cor:ch23-continuity-criterion}, there are $\omega$-psh model
functions $\chi_j$ with $\lVert\chi_j-u\rVert_\infty\to0$.  Energy is
$V$-Lipschitz for the uniform norm on bounded $\omega$-psh functions by
\Cref{prop:ch27-F-elementary}\,(7); hence
$E_\omega(\chi_j)\to E_\omega(u)$, and
$\int\chi_j\,\rd\mu\to\int u\,\rd\mu$ because $\mu$ is finite.  Each
$\chi_j$ is an admissible competitor in \eqref{eq:dual-energy-measure}, so
$E_\omega(u)-\int u\,\rd\mu\leq E^\vee_\omega(\mu)$.

For the reverse inequality, every $\omega$-psh model function $\chi$ is
continuous and satisfies $P_\omega\chi=\chi$, so
$F(\chi)-\int\chi\,\rd\mu=E_\omega(\chi)-\int\chi\,\rd\mu$; taking the
supremum over $\chi$ gives $E^\vee_\omega(\mu)\leq F^\vee(\mu)$.

(3) By definition, $\mu\in\partial F(f)$ means
$F(g)-\int g\,\rd\mu\leq F(f)-\int f\,\rd\mu$ for every $g\in C^0$, that
is, the supremum in \eqref{eq:ch27-F-conjugate} is attained at $f$.  Now
apply (2).
\end{proof}

\begin{remark}[Where this sits in the Legendre picture]
\label{rem:ch27-legendre-picture}
\Cref{prop:ch27-conjugate-identification} is the infinite-dimensional
Fenchel equality: equality is equivalent to the supergradient relation.
Here biconjugation follows as well from the already-proved nonemptiness
of the superdifferential:
\[
 F(f)=\min_{\mu(X^{\An})=V,\ \mu\geq0}
        \left(E_\omega^\vee(\mu)+\int f\,\rd\mu\right).
\]
Indeed the definition of $F^\vee$ makes every expression on the right at
least $F(f)$.  Choose $\mu\in\partial F(f)$ using
\Cref{thm:ch27-superdifferential}; the Fenchel equality
\eqref{eq:ch27-fenchel-equality} gives equality and attainment.
The obstacle functional is finite, concave and Lipschitz on a Banach space,
so the separation argument is available at every obstacle.

This measure-side minimum does not assert that $E_\omega^\vee(\mu)$
is attained at a finite-energy potential for a prescribed $\mu$.
That different assertion is
\Cref{cor:finite-energy-variational-characterization}, whose proof needs
\Cref{chap:bfj-finite-energy-measures}.

Part (2) of the proposition uses
\Cref{thm:bfj-envelope-regularity} and the approximation criterion:
uniform approximation of $P_\omega f$ by psh model functions compares
the continuous-obstacle supremum \eqref{eq:ch27-F-conjugate} with the
model-test supremum \eqref{eq:dual-energy-measure}.
This is group \textup{(EC5)} of \Cref{sec:ch27-local-uses}.
\end{remark}

\section{Orthogonality of envelopes}
\label{sec:ch27-orthogonality}

\subsection{The definition and its five faces}
\label{sec:ch27-orthogonality-definition}

\begin{definition}[Orthogonality property]
\label{def:orthogonality-property}
The form $\omega$ has the \emph{orthogonality property} if, for every
$f\in C^0(X^{\An})$,
\begin{equation}\label{eq:orthogonality}
 \int_{X^{\An}}
     (f-P_\omega(f))\,
     \MA_\omega(P_\omega(f))=0.
\end{equation}
\end{definition}

The integrand is nonnegative and Borel, and it is continuous as soon as
$P_\omega(f)$ is.  Thus \eqref{eq:orthogonality} is equivalent to saying
that $\MA_\omega(P_\omega(f))$ is concentrated on the contact locus
$\{P_\omega(f)=f\}$.  We say that \emph{orthogonality holds at $f$} when
\eqref{eq:orthogonality} holds for that particular $f$.

\begin{proposition}
\label{prop:ch27-orthogonality-equivalences}
Let $f\in C^0(X^{\An})$, put $u:=P_\omega(f)$,
$\mu:=\MA_\omega(u)$, and let $C_f$ be the contact set
\eqref{eq:ch27-contact-set}.  The following are equivalent.
\begin{enumerate}
\item $\displaystyle\int(f-u)\,\rd\mu=0$, that is, orthogonality holds at
$f$;
\item $\mu(X^{\An}\setminus C_f)=0$;
\item $\displaystyle\int f\,\rd\mu=\int u\,\rd\mu$;
\item $\mu\in\partial F(f)$, that is,
$F(g)\leq F(f)+\int(g-f)\,\rd\mu$ for every $g\in C^0(X^{\An})$;
\item $F(f)-\int f\,\rd\mu=E^\vee_\omega(\mu)$.
\end{enumerate}
The implications
$\textup{(1)}\Leftrightarrow\textup{(2)}\Leftrightarrow\textup{(3)}
\Rightarrow\textup{(4)}$
use no continuity property of $P_\omega(f)$ beyond boundedness.  The
equivalence $\textup{(4)}\Leftrightarrow\textup{(5)}$ does use it: it runs
through \Cref{prop:ch27-conjugate-identification}\,(2), whose proof
approximates $P_\omega f$ uniformly by $\omega$-psh model functions and
so spends \Cref{thm:bfj-envelope-regularity} --- this is item
\textup{(EC5)} of \Cref{sec:ch27-local-uses}.  So does the implication
$\textup{(4)}\Rightarrow\textup{(1)}$.
\end{proposition}

\begin{proof}
$(1)\Leftrightarrow(2)\Leftrightarrow(3)$.  The function $f-u$ is
nonnegative and Borel, and $\{f-u>0\}=X^{\An}\setminus C_f$; a nonnegative
Borel function has vanishing integral against a positive measure if and only
if it vanishes almost everywhere.  Both integrals in (3) are finite because
$f$ and $u$ are bounded and $\mu$ has mass $V$.

$(1)\Rightarrow(4)$.  Let $g\in C^0(X^{\An})$ and put
$v:=P_\omega(g)$.  Both $u$ and $v$ are bounded $\omega$-psh, hence
lie in $\mathcal E^1(X,\omega)$, so the right-hand inequality of the
two-sided sandwich \eqref{eq:ch27-two-sided-sandwich} gives
\[
 F(g)-F(f)=E_\omega(v)-E_\omega(u)\ \leq\ \int(v-u)\,\rd\mu .
\]
By (2) the measure $\mu$ is carried by $C_f$, and on $C_f$ we have
$u=f$ and $v\leq g$; hence $v-u\leq g-f$ $\mu$-almost everywhere,
and $\int(v-u)\,\rd\mu\leq\int(g-f)\,\rd\mu$.  This is (4).

$(4)\Rightarrow(1)$.  This is
\Cref{prop:ch27-supergradients-on-contact} applied to
$\ell=\mu$.

$(4)\Leftrightarrow(5)$ is
\Cref{prop:ch27-conjugate-identification}\,(3).
\end{proof}

\begin{remark}[The asymmetry of the equivalence]
\label{rem:ch27-orthogonality-asymmetry}
The direction that matters for the variational method is
$(1)\Rightarrow(4)$, and it is elementary once
\eqref{eq:ch27-two-sided-sandwich} is available: orthogonality upgrades the
Monge--Amp\`ere measure of the envelope from ``a measure attached to $u$''
to ``a supergradient of $F$ at $f$''.  The reverse direction is a
consequence of the general \Cref{prop:ch27-supergradients-on-contact},
which says that \emph{every} supergradient has that property; so
orthogonality is exactly the statement that the set
$\partial F(f)$, which is never empty, contains the one measure we can
compute.  This is the cleanest way to see that orthogonality is a genuine
theorem and not a formality: nothing in the soft theory identifies a
supergradient.
\end{remark}

\subsection{Reductions: uniform limits and reference change}
\label{sec:ch27-orthogonality-reductions}

Two elementary reductions cut the orthogonality property down to a much
smaller family of test data.  We first need a convergence lemma which is
used repeatedly below.

\begin{lemma}
\label{lem:ch27-uniform-ma}
Let $u$ and $u_j$ ($j\in\mathbf N$) be bounded $\omega$-psh
functions with $\lVert u_j-u\rVert_\infty\to0$.  Then
$\MA_\omega(u_j)\to\MA_\omega(u)$ weakly.  The same holds for a family
$(u_t)_{t\in I}$ depending on a real parameter with
$\lVert u_t-u\rVert_\infty\to0$ as $t\to t_0$.
\end{lemma}

\begin{proof}
It suffices to show that every subsequence has a further subsequence along
which the conclusion holds, since the assertion concerns convergence of the
real numbers $\int h\,\MA_\omega(u_j)$ for $h\in C^0(X^{\An})$.  So pass
to a subsequence with $\varepsilon_j:=\lVert u_j-u\rVert_\infty$
summable, and put
\[
 d_j:=\sum_{k\geq j}(\varepsilon_k+\varepsilon_{k+1}),
 \qquad v_j:=u_j+d_j .
\]
Then $v_j$ is $\omega$-psh, being $u_j$ plus a constant, and
$v_j\geq u_j+\varepsilon_j\geq u$.  Moreover
\[
 v_j-v_{j+1}=(u_j-u_{j+1})+(\varepsilon_j+\varepsilon_{j+1})\geq0,
\]
so $(v_j)$ is decreasing, and $d_j\to0$ gives $v_j\searrow u$.  All
these functions are bounded, hence lie in $\mathcal E^1(X,\omega)$, so
\Cref{prop:e1-monotone-continuity} gives
$\MA_\omega(v_j)\to\MA_\omega(u)$ weakly.  The bounded decreasing-net continuity
\eqref{eq:weighted-monotone-ma} also gives this directly.
Finally
$\MA_\omega(v_j)=\MA_\omega(u_j)$, because adding a constant does not
change $\ddc$.  The parametric statement follows by applying the
sequential one to arbitrary sequences $t_k\to t_0$.  This is the same
device as in the proof of \Cref{thm:bfj-cld-ma-common-locus}.
\end{proof}

\begin{proposition}
\label{prop:ch27-orthogonality-uniform-limits}
Let $f_j,f\in C^0(X^{\An})$ with $\lVert f_j-f\rVert_\infty\to0$.  If
orthogonality holds at every $f_j$, it holds at $f$.  Consequently
$\omega$ has the orthogonality property if and only if
\eqref{eq:orthogonality} holds for every \emph{model} function $f$.
\end{proposition}

\begin{proof}
By the contraction estimate \eqref{eq:ch23-envelope-contraction},
$\lVert P_\omega(f_j)-P_\omega(f)\rVert_\infty\leq\lVert
f_j-f\rVert_\infty\to0$.  Hence the continuous functions
$g_j:=f_j-P_\omega(f_j)$ converge uniformly to
$g:=f-P_\omega(f)$; here we use \Cref{thm:bfj-envelope-regularity} to know
that these functions are continuous, so that
\Cref{lem:joint-weak-uniform} applies.  By
\Cref{lem:ch27-uniform-ma}, $\MA_\omega(P_\omega f_j)\to
\MA_\omega(P_\omega f)$ weakly, and all these measures have mass $V$.
Therefore \Cref{lem:joint-weak-uniform} gives
\[
 0=\int g_j\,\MA_\omega(P_\omega f_j)
 \ \longrightarrow\
 \int g\,\MA_\omega(P_\omega f),
\]
which is \eqref{eq:orthogonality} at $f$.  The final assertion follows
because model functions are uniformly dense in $C^0(X^{\An})$ by
\Cref{lem:ch23-model-density}.
\end{proof}

\begin{proposition}
\label{prop:ch27-orthogonality-reference-change}
Let $w$ be a model function and set $\omega':=\omega+\ddc w$.  Then for
every $f\in C^0(X^{\An})$,
\begin{equation}\label{eq:ch27-orthogonality-reference}
 \int(f-P_\omega f)\,\MA_\omega(P_\omega f)
 =\int\bigl((f-w)-P_{\omega'}(f-w)\bigr)\,
   \MA_{\omega'}\bigl(P_{\omega'}(f-w)\bigr).
\end{equation}
In particular orthogonality holds at $f$ for $\omega$ if and only if it
holds at $f-w$ for $\omega'$; and $\omega$ has the orthogonality
property if and only if
\[
 \int_{X^{\An}}\bigl(-P_{\omega'}(0)\bigr)\,
   \MA_{\omega'}\bigl(P_{\omega'}(0)\bigr)=0
\]
for every closed model form $\omega'$ with $\{\omega'\}=\{\omega\}$.
\end{proposition}

\begin{proof}
The form $\omega'$ need not be semipositive.  Its psh class and measures
are interpreted by the model-potential change of reference of
\Cref{cor:ch21-reference-independence}; in particular we do not assume
that zero is $\omega'$-psh.
The reference-change identity \eqref{eq:ch23-reference-change} gives
$P_{\omega'}(f-w)=P_\omega(f)-w$.  Consequently
\[
 \omega'+\ddc\bigl(P_{\omega'}(f-w)\bigr)
 =\omega+\ddc w+\ddc P_\omega(f)-\ddc w
 =\omega+\ddc P_\omega(f),
\]
so the two Monge--Amp\`ere measures in
\eqref{eq:ch27-orthogonality-reference} coincide.  The two integrands
coincide as well, since
$(f-w)-\bigl(P_\omega(f)-w\bigr)=f-P_\omega(f)$.  This proves
\eqref{eq:ch27-orthogonality-reference}.

For the last assertion, combine
\Cref{prop:ch27-orthogonality-uniform-limits} with the case $w=f$ of
\eqref{eq:ch27-orthogonality-reference}: for a model obstacle $f$,
orthogonality at $f$ for $\omega$ is orthogonality at $0$ for
$\omega+\ddc f$, and by \Cref{thm:ch20-ddc-exact} the forms
$\omega+\ddc f$, for $f$ running over model functions, are exactly the
closed model forms with the same de Rham class as $\omega$.
\end{proof}

\begin{remark}[Why the reduction to a zero obstacle is the right one]
\label{rem:ch27-zero-obstacle}
\Cref{prop:ch27-orthogonality-reference-change} explains the shape of the
literature.  Orthogonality is never proved obstacle by obstacle; it is
proved once, for the zero obstacle, for an arbitrary form in the given
ample de Rham class.  In that formulation $P_{\omega'}(0)$ is the
\emph{semipositive envelope} of the reference metric, the quantity
$-P_{\omega'}(0)$ measures how far the reference is from being
semipositive, and \eqref{eq:orthogonality} is the assertion that this defect
is invisible to the Monge--Amp\`ere measure of the envelope.  Written in
terms of line bundles it becomes a statement about the vertical part of an
approximate Zariski decomposition, which is where the real proof happens;
see \Cref{sec:ch27-orthogonality-mechanism}.  It is also the shape in which
\Cref{lem:ch23-reduction-to-base-ideals} reduces envelope continuity, which
is why the two theorems have the same reduction step and completely
different engines.
\end{remark}

\subsection{Why a formal slackness argument is not a proof}
\label{sec:ch27-orthogonality-heuristic}

\begin{remark}[Why a formal slackness argument is not a proof]
\label{rem:model-slackness-heuristic}
On a curve, or in the polyhedral toric obstacle problem, one can prove
directly that the Laplacian or Alexandrov measure of an envelope is supported
on its contact set; the examples at the end of the chapter carry this out.
In higher-dimensional model geometry it is tempting to argue as follows.  If
the vertical divisor measuring $f-P_\omega(f)$ has a positive coefficient
along a component $E$, try to raise the nef envelope along $E$; maximality
would then appear to force vanishing of the top intersection on $E$.

That perturbation need not remain nef: a curve on which nefness is already
an equality can obstruct it immediately.  Thus this argument neither proves
a componentwise complementary-slackness statement nor
\eqref{eq:orthogonality}.  BFJ instead use approximate Zariski decompositions:
if $mL=M_m+F_m$ is the normalized base-ideal decomposition, the required
integral is the limit of the normalized intersections
\[
                  m^{-n}M_m^n\cdot m^{-1}F_m.
\]
Their vanishing is the global asymptotic orthogonality theorem, not a formal
consequence of nefness.  \Cref{thm:orthogonality-provenance} records
the hypotheses under which that theorem is known.
\end{remark}

The reader who wants to see exactly which step of the heuristic collapses
should compare two computations already available.  In
\Cref{thm:ch23-model-orthogonality-curve} the constraint set is the
polyhedron \eqref{eq:ch23-curve-polyhedron} cut out by the \emph{affine}
functions $A_i$ of \eqref{eq:ch23-curve-residual-degree}, so that raising a
single coordinate changes each constraint by a computable amount of a fixed
sign, and the model measure \emph{is} the vector $(A_i)$.  In dimension
$n\geq2$ neither feature survives: relative nefness is tested on
all vertical curves, rather than only the finitely many components; and the
atom of the model measure at a divisorial point is a polynomial of degree
$n$ in the vertex values, so its positivity says nothing about a
directional derivative of the constraint.  This is
\Cref{rem:ch23-model-slackness-fails}, and it is the precise reason why the
present chapter has to import \Cref{thm:orthogonality-provenance}.

\subsection{The mechanism: asymptotic orthogonality}
\label{sec:ch27-orthogonality-mechanism}

\index{Zariski decomposition!approximate}
\index{orthogonality!asymptotic}
It is worth describing, without proof, the shape of the argument that does
work, because it is the reason orthogonality carries the hypotheses it
does.

By \Cref{prop:ch27-orthogonality-reference-change} we may take the obstacle
to be zero and let the reference form vary in a fixed ample de Rham class.
Suppose that the class is that of a model line bundle $\mathcal L$ on a
normal model $\mathcal X$, with $L$ ample on $X$.  For each sufficiently large $m$
let $\mathfrak a_m\subset\cO_{\mathcal X}$ be the base ideal of
$\mathcal L^{\otimes m}$, that is, the image of
$H^0(\mathcal X,\mathcal L^{\otimes m})\otimes\mathcal L^{-m}
 \to\cO_{\mathcal X}$, and let
\[
 \varphi_m:=\tfrac1m\log|\mathfrak a_m|
\]
be the associated model function, as in
\eqref{eq:ch23-base-potential}.  Blowing up $\mathfrak a_m$ writes
\[
 m\,\mathcal L=M_m+F_m,
\]
with $M_m$ globally generated (the ``moving part'') and $F_m$ effective
and vertical (the ``fixed part''); this is the \emph{approximate Zariski
decomposition} at level $m$.  Two facts are then in play.

\emph{First}, $\varphi_m$ converges uniformly to $P_\omega(0)$, and
may be taken increasing along a suitable divisibility sequence; this is
\Cref{thm:ch23-base-ideal-approximation}, and it is the same statement that
yields continuity of the envelope in \Cref{chap:bfj-envelopes}.  In
particular the model functions of $m^{-1}F_m$ converge uniformly to
$-P_\omega(0)$.  The moving metrics of $m^{-1}M_m$ converge to the
envelope metric; the divisors live on varying models.

\emph{Second}, the integral \eqref{eq:orthogonality} for the zero obstacle
is the limit of the normalized intersection numbers
\begin{equation}\label{eq:ch27-asymptotic-orthogonality}
 \bigl(m^{-1}M_m\bigr)^n\cdot\bigl(m^{-1}F_m\bigr)
 \ \longrightarrow\ 0 .
\end{equation}
This is the \emph{asymptotic orthogonality} statement, and it is not formal.
Its absolute analogue---for a big line bundle on a projective variety over a
field, with $M_m,F_m$ the Fujita approximation data---is a theorem of
independent standing, proved through differentiability of the volume
function \cite{BFJ09}, restricted volumes \cite{ELMNP09}, Okounkov bodies
\cite{LM09}, and the duality between the pseudoeffective and movable cones
\cite{BDPP13}.  The relative statement needed here is
\cite[Appendix~A]{BFJ15}, and the mechanism by which BFJ obtain it is to
\emph{spread the model out over a projective curve} and apply the absolute
theorem on the resulting projective variety over the residue field.  That spreading out
is exactly where algebraizability enters, and it is the only place.

The later route uses a \emph{relative} non-Archimedean volume
$\operatorname{Vol}(L,\phi,\phi_0)$ of two metrics on the same line
bundle.  Fixing $\phi_0$ turns it into a function of $\phi$; the relative
volume is a normalized asymptotic determinant-lattice volume, with no
canonical absolute reference suppressed.  In additive metric-weight
notation, \cite[Theorem~5.4.3]{BGGJKM20} gives, at a continuous
semipositive weight $u$,
\[
 \operatorname{Vol}(L,u+\varepsilon g,u)
   =\varepsilon\int g\,\MA(u)+o(\varepsilon).
\]
If $u=P_\omega f$ is the continuous envelope of an obstacle $f$, then
$u\leq u+\varepsilon(f-u)\leq f$ for $0\leq\varepsilon\leq1$.
Every intermediate obstacle has envelope $u$.
Passing to a continuous semipositive envelope preserves the small-section
lattices \cite[Proposition~6.2.1]{BGGJKM20}, so
\[
 \operatorname{Vol}(L,u+\varepsilon(f-u),u)=0.
\]
The derivative formula therefore yields
$\int(f-u)\MA(u)=0$.  This is
\cite[Theorem~6.3.2]{BGGJKM20}.  It needs continuity of the envelope to
apply the derivative theorem at $u$; the derivative theorem itself does
not assume continuity of envelopes for arbitrary obstacles.
The real ample-class extension uses rational approximation, homogeneity
and the closedness of the orthogonality locus in the ample cone, as in
\cite[Theorem~6.3.3]{BGGJKM20}.  The source keeps this volume derivative
separate from the projected-energy derivative of the present chapter.

\subsection{Provenance}
\label{sec:ch27-orthogonality-provenance}

\begin{theorem}\label{thm:orthogonality-provenance}
The following statements must be distinguished.
\begin{enumerate}
\item In the original proof of \cite[Theorem~A.6]{BFJ15}, $X$ is assumed
\emph{algebraizable}: there are a one-variable function field $F$ admitting
$K$ as a completion and a smooth projective $F$-variety $Y$ with
$X\simeq Y_K$. Under this additional hypothesis, every ample class has the
orthogonality property.
\item Burgos Gil--Gubler--Jell--K\"unnemann--Martin later proved that if
$K$ is complete discretely valued of residue characteristic zero and $X$
is smooth projective over $K$, then every model form with ample de Rham
class has the orthogonality property, without algebraizability.
\end{enumerate}
\par
\end{theorem}

\begin{proof}[Source architecture]
For (1), BFJ identify \eqref{eq:orthogonality} with asymptotic orthogonality
in approximate Zariski decompositions on models. Algebraizability lets them
spread a model over a projective curve and apply absolute volume
differentiability and orthogonality.

For (2), the differentiability input is
\cite[Theorem~5.4.3]{BGGJKM20}; it identifies the derivative of the
non-Archimedean volume with the Chambert--Loir measure.  Since passing to
the semipositive envelope does not change the small-section volume,
\cite[Theorem~6.3.2]{BGGJKM20} turns that derivative formula into
orthogonality for a continuous semipositive envelope, and
\cite[Theorem~6.3.3]{BGGJKM20} records the resulting BFJ statement in
residue characteristic zero.  Residue characteristic zero supplies continuity of envelopes there.
The extension from line bundles to real ample classes also uses rational
approximation and homogeneity, as stated in Theorem~6.3.3.  This later route
removes algebraizability; it does not remove the discrete valuation,
smoothness, or ample-class assumptions in the statement used here.
\end{proof}

Three supplementary remarks fix the scope of the theorem.

\begin{remark}[What each route costs]
\label{rem:ch27-two-routes}
Within the residue-characteristic-zero framework of \cite{BFJ15}, route
(1) costs algebraizability, an assumption about the \emph{origin} of $X$:
it must descend to a variety over a one-variable function field.  This is
an additional descent hypothesis, different in kind from envelope
regularity.
Route (2) removes that origin condition and instead passes through
continuity and regularization of the envelope.  Its differentiability-to-
orthogonality step is characteristic-neutral once those inputs and the
required model theory are available.  This is the route used in equal
characteristic $p$ by \cite[Theorem~9.2]{GJKM19}, after test ideals supply
envelope continuity and the stated resolution hypotheses supply the model
calculus.  Thus the positive-characteristic result is not an extension of
the original asymptotic route; it is an application of the later volume
route with different regularity input.  The final ledger is
\Cref{sec:ch27-where-continuity}.
\end{remark}

\begin{remark}[Orthogonality is a property of the class, not of the form]
\label{rem:ch27-orthogonality-class}
By \Cref{prop:ch27-orthogonality-reference-change}, if
$\{\omega'\}=\{\omega\}$ then $\omega$ has the orthogonality property if
and only if $\omega'$ does.  This is why
\Cref{thm:orthogonality-provenance} may speak of an ample \emph{class}.  It
also explains why the property is never stated for a nef but non-ample
class in this chapter: \Cref{lem:ch23-envelope-finiteness}, on which the
whole envelope calculus rests, uses ampleness through
\Cref{prop:ch20-positive-representative}.
\end{remark}

\begin{center}
\fbox{\begin{minipage}{0.91\textwidth}
\textbf{Warning (historical hypotheses).}
Envelope continuity in \Cref{chap:bfj-envelopes} does not by itself
imply orthogonality. In BFJ's 2015 Monge--Amp\`ere paper, orthogonality was
proved with algebraizability. The removal of that hypothesis is a later
theorem of \cite{BGGJKM20}, not part of the original proof.
\end{minipage}}
\end{center}

\begin{warning}{Orthogonality is not a maximum principle}
Both \eqref{eq:orthogonality} and the bounded maximum principle
\eqref{eq:bounded-maximum-principle} assert that a Monge--Amp\`ere measure
lives where two functions agree, but their hypotheses and mechanisms
differ.  The maximum principle compares
$\max\{\varphi,\psi\}$ with $\varphi$ on the Borel locus
$\{\varphi>\psi\}$, where \emph{both} functions are psh and the identity
is local; its proof is \Cref{prop:model-maximum-principle} followed by a
capacity argument.  Orthogonality compares a psh function with a
\emph{non-psh} obstacle on the set where they agree, and there is no local
identity to exploit: the contact set can be totally disconnected and the
statement is global.  In dimension one the two do come together, because
there the constraint set is linear; that coincidence is
\Cref{thm:ch23-model-orthogonality-curve} and it is misleading.
\end{warning}

\subsection{Curves: from model level to continuous obstacles}
\label{sec:ch27-curve-orthogonality}

\Cref{chap:bfj-envelopes} proves orthogonality at model level on a curve, by
complementary slackness on the dual graph.  We record what that gives here,
and exactly what it does not.

\begin{proposition}
\label{prop:ch27-curve-orthogonality}
Let $n=1$, let $\mathcal X$ be a regular projective model of $X$ on
which $\omega$ is determined and semipositive, and let $f$ be a real
model function determined on $\mathcal X$.  Assume that the model-level
envelope of \Cref{def:ch23-fixed-model-envelope} computes the global
envelope, that is,
\begin{equation}\label{eq:ch27-curve-model-level-hypothesis}
 P^{\mathcal X}_\omega(f)=P_\omega(f).
\end{equation}
Then orthogonality holds at $f$:
\[
 \int_{X^{\An}}\bigl(f-P_\omega(f)\bigr)\,
   \MA_\omega\bigl(P_\omega(f)\bigr)=0 .
\]
Hypothesis \eqref{eq:ch27-curve-model-level-hypothesis} holds whenever
$\mathcal X$ is strictly semistable.
\end{proposition}

\begin{proof}
Write $P:=P^{\mathcal X}_\omega(f)$.  By
\Cref{thm:ch23-model-orthogonality-curve}\,(1) the model-level envelope
exists, and by part (3) of that theorem
$\int(f-P)\MA_\omega(P)=0$.  Under
\eqref{eq:ch27-curve-model-level-hypothesis} this is the assertion.  The
final sentence is \Cref{rem:ch23-model-dependence}, which proves
$P^{\mathcal X}_\omega(f)=P_\omega(f)$ for a strictly semistable model
determining $\omega$ and $f$, by the chord argument of the proof of
\Cref{thm:curves-arbitrary-field-continuity}.
\end{proof}

\begin{corollary}
\label{cor:ch27-curve-orthogonality-global}
Let $n=1$ and suppose that the strictly semistable models of $X$
determining $\omega$ are cofinal among all models determining
$\omega$.  Then $\omega$ has the orthogonality property.
\end{corollary}

\begin{proof}
By \Cref{prop:ch27-orthogonality-uniform-limits} it suffices to treat model
obstacles.  A model function $f$ is determined on some model
$\mathcal X_0$; by hypothesis there is a strictly semistable
$\mathcal X\geq\mathcal X_0$ determining $\omega$, and $f$ is then
determined on $\mathcal X$ as well, since a function determined on a model
is determined on every model dominating it.  Now apply
\Cref{prop:ch27-curve-orthogonality}.
\end{proof}

\begin{remark}[What is missing in the general curve case, and where to get
it]
\label{rem:ch27-curve-base-change}
The cofinality hypothesis in
\Cref{cor:ch27-curve-orthogonality-global} is an additional conditional
hypothesis, not a consequence of semistable reduction.  For example,
blowing up a node of a reduced semistable fibre introduces a component of
multiplicity two.  A reduced semistable model cannot dominate that
model while retaining this divisorial valuation.  Semistable reduction
itself may also require a finite extension $K'/K$.  Transporting
\Cref{cor:ch27-curve-orthogonality-global} back to $K$ therefore requires
two compatibilities: that the envelope commutes with the finite base change,
and that Monge--Amp\`ere measures do.  The second is
\Cref{thm:base-change-ma}, in the form
$\pi_*\MA_{\omega_{K'}}(\varphi\circ\pi)=\MA_\omega(\varphi)$, together
with \eqref{eq:base-change-weight}.  For the rational-model envelope,
the first is \cite[Proposition~2.11]{GJKM19} after enlarging to a finite
normal extension, including its inseparable part.  Applying this identity
to the raw all-psh envelope additionally requires the appropriate envelope
comparison.  We do not infer that comparison from notation, or state an
unconditional curve orthogonality theorem in this chapter: what is proved
here is \Cref{prop:ch27-curve-orthogonality} and
\Cref{cor:ch27-curve-orthogonality-global}, and the general curve case is
covered, with its own hypotheses and its own source, by
\Cref{thm:orthogonality-provenance} in the residue characteristic zero
regime and by the sources listed in
\Cref{sec:ch27-where-continuity} otherwise.
\end{remark}

\subsection{The Legendre form of the obstacle problem}
\label{sec:ch27-legendre}

The second case in which orthogonality is provable outright is the convex
one.  Here it becomes a statement about Legendre transforms, and---this is
the point of the subsection---the \emph{same} computation simultaneously
yields the derivative of the energy.  The convex statement is proved in
full; its transport to invariant metrics on a toric variety is
\Cref{thm:ch27-toric-orthogonality}.

Throughout, $N_{\mathbf R}\cong\mathbf R^n$ and $M_{\mathbf R}$ are dual
real vector spaces with dual lattices $N$ and $M$, and
$\operatorname{Leb}_M$ is the lattice Lebesgue measure on
$M_{\mathbf R}$ normalized by $M$, as in
\eqref{eq:appD-alexandrov-measure}.  For a compact convex body
$Q\subset M_{\mathbf R}$ let
\[
 h_Q(x):=\sup_{m\in Q}\langle m,x\rangle,
 \qquad x\in N_{\mathbf R},
\]
be its support function.

\begin{theorem}
\label{thm:ch27-legendre-orthogonality}
\index{Legendre duality!and orthogonality}
Let $Q\subset M_{\mathbf R}$ be a compact convex body with nonempty
interior, and let $f:N_{\mathbf R}\to\mathbf R$ be continuous with
\[
 C:=\lVert f-h_Q\rVert_\infty<\infty .
\]
Let $u$ be the convex envelope of $f$, that is, the largest convex
function on $N_{\mathbf R}$ with $u\leq f$, and let
$C_f:=\{u=f\}$ be the contact set.  Then:
\begin{enumerate}
\item $f^*(m)=\sup_{x}(\langle m,x\rangle-f(x))$ is finite exactly on
$Q$, where $|f^*|\leq C$; moreover $u=f^{**}$, $u^*=f^*$, and
$\lVert u-h_Q\rVert_\infty\leq C$;
\item for every $m$ in the interior of $Q$ the supremum defining
$f^*(m)$ is attained, the set $\mathcal M(m)$ of maximizers is compact
and nonempty, and
\[
 \mathcal M(m)\subset C_f\cap\partial u^*(m);
\]
in particular $\mathcal M(m)=\{\nabla u^*(m)\}$ for
$\operatorname{Leb}_M$-almost every $m\in Q$;
\item after extending the almost-everywhere-defined gradient $\nabla u^*$
arbitrarily across its Lebesgue-null exceptional set by a Borel map,
$\MA_{\mathbf R}(u)=(\nabla u^*)_*\bigl(
\operatorname{Leb}_M|_Q\bigr)$, of total mass
$\operatorname{vol}_M(Q)$, it is carried by $C_f$, and
\begin{equation}\label{eq:ch27-convex-orthogonality}
 \int_{N_{\mathbf R}}\bigl(f-u\bigr)\,\MA_{\mathbf R}(u)=0;
\end{equation}
\item for every bounded continuous $g:N_{\mathbf R}\to\mathbf R$ the
function
\[
 \Lambda(t):=-\int_Q(f+tg)^*\,\rd\operatorname{Leb}_M
\]
is concave and $\operatorname{vol}_M(Q)\lVert g\rVert_\infty$-Lipschitz on
$\mathbf R$, and it is differentiable at $t=0$ with
\begin{equation}\label{eq:ch27-convex-derivative}
 \Lambda'(0)=\int_{N_{\mathbf R}}g\,\MA_{\mathbf R}(u).
\end{equation}
\end{enumerate}
\end{theorem}

\begin{proof}
\emph{(1).}  For $m\in Q$ we have $\langle m,x\rangle\leq h_Q(x)$ for all
$x$, so
\[
 \langle m,x\rangle-f(x)\leq\langle m,x\rangle-h_Q(x)+C\leq C,
\]
whence $f^*(m)\leq C$; and $f^*(m)\geq\langle m,0\rangle-f(0)\geq
-h_Q(0)-C=-C$ since $h_Q(0)=0$.  If $m\notin Q$, separate: there is
$x_0$ with $\langle m,x_0\rangle>h_Q(x_0)$, and then
$\langle m,sx_0\rangle-f(sx_0)\geq
s(\langle m,x_0\rangle-h_Q(x_0))-C\to+\infty$ as $s\to+\infty$.  Thus
$\operatorname{dom}f^*=Q$.

Since $f\geq h_Q-C$ admits affine minorants, $f^{**}$ is a finite convex
function, and by the Fenchel--Moreau theorem
\Cref{thm:appD-fenchel-moreau} it is the largest lower semicontinuous convex
minorant of $f$.  A finite convex function on $N_{\mathbf R}$ is
continuous, hence lower semicontinuous; so $f^{**}$ is also the largest
convex minorant, that is, $u=f^{**}$.  Conjugating once more gives
$u^*=f^{***}=f^*$.  Finally
\[
 u(x)=\sup_{m\in Q}\bigl(\langle m,x\rangle-f^*(m)\bigr),
 \qquad
 h_Q(x)=\sup_{m\in Q}\langle m,x\rangle,
\]
and $|f^*|\leq C$ on $Q$ gives
$\lVert u-h_Q\rVert_\infty\leq C$.

\emph{(2).}  Let $m$ be interior to $Q$, and choose $\delta>0$ with
$m+\delta \bar B\subset Q$, where $\bar B$ is the closed unit ball for
some norm $|\cdot|$ on $M_{\mathbf R}$ with dual norm again written
$|\cdot|$.  Then $h_Q(x)\geq\langle m,x\rangle+\delta|x|$, so
\[
 \langle m,x\rangle-f(x)\leq-\delta|x|+C .
\]
Since $f^*(m)\geq-C$, every maximizing sequence eventually lies in the
compact set $\{|x|\leq(2C+1)/\delta\}$; by continuity of $f$ the supremum is
attained there and $\mathcal M(m)$ is a nonempty compact set.

Let $x\in\mathcal M(m)$.  Using $u\leq f$ and $u^*=f^*$,
\[
 f^*(m)=\langle m,x\rangle-f(x)
       \leq\langle m,x\rangle-u(x)
       \leq u^*(m)=f^*(m),
\]
so both inequalities are equalities.  The first gives $u(x)=f(x)$, that is,
$x\in C_f$.  The second is the equality case of the Fenchel inequality for
the pair $(u,u^*)$, which by
\Cref{prop:appD-fenchel-equality} and \eqref{eq:appD-subgradient} means
$m\in\partial u(x)$, equivalently $x\in\partial u^*(m)$.  Since $u^*$
is a finite convex function on the interior of $Q$, it is differentiable
outside a Lebesgue-null set, and at a differentiability point
$\partial u^*(m)=\{\nabla u^*(m)\}$; the boundary of $Q$ is also
Lebesgue-null.

\emph{(3).}  By part (1), $\operatorname{dom}u^*=Q$, so
$\overline{\partial u(N_{\mathbf R})}=Q$ as in the proof of
\Cref{prop:appD-legendre-pushforward}; that proposition then gives both the
pushforward formula \eqref{eq:appD-legendre-pushforward} and the mass
\eqref{eq:appD-real-ma-mass}.  By part (2), for almost every $m\in Q$ the
point $\nabla u^*(m)$ is a maximizer, hence lies in $C_f$.  Therefore
the pushforward measure gives no mass to $N_{\mathbf R}\setminus C_f$,
and, $f-u$ being nonnegative and vanishing on $C_f$,
\eqref{eq:ch27-convex-orthogonality} follows.  (This reproves
\Cref{thm:ch23-convex-envelope-orthogonality} in the present global
situation; the argument there, through nondifferentiability of the
conjugate, is the local form of the same duality.)

\emph{(4).}  Write $\varphi_t(m):=(f+tg)^*(m)$.  Since
$\lVert(f+tg)-h_Q\rVert_\infty\leq C+|t|\lVert g\rVert_\infty$, part (1)
applies to $f+tg$ for every $t$; in particular $\varphi_t$ is finite on
$Q$ and $\Lambda$ is finite.  As a supremum of affine functions of
$t$,
\[
 \varphi_t(m)=\sup_x\bigl(\langle m,x\rangle-f(x)-tg(x)\bigr)
\]
is convex in $t$; hence $\Lambda$ is concave.  From
$f+tg\geq f-|t|\lVert g\rVert_\infty$ and its companion we get
$|\varphi_t-\varphi_0|\leq|t|\lVert g\rVert_\infty$ pointwise, whence the
Lipschitz bound for $\Lambda$.

Fix $m$ in the interior of $Q$ with $\mathcal M(m)=\{x_m\}$ a
singleton, which by (2) is the case for almost every $m$.  We show
\begin{equation}\label{eq:ch27-danskin}
 \lim_{t\to0}\frac{\varphi_t(m)-\varphi_0(m)}{t}=-g(x_m).
\end{equation}
The coercivity estimate of (2), applied to $f+tg$ with
$|t|\leq1$, shows that all maximizers for $f+tg$ at $m$ lie in the
fixed compact set
$S:=\{|x|\leq2(C+\lVert g\rVert_\infty)/\delta\}$.  Let $x_t\in S$ be
one of them.  Testing the supremum for $\varphi_t$ at the point $x_m$
gives
\[
 \varphi_t(m)\geq\langle m,x_m\rangle-f(x_m)-tg(x_m)
 =\varphi_0(m)-tg(x_m),
\]
while testing the supremum for $\varphi_0$ at $x_t$ gives
\[
 \varphi_t(m)=\langle m,x_t\rangle-f(x_t)-tg(x_t)
 \leq\varphi_0(m)-tg(x_t).
\]
For $t>0$ these read
$-g(x_m)\leq(\varphi_t(m)-\varphi_0(m))/t\leq-g(x_t)$, and for $t<0$ the
two inequalities are reversed.  It remains to check that $x_t\to x_m$ as
$t\to0$.  Let $x_0$ be a limit point of $(x_t)$ in the compact set
$S$; passing to the limit in
$\varphi_t(m)=\langle m,x_t\rangle-f(x_t)-tg(x_t)$, and using continuity
of $f$ and $g$ and $\varphi_t(m)\to\varphi_0(m)$, gives
$\varphi_0(m)=\langle m,x_0\rangle-f(x_0)$, so $x_0\in\mathcal M(m)=
\{x_m\}$.  Hence $x_t\to x_m$, $g(x_t)\to g(x_m)$, and
\eqref{eq:ch27-danskin} follows.

Finally, the difference quotients
$(\varphi_t-\varphi_0)/t$ are bounded in absolute value by
$\lVert g\rVert_\infty$, and $Q$ has finite measure; so dominated
convergence along an arbitrary sequence $t_k\to0$ gives
\[
 \Lambda'(0)
 =-\lim_{t\to0}\int_Q\frac{\varphi_t-\varphi_0}{t}\,
       \rd\operatorname{Leb}_M
 =\int_Qg\bigl(\nabla u^*(m)\bigr)\,\rd\operatorname{Leb}_M(m)
 =\int_{N_{\mathbf R}}g\,\MA_{\mathbf R}(u),
\]
the last equality being \eqref{eq:appD-legendre-pushforward}.
\end{proof}

\begin{remark}[One computation, two theorems]
\label{rem:ch27-one-computation}
The proof above is worth reading twice.  Part (2)---``every maximizer of the
Legendre supremum lies in the contact set''---is orthogonality, and part
(4)---``the derivative of the conjugate integral is integration of the
direction against the pushforward measure''---is differentiability of the
envelope energy.  They are the \emph{same} sentence: the maximizer
$\nabla u^*(m)$ appears in both, once as a point of $C_f$ and once as the
evaluation point of $g$.  In the general non-Archimedean situation the
Legendre transform is unavailable, and this single sentence splits into the
two logically separate statements \eqref{eq:orthogonality} and
\eqref{eq:energy-envelope-derivative}, whose equivalence is
\Cref{cor:ch27-orthogonality-iff-differentiability}.  The convex case is
therefore not merely an example; it is a proof, in one family, of the
theorem this chapter cannot otherwise prove, and it shows what structure a
proof needs.
\end{remark}

\subsection{Transport to invariant toric metrics}
\label{sec:ch27-toric-transport}

We now translate.  Let $X_\Sigma$ be a \emph{smooth} projective toric
variety of dimension $n$ over $K$ with dense torus $T=\Spec K[M]$, so
that the standing hypotheses of this chapter are in force, let $L$ be an
ample toric line bundle with polytope $P\subset M_{\mathbf R}$, and put
\[
 Q:=-P .
\]
With the tropical coordinate \eqref{eq:appJ-toric-coordinate}, the recession
function of a toric metric weight is
$\check h_P=h_Q$ by \eqref{eq:appJ-toric-recession}, so the convex
functions of \Cref{thm:ch27-legendre-orthogonality} with body $Q$ are
exactly the weights of invariant semipositive metrics on $L$, by
\Cref{cor:ch19-toric-convexity-test}.  Let $\trop:X_\Sigma^{\An}\to
N_\Sigma$ be the extended tropicalization and
$\iota:N_{\mathbf R}\to\Sk(T)\subset X_\Sigma^{\An}$ the Gauss section, so
that $\trop\circ\iota=\operatorname{id}$ by
\Cref{lem:ch19-toric-chart}\,(1).  Call a function on $X_\Sigma^{\An}$
\emph{invariant} if it factors through $\trop$.

Fix an invariant semipositive \emph{model} metric $\lVert\cdot\rVert_0$ on
$L$, with weight $w_0$ a rational piecewise affine convex function with
$w_0-h_Q$ bounded, and set $\omega:=c_1(L,\lVert\cdot\rVert_0)$; thus
$V=(L^n)=n!\operatorname{vol}_M(P)=n!\operatorname{vol}_M(Q)$ by
\eqref{eq:appJ-toric-degree}.  For an invariant $f\in C^0(X_\Sigma^{\An})$
write $w_f:=w_0+f\circ\iota$ for the corresponding function on
$N_{\mathbf R}$, and let
\begin{equation}\label{eq:ch27-invariant-envelope}
 P^{\mathrm{inv}}_\omega(f):=
 \sup\{\varphi\in\PSH(X,\omega)\ \text{invariant}:\varphi\leq f\}
\end{equation}
be the \emph{invariant} envelope.

\begin{theorem}
\label{thm:ch27-toric-orthogonality}
In the above situation, let $f$ be an invariant rational model function.  Then
$w_f$ is rational piecewise affine with $w_f-h_Q$ bounded, its convex
envelope $u_f$ is again rational piecewise affine convex with the same
recession function, $P^{\mathrm{inv}}_\omega(f)=u_f\circ\trop-w_0$ is an invariant
semipositive model metric weight, and
\begin{equation}\label{eq:ch27-toric-orthogonality}
 \int_{X_\Sigma^{\An}}
   \bigl(f-P^{\mathrm{inv}}_\omega(f)\bigr)\,
   \MA_\omega\bigl(P^{\mathrm{inv}}_\omega(f)\bigr)=0 .
\end{equation}
Moreover
$\trop_*\MA_\omega(P^{\mathrm{inv}}_\omega(f))=n!\,\MA_{\mathbf R}(u_f)$, a
measure of total mass $V$ carried by the contact set
$\{u_f=w_f\}$.
\end{theorem}

\begin{proof}
Since $f$ is a rational model function it is continuous and bounded, so
$w_f-h_Q=(w_0-h_Q)+f\circ\iota$ is bounded; and $w_f$ is rational
piecewise affine because $w_0$ and $f\circ\iota$ are.  The epigraph of
$w_f$ is a finite union of rational polyhedra; its closed convex hull is a
rational polyhedron, hence $u_f:=\operatorname{conv}(w_f)$
is rational piecewise affine and convex, and by
\Cref{thm:ch27-legendre-orthogonality}\,(1) it satisfies
$\lVert u_f-h_Q\rVert_\infty\leq\lVert w_f-h_Q\rVert_\infty$.  By
\Cref{thm:toric-positivity} and \Cref{cor:ch19-toric-convexity-test},
$u_f\circ\trop-w_0$ is therefore the weight of an invariant semipositive
model metric, that is, an invariant $\omega$-psh model function; and it is
$\leq f$ because $u_f\leq w_f$.  Conversely every invariant
$\omega$-psh $\varphi\leq f$ has convex weight $w_0+\varphi\circ\iota
\leq w_f$, hence $w_0+\varphi\circ\iota\leq u_f$.  This proves
$P^{\mathrm{inv}}_\omega(f)=u_f\circ\trop-w_0$.

By \Cref{thm:ch19-toric-ma}, applied to the invariant semipositive metric
with weight $u_f$---and using that for a model metric the BFJ measure
\eqref{eq:model-mixed-measure} and the Chambert--Loir measure agree, by
\Cref{thm:bfj-cld-ma-common-locus}---the measure
$\MA_\omega(P^{\mathrm{inv}}_\omega(f))$ is supported on the toric skeleton and
\[
 \trop_*\MA_\omega\bigl(P^{\mathrm{inv}}_\omega(f)\bigr)=n!\,\MA_{\mathbf R}(u_f),
\]
of total mass $n!\operatorname{vol}_M(Q)=V$.  The function
$f-P^{\mathrm{inv}}_\omega(f)$ is invariant and its restriction to the skeleton is
$(w_f-u_f)\circ\trop$; since the measure lives on the skeleton, on which
$\trop$ is a homeomorphism onto $N_{\mathbf R}$ by
\Cref{lem:ch19-toric-chart}\,(1),
\[
 \int_{X_\Sigma^{\An}}\bigl(f-P^{\mathrm{inv}}_\omega(f)\bigr)\,
   \MA_\omega\bigl(P^{\mathrm{inv}}_\omega(f)\bigr)
 =n!\int_{N_{\mathbf R}}(w_f-u_f)\,\MA_{\mathbf R}(u_f)=0
\]
by \eqref{eq:ch27-convex-orthogonality}.  The last assertion is
\Cref{thm:ch27-legendre-orthogonality}\,(3).
\end{proof}

\begin{corollary}
\label{cor:ch27-toric-energy}
Let $\varphi,\varphi'$ be bounded continuous invariant $\omega$-psh
functions, in particular model functions, with
convex weights $w=w_0+\varphi\circ\iota$ and
$w'=w_0+\varphi'\circ\iota$.  Then
\begin{equation}\label{eq:ch27-toric-energy}
 E_\omega(\varphi)-E_\omega(\varphi')
 =n!\int_Q\bigl(w'^{\,*}-w^{*}\bigr)\,\rd\operatorname{Leb}_M .
\end{equation}
\end{corollary}

\begin{proof}
Apply \Cref{thm:ch26-toric-energy} to $w,w'$ with the fixed model
reference $w_0$.  It applies to bounded continuous invariant semipositive
metrics; the measures agree with the BFJ continuous product by
\Cref{prop:ch24-continuous-uniform}.
In the roof convention, $g_w(y)=-w^*(-y)$ and $Q=-P$.
Changing variables $m=-y$ and subtracting the two roof identities gives
\eqref{eq:ch27-toric-energy}.
This proof also establishes the formula for bounded continuous invariant
psh weights, including real parameter values along the examples below.
\end{proof}

\begin{remark}[Scope of the toric statement]
\label{rem:ch27-toric-scope}
Three restrictions in \Cref{thm:ch27-toric-orthogonality} are real and must
not be dropped.

\emph{Invariance of the obstacle.}  The theorem computes the
\emph{invariant} envelope \eqref{eq:ch27-invariant-envelope}.  It does not
assert $P^{\mathrm{inv}}_\omega(f)=P_\omega(f)$, and no averaging argument
supplies that equality: invariance is a hypothesis and not a conclusion in
the toric dictionary, as \Cref{sec:ch19-toric-continuous} states explicitly
in its item on invariance.  Consequently
\Cref{thm:ch27-toric-orthogonality} is a verification of
\eqref{eq:orthogonality} in a subcategory, not a proof of the orthogonality
property of $\omega$ in the sense of
\Cref{def:orthogonality-property}.

\emph{Model obstacles.}  The proof used piecewise affineness twice, to know
that $u_f$ is again a model weight and that
\Cref{thm:ch19-toric-ma} applies to it.  The passage to an arbitrary
invariant continuous obstacle is the argument of
\Cref{prop:ch27-orthogonality-uniform-limits} verbatim, and it goes through
as soon as the obstacle is a uniform limit of invariant model obstacles.

\emph{Ampleness.}  $Q$ must have nonempty interior; this is where
ampleness of $L$ enters, and it is what makes the coercivity estimate in
part (2) of \Cref{thm:ch27-legendre-orthogonality} available.  For a nef but
non-ample toric class the body $Q$ may be lower dimensional and the
Legendre supremum need not be attained.
\end{remark}

\section{Differentiability along model directions}
\label{sec:ch27-model-directions}

\subsection{The theorem and its two-line mechanism}
\label{sec:ch27-model-mechanism}

\begin{theorem}
\label{thm:ch27-model-direction-differentiability}
Assume that $\omega$ has the orthogonality property.  Let $f,g$ be real
model functions and put
\[
 u_t:=P_\omega(f+tg),
 \qquad \mu_t:=\MA_\omega(u_t),
 \qquad t\in\mathbf R .
\]
Then $t\mapsto F(f+tg)=E_\omega(u_t)$ is differentiable on all of
$\mathbf R$, with
\begin{equation}\label{eq:ch27-model-direction-derivative}
 \frac{\rd}{\rd t}E_\omega\bigl(P_\omega(f+tg)\bigr)
 =\int_{X^{\An}}g\,\rd\mu_t
 \qquad\text{for every }t\in\mathbf R,
\end{equation}
and the derivative is a nonincreasing continuous function of $t$.
\end{theorem}

\begin{proof}
It suffices to prove differentiability at $t=0$ with derivative
$\int g\,\rd\mu_0$; the general $t$ follows by replacing $f$ with
$f+tg$, which is again a real model function.

\emph{Step 1: the family $u_t$ is uniformly close to $u_0$.}
Each $f+tg$ is a real model function, hence continuous; so
$u_t$ is a bounded $\omega$-psh function by
\Cref{cor:ch23-continuous-obstacle-psh} and
\Cref{cor:ch23-continuous-finite}, and $\mu_t$ is a positive measure of
mass $V$ by \Cref{thm:bounded-ma-extension}.  The contraction estimate
\eqref{eq:ch23-envelope-contraction} gives
\begin{equation}\label{eq:ch27-envelope-lipschitz-in-t}
 \lVert u_t-u_0\rVert_\infty\ \leq\ |t|\,\lVert g\rVert_\infty ,
\end{equation}
so by \Cref{lem:ch27-uniform-ma},
\begin{equation}\label{eq:ch27-measures-converge}
 \mu_t\longrightarrow\mu_0
 \quad\text{weakly as }t\to0;
 \qquad\text{in particular}\qquad
 \int g\,\rd\mu_t\longrightarrow\int g\,\rd\mu_0 .
\end{equation}

\emph{Step 2: what orthogonality gives.}
Applying \Cref{def:orthogonality-property} to the continuous obstacle
$f+tg$ and using
\Cref{prop:ch27-orthogonality-equivalences}\,$(1)\Leftrightarrow(2)$, the
measure $\mu_t$ is carried by the contact set of $f+tg$; that is,
\begin{equation}\label{eq:ch27-contact-at-t}
 u_t=f+tg
 \qquad\mu_t\text{-almost everywhere},
 \qquad\text{for every }t\in\mathbf R .
\end{equation}
In particular $u_0=f$ $\mu_0$-almost everywhere.

\emph{Step 3: the upper estimate.}
Both $u_0$ and $u_t$ are bounded $\omega$-psh, hence lie in
$\mathcal E^1(X,\omega)$.  The right-hand inequality of the two-sided
sandwich \eqref{eq:ch27-two-sided-sandwich} gives
\[
 E_\omega(u_t)-E_\omega(u_0)\ \leq\ \int(u_t-u_0)\,\rd\mu_0 .
\]
Now $\mu_0$-almost everywhere $u_0=f$, by \eqref{eq:ch27-contact-at-t}
at $t=0$, while $u_t\leq f+tg$ everywhere by definition of the envelope.
Hence $u_t-u_0\leq tg$ $\mu_0$-almost everywhere, and
\begin{equation}\label{eq:ch27-upper-estimate}
 E_\omega(u_t)-E_\omega(u_0)\ \leq\ t\int g\,\rd\mu_0 .
\end{equation}

\emph{Step 4: the lower estimate.}
The left-hand inequality of \eqref{eq:ch27-two-sided-sandwich} gives
\[
 E_\omega(u_t)-E_\omega(u_0)\ \geq\ \int(u_t-u_0)\,\rd\mu_t .
\]
Now $\mu_t$-almost everywhere $u_t=f+tg$, by
\eqref{eq:ch27-contact-at-t}, while $u_0\leq f$ everywhere.  Hence
$u_t-u_0\geq f+tg-f=tg$ $\mu_t$-almost everywhere, and
\begin{equation}\label{eq:ch27-lower-estimate}
 E_\omega(u_t)-E_\omega(u_0)\ \geq\ t\int g\,\rd\mu_t .
\end{equation}

\emph{Step 5: conclusion.}
For $t>0$, dividing \eqref{eq:ch27-upper-estimate} and
\eqref{eq:ch27-lower-estimate} by $t$ gives
\[
 \int g\,\rd\mu_t
 \ \leq\ \frac{F(f+tg)-F(f)}{t}
 \ \leq\ \int g\,\rd\mu_0 ;
\]
for $t<0$ the two inequalities are reversed.  In both cases
\eqref{eq:ch27-measures-converge} squeezes the difference quotient, so the
two one-sided derivatives at $0$ exist and both equal
$\int g\,\rd\mu_0$.  This is
\eqref{eq:ch27-model-direction-derivative} at $t=0$.

Finally, $t\mapsto F(f+tg)$ is concave by
\Cref{prop:ch27-F-elementary}\,(4), so its derivative is nonincreasing; and
it is continuous by \eqref{eq:ch27-measures-converge} applied at each base
point.
\end{proof}

\begin{center}
\fbox{\begin{minipage}{0.91\textwidth}
\textbf{Where each hypothesis is spent.}
Orthogonality is used \emph{twice}: at the obstacle $f$, to know that
$\mu_0$ sees only the contact set, and at the obstacle $f+tg$, to know
that $\mu_t$ does.  Dropping either use leaves only one of
\eqref{eq:ch27-upper-estimate} and \eqref{eq:ch27-lower-estimate}, hence
only a one-sided bound, hence---by
\Cref{ex:ch27-corner-without-orthogonality}---no information at all.  The
energy calculus is used through the two-sided sandwich
\eqref{eq:ch27-two-sided-sandwich}, which is \textup{(D6)} and is
independent of orthogonality.  Continuity of envelopes is \emph{not} used
in this proof, beyond its role upstream in constructing
$\MA_\omega$ and $E_\omega$ at all.
\end{minipage}}
\end{center}

\subsection{The finite-dimensional shadow}
\label{sec:ch27-finite-shadow}

Before passing to continuous directions it is worth seeing
\Cref{thm:ch27-model-direction-differentiability} in the one situation where
the whole obstacle problem is a linear program, because there every step of
the proof becomes a line of linear algebra and the role of complementary
slackness is visible.

\begin{proposition}
\label{prop:ch27-curve-shadow}
Let $n=1$, let $\mathcal X$ be a regular model with dual graph
$\Gamma$ as in \eqref{eq:ch23-curve-special-fibre}, let $\omega$ be a
semipositive model form determined on $\mathcal X$, and let $u$ be a
real model function determined on $\mathcal X$.  Identify model functions
determined on $\mathcal X$ with vectors $v\in\mathbf R^I$, write
$A_i(v)$ as in \eqref{eq:ch23-curve-residual-degree}, and let
$K=K_{\mathcal X}(\omega,u)$ be the polyhedron
\eqref{eq:ch23-curve-polyhedron}.  Put
\[
 \Phi(u):=E_\omega\bigl(P^{\mathcal X}_\omega(u)\bigr).
\]
Then $\Phi$ is concave and $\deg_X\{\omega\}$-Lipschitz on
$\mathbf R^I$ for the maximum norm, and for every $g\in\mathbf R^I$,
\begin{equation}\label{eq:ch27-curve-shadow-derivative}
 \frac{\rd}{\rd t}\Big|_{t=0}\Phi(u+tg)
 =\sum_{i\in I}A_i(P)\,g_i
 =\int_{X^{\An}}g\,\MA_\omega\bigl(P^{\mathcal X}_\omega(u)\bigr),
 \qquad P:=P^{\mathcal X}_\omega(u),
\end{equation}
the derivative existing for every $t$, not only at $t=0$.
\end{proposition}

\begin{proof}
Existence and maximality of $P^{\mathcal X}_\omega(u)$ are
\Cref{thm:ch23-model-orthogonality-curve}\,(1).  Concavity of
$u\mapsto P^{\mathcal X}_\omega(u)$ is immediate from
\eqref{eq:ch23-curve-polyhedron}: if $v\in K(u)$ and $v'\in K(u')$ then
$\lambda v+(1-\lambda)v'\in K(\lambda u+(1-\lambda)u')$, because the
obstacle constraints are affine and $A_i$ is affine.  Composing with the
concave increasing functional $E_\omega$ gives concavity of $\Phi$, and
the Lipschitz bound follows from $P^{\mathcal X}_\omega(u+c\one)=
P^{\mathcal X}_\omega(u)+c\one$ and monotonicity, exactly as in
\Cref{prop:ch27-F-elementary}\,(3), the mass being
$\deg_X\{\omega\}$ by \eqref{eq:ch23-curve-model-measure}.

For \eqref{eq:ch27-curve-shadow-derivative}, note first that the second
equality is \eqref{eq:ch23-curve-model-measure}.  For the first, run the
proof of \Cref{thm:ch27-model-direction-differentiability} with
$P^{\mathcal X}_\omega$ in place of $P_\omega$: Steps~3 and~4 are
unchanged, and Step~2, which is the only place orthogonality was invoked, is
now the elementary complementary slackness
\eqref{eq:ch23-complementary-slackness} proved in
\Cref{thm:ch23-model-orthogonality-curve}\,(2), applied to the obstacles
$u$ and $u+tg$.  Step~1 holds because
$\lVert P^{\mathcal X}_\omega(u+tg)-P^{\mathcal X}_\omega(u)\rVert_\infty
\leq|t|\lVert g\rVert_\infty$, which follows from
\eqref{eq:ch23-curve-polyhedron} by the same translation argument as
\eqref{eq:ch23-envelope-contraction}; and the measures depend continuously
on the vertex data by \eqref{eq:ch23-curve-model-measure}, since
$A_i$ is affine; so the squeeze in Step~5 of that proof closes and the
two-sided derivative exists at every $t$.  This proves
\eqref{eq:ch27-curve-shadow-derivative}.

It is worth recording the finite-dimensional reason for the regularity just
used.  Each coordinate
$t\mapsto P^{\mathcal X}_\omega(u+tg)_i
 =\max\{v_i:v\leq u+tg,\ A(v)\geq0\}$
is the optimal value of a linear program whose objective is fixed and whose
right-hand side depends affinely on $t$; such a value function is concave
and piecewise affine in $t$.  Hence $t\mapsto P^{\mathcal
X}_\omega(u+tg)$ is piecewise affine and $\Phi(u+tg)$ is piecewise
polynomial, in agreement with the explicit formulas of
\Cref{ex:ch27-curve-derivative}.
\end{proof}

\begin{remark}[What the shadow does and does not model]
\label{rem:ch27-shadow-scope}
\Cref{prop:ch27-curve-shadow} is a genuine theorem, but it is a theorem
about a linear program.  The two features that make it elementary are
exactly the two that
\Cref{rem:ch23-model-slackness-fails} identifies as one dimensional: the
constraint $A_i(v)\geq0$ is affine in $v$, and the measure \emph{is} the
vector $(A_i(v))_i$ of constraint values, so complementary slackness and
the derivative formula are the same statement.  For $n\geq2$ the
constraint set is not polyhedral in the vertex data and the measure is a
degree-$n$ polynomial in them, and no part of the argument survives.  The
value of the shadow is that it shows precisely which structural fact
orthogonality has to replace.
\end{remark}

\section{Differentiability along continuous directions}
\label{sec:ch27-continuous-directions}

\subsection{Reduction of the continuous case to the model case}
\label{sec:ch27-reduction}

We now prove the theorem in the form in which it is used.  The proof takes
\Cref{thm:ch27-model-direction-differentiability} as a black box and uses no
further geometric input; in particular \emph{no further use of
orthogonality is made}.  This modularity is the reason for separating the
two statements: if orthogonality were available only for model obstacles,
the theorem below would still follow.

\begin{theorem}\label{thm:energy-envelope-differentiability}
Assume that $\omega$ has the orthogonality property. Then
\[
                E_\omega\circ P_\omega:C^0(X^{\An})\longrightarrow\mathbf R
\]
is G\^ateaux differentiable, and for $f,g\in C^0(X^{\An})$,
\begin{equation}\label{eq:energy-envelope-derivative}
 \left.\frac{\rd}{\rd t}\right|_{t=0}
 E_\omega(P_\omega(f+tg))
 =\int g\,\MA_\omega(P_\omega(f)).
\end{equation}
This is \cite[Theorem~7.2]{BFJ15}.
\end{theorem}

\begin{proof}
Fix $f,g\in C^0(X^{\An})$ and set
\[
 h(t):=F(f+tg)=E_\omega\bigl(P_\omega(f+tg)\bigr),
 \qquad
 \gamma(s):=\int_{X^{\An}}g\,\MA_\omega\bigl(P_\omega(f+sg)\bigr).
\]
By \Cref{lem:ch23-model-density} choose real model functions $f_j$ and
$g_j$ with
$\lVert f_j-f\rVert_\infty\to0$ and $\lVert g_j-g\rVert_\infty\to0$, and
put $h_j(t):=F(f_j+tg_j)$.

\emph{Step 1: $h_j\to h$ uniformly on bounded intervals.}
By \Cref{prop:ch27-F-elementary}\,(3),
\[
 |h_j(t)-h(t)|
 \leq V\lVert (f_j-f)+t(g_j-g)\rVert_\infty
 \leq V\bigl(\lVert f_j-f\rVert_\infty
       +|t|\,\lVert g_j-g\rVert_\infty\bigr),
\]
which tends to $0$ uniformly for $t$ in a bounded interval.  In
particular $h$ is concave, being a pointwise limit of the concave
functions $h_j$.

\emph{Step 2: the derivatives converge pointwise.}
Each $h_j$ is differentiable on $\mathbf R$ with
$h_j'(s)=\int g_j\,\MA_\omega(P_\omega(f_j+sg_j))$, by
\Cref{thm:ch27-model-direction-differentiability}.  Fix $s$.  Since
\[
 \lVert (f_j+sg_j)-(f+sg)\rVert_\infty\longrightarrow0,
\]
the contraction estimate \eqref{eq:ch23-envelope-contraction} gives
$\lVert P_\omega(f_j+sg_j)-P_\omega(f+sg)\rVert_\infty\to0$, and then
\Cref{lem:ch27-uniform-ma} gives
$\MA_\omega(P_\omega(f_j+sg_j))\to\MA_\omega(P_\omega(f+sg))$ weakly.  All
these measures have mass $V$, and $g_j\to g$ uniformly; therefore
\Cref{lem:joint-weak-uniform} applies and yields
\begin{equation}\label{eq:ch27-derivatives-converge}
 h_j'(s)\longrightarrow\gamma(s)
 \qquad\text{for every }s\in\mathbf R .
\end{equation}

\emph{Step 3: passage to the limit in the integral.}
Each $h_j$ is concave and differentiable, hence locally Lipschitz and
absolutely continuous, so
\[
 h_j(t)-h_j(0)=\int_0^th_j'(s)\,\rd s
 \qquad(t\in\mathbf R).
\]
The functions $h_j'$ are nonincreasing, by concavity of $h_j$, and
uniformly bounded:
$|h_j'(s)|\leq V\lVert g_j\rVert_\infty\leq M$ for a constant $M$
independent of $j$, because $\lVert g_j\rVert_\infty\to\lVert
g\rVert_\infty$.  By \eqref{eq:ch27-derivatives-converge} and
\Cref{lem:ch27-concave-nets}\,(3),
\[
 \int_0^th_j'(s)\,\rd s\ \longrightarrow\ \int_0^t\gamma(s)\,\rd s .
\]
Combining with Step~1,
\begin{equation}\label{eq:ch27-integrated-derivative}
 h(t)-h(0)=\int_0^t\gamma(s)\,\rd s
 \qquad\text{for every }t\in\mathbf R .
\end{equation}

\emph{Step 4: continuity of $\gamma$.}
By \eqref{eq:ch23-envelope-contraction},
$\lVert P_\omega(f+sg)-P_\omega(f+s_0g)\rVert_\infty
\leq|s-s_0|\lVert g\rVert_\infty$, so
\Cref{lem:ch27-uniform-ma} gives weak convergence of the corresponding
Monge--Amp\`ere measures as $s\to s_0$, and hence
$\gamma(s)\to\gamma(s_0)$.  Thus $\gamma$ is continuous on
$\mathbf R$.

\emph{Step 5: conclusion.}
By \eqref{eq:ch27-integrated-derivative} and Step~4, $h$ is continuously
differentiable with $h'(t)=\gamma(t)$; in particular
\[
 \left.\frac{\rd}{\rd t}\right|_{t=0}E_\omega(P_\omega(f+tg))
 =\gamma(0)=\int g\,\MA_\omega(P_\omega f),
\]
which is \eqref{eq:energy-envelope-derivative}.  The right-hand side is
linear in $g$, so $F$ is G\^ateaux differentiable at $f$ with
differential $\MA_\omega(P_\omega f)$.
\end{proof}

\begin{corollary}
\label{cor:ch27-singleton}
Assume orthogonality.  Then for every $f\in C^0(X^{\An})$,
\[
 \partial F(f)=\bigl\{\MA_\omega(P_\omega f)\bigr\}.
\]
\end{corollary}

\begin{proof}
By \Cref{thm:energy-envelope-differentiability}, the functional $F$ is
G\^ateaux differentiable at $f$ with differential
$\MA_\omega(P_\omega f)$.  The singleton criterion in
\Cref{thm:ch27-superdifferential}\,(4) therefore shows that
$\partial F(f)$ has exactly one element, and the min--max formula
\eqref{eq:ch27-minmax-formula} identifies that element with the displayed
differential.
\end{proof}

\begin{remark}[Two proofs, and why the second one is written down]
\label{rem:ch27-two-proofs}
The direct argument of
\Cref{thm:ch27-model-direction-differentiability} never used that $f$ and
$g$ are model functions; it used only that $f+tg$ is continuous, so that
\Cref{def:orthogonality-property} applies to it.  Read that way, it proves
\Cref{thm:energy-envelope-differentiability} in five lines.  The
approximation proof given above is longer, and it is written down because it
proves a strictly stronger statement: \emph{differentiability along
continuous directions follows from differentiability along model directions
by soft arguments alone}.  In particular, in any situation where the
orthogonality property is established only for model obstacles---which is
how it arises from an intersection-theoretic argument, since
\eqref{eq:ch27-asymptotic-orthogonality} is a statement about model line
bundles---one still obtains
\eqref{eq:energy-envelope-derivative} for all continuous data.  The two
convergence inputs are \Cref{lem:ch27-uniform-ma}, which upgrades uniform
convergence of potentials to weak convergence of measures, and
\Cref{lem:ch27-concave-nets}\,(3), which passes the integral identity through
pointwise convergence of uniformly bounded monotone derivatives.  Here the
model approximation was chosen as a sequence, so ordinary dominated
convergence would also suffice; the net form is recorded for the chosen finite-energy regularization in
\Cref{cor:e1-projected-variation}; a potential-specific sequence is an
alternative there.
\end{remark}

\begin{corollary}
\label{cor:ch27-linearization}
Assume orthogonality and let $f,g\in C^0(X^{\An})$.  Put
\[
 \mu:=\MA_\omega(P_\omega f).
\]
Then the concave function
\[
 q(t):=\int_{X^{\An}}P_\omega(f+tg)\,\rd\mu
\]
is differentiable at $t=0$ and
\begin{equation}\label{eq:linearized-envelope-energy}
 \left.\frac{\rd}{\rd t}\right|_{0+}E_\omega(P_\omega(f+tg))
 =
 \left.\frac{\rd}{\rd t}\right|_{0+}
 \int P_\omega(f+tg)\,\rd \mu ,
\end{equation}
both sides being equal to $\int g\,\rd\mu$.
\end{corollary}

\begin{proof}
Concavity of $q$ follows from concavity of the envelope
\eqref{eq:ch23-joint-concavity} and positivity of $\mu$; so $q'(0+)$
exists by \Cref{lem:ch27-concave-quotients}.

\emph{Upper bound.}  By orthogonality, $P_\omega f=f$ $\mu$-almost
everywhere, while $P_\omega(f+tg)\leq f+tg$ everywhere.  Hence
$\mu$-almost everywhere
$P_\omega(f+tg)\leq P_\omega(f)+tg$, and integrating,
$q(t)\leq q(0)+t\int g\,\rd\mu$.  Dividing by $t>0$ and letting
$t\downarrow0$ gives $q'(0+)\leq\int g\,\rd\mu$.

\emph{Lower bound.}  The right-hand inequality of
\eqref{eq:ch27-two-sided-sandwich}, applied to
$u=P_\omega f$ and $v=P_\omega(f+tg)$, gives
\[
 F(f+tg)-F(f)\ \leq\ \int\bigl(P_\omega(f+tg)-P_\omega f\bigr)\,\rd\mu
 =q(t)-q(0).
\]
Dividing by $t>0$ and letting $t\downarrow0$, and using
\Cref{thm:energy-envelope-differentiability},
\[
 \int g\,\rd\mu
 =\left.\frac{\rd}{\rd t}\right|_{0+}F(f+tg)
 \ \leq\ q'(0+).
\]
The two bounds coincide, proving
\eqref{eq:linearized-envelope-energy} for the right derivative.
Apply the same argument to $-g$.  Since the resulting function is $q(-t)$,
its right derivative is $-q'(0-)$, and its computed value is
$-\int g\,\rd\mu$.  Thus the left derivative has the same value and
$q$ is differentiable at zero, as stated.
\end{proof}

\subsection{The converse: differentiability forces orthogonality}
\label{sec:ch27-converse}

\begin{proposition}\label{prop:differentiability-iff-orthogonality}
If \eqref{eq:energy-envelope-derivative} holds for every continuous
direction, then $\omega$ has the orthogonality property.
\end{proposition}

\begin{proof}
Fix $f\in C^0(X^{\An})$, put $\phi:=P_\omega(f)$, and set
\[
 g:=\phi-f\ \leq\ 0 .
\]
The function $g$ is continuous: this is the point at which continuity of
envelopes, \Cref{thm:bfj-envelope-regularity}, is used, and it is used only
to make $g$ an admissible direction in
\eqref{eq:energy-envelope-derivative}.

For $0\leq\varepsilon\leq1$,
\[
 f+\varepsilon g=(1-\varepsilon)f+\varepsilon\phi\ \geq\ \phi ,
\]
because $\phi\leq f$.  Thus $\phi$ is an $\omega$-psh function
dominated by $f+\varepsilon g$, whence
$P_\omega(f+\varepsilon g)\geq\phi$, and monotonicity of energy
(\Cref{cor:energy-monotone-concave}) gives
\[
 E_\omega\bigl(P_\omega(f+\varepsilon g)\bigr)\ \geq\ E_\omega(\phi)
 =E_\omega\bigl(P_\omega(f)\bigr)
 \qquad(0\leq\varepsilon\leq1).
\]
Therefore the right derivative at $\varepsilon=0$ of
$\varepsilon\mapsto E_\omega(P_\omega(f+\varepsilon g))$ is nonnegative.
By hypothesis that derivative equals
\[
 \int g\,\MA_\omega\bigl(P_\omega f\bigr)
 =\int(\phi-f)\,\MA_\omega(\phi)\ \leq\ 0,
\]
the inequality because $\phi-f\leq0$ and $\MA_\omega(\phi)\geq0$.  Hence
$\int(f-\phi)\MA_\omega(\phi)=0$, which is
\eqref{eq:orthogonality} at $f$.  As $f$ was arbitrary, $\omega$ has
the orthogonality property.
\end{proof}

\begin{corollary}
\label{cor:ch27-orthogonality-iff-differentiability}
The following are equivalent.
\begin{enumerate}
\item $\omega$ has the orthogonality property;
\item $E_\omega\circ P_\omega$ is G\^ateaux differentiable on
$C^0(X^{\An})$, with differential $\MA_\omega(P_\omega f)$ at $f$;
\item $\partial F(f)=\{\MA_\omega(P_\omega f)\}$ for every
$f\in C^0(X^{\An})$;
\item $\MA_\omega(P_\omega f)\in\partial F(f)$ for every
$f\in C^0(X^{\An})$;
\item \eqref{eq:orthogonality} holds for every model function $f$.
\end{enumerate}
\end{corollary}

\begin{proof}
$(1)\Rightarrow(2)$ is
\Cref{thm:energy-envelope-differentiability};
$(2)\Rightarrow(1)$ is
\Cref{prop:differentiability-iff-orthogonality};
$(2)\Leftrightarrow(3)$ is
\Cref{thm:ch27-superdifferential}\,(4) together with
\Cref{cor:ch27-singleton}; $(1)\Leftrightarrow(4)$ is
\Cref{prop:ch27-orthogonality-equivalences}\,$(1)\Leftrightarrow(4)$;
and $(1)\Leftrightarrow(5)$ is
\Cref{prop:ch27-orthogonality-uniform-limits}.
\end{proof}

\begin{center}
\fbox{\begin{minipage}{0.91\textwidth}
\textbf{Reading \Cref{cor:ch27-orthogonality-iff-differentiability}
correctly.}
The corollary is an equivalence of two statements, not a proof of either.
The forward implication has orthogonality as a \emph{hypothesis}, and the
backward implication has differentiability as a \emph{hypothesis}; no
statement of \Cref{sec:ch27-directional} or
\Cref{sec:ch27-superdifferential} is used in a way that presupposes either.
It is precisely because the two are equivalent that the historical routes to
orthogonality recorded in \Cref{thm:orthogonality-provenance} proceed
through a \emph{third} functional, the non-Archimedean volume, whose
differentiability is proved by intersection theory.  A proof of
orthogonality that used
\Cref{thm:energy-envelope-differentiability} would be circular; none of the
sources cited here does so.
\end{minipage}}
\end{center}

\section{Finite-energy envelope variations}
\label{sec:ch27-e1-variations}

The variational method perturbs a \emph{finite-energy} potential, not a
continuous one, so the derivative formula must be extended from
$C^0(X^{\An})$ to obstacles of the form $\varphi+tf$ with
$\varphi\in\mathcal E^1(X,\omega)$.  Such an obstacle is upper
semicontinuous and may take the value $-\infty$; the extension is
therefore not formal, but it needs no new geometric input.

\subsection{The envelope of a finite-energy obstacle}
\label{sec:ch27-e1-envelope}

\begin{lemma}
\label{lem:ch27-e1-envelope-calculus}
Let $\varphi\in\PSH(X,\omega)$ and $f\in C^0(X^{\An})$, and let
$t\in\mathbf R$.  Then:
\begin{enumerate}
\item $P_\omega(\varphi+tf)\in\PSH(X,\omega)$, and it is the largest
$\omega$-psh function dominated by $\varphi+tf$;
\item $\varphi-|t|\lVert f\rVert_\infty\ \leq\ P_\omega(\varphi+tf)
\ \leq\ \varphi+|t|\lVert f\rVert_\infty$;
\item for $s,t\in\mathbf R$,
\begin{equation}\label{eq:ch27-e1-contraction}
 P_\omega(\varphi+sf)\ \leq\ P_\omega(\varphi+tf)
   +|s-t|\,\lVert f\rVert_\infty ;
\end{equation}
\item if $\varphi\in\mathcal E^1(X,\omega)$, then
$P_\omega(\varphi+tf)\in\mathcal E^1(X,\omega)$ and
\begin{equation}\label{eq:ch27-e1-energy-lipschitz}
 \bigl|E_\omega\bigl(P_\omega(\varphi+tf)\bigr)-E_\omega(\varphi)\bigr|
 \ \leq\ |t|\,\lVert f\rVert_\infty\,V .
\end{equation}
\end{enumerate}
\end{lemma}

\begin{proof}
(1) The obstacle $\varphi+tf$ is upper semicontinuous, being the sum of the
upper semicontinuous function $\varphi$ and the continuous function
$tf$, and it is bounded above because $\varphi$ is
(\Cref{def:ch21-bfj-psh}) and $f$ is.  The admissible family is nonempty:
$\varphi-|t|\lVert f\rVert_\infty$ is $\omega$-psh by
\Cref{prop:ch21-elementary-stability} and lies below $\varphi+tf$.  Now
apply \Cref{cor:ch23-usc-obstacle}.

(2) The lower bound is the competitor just exhibited; the upper bound holds
because $P_\omega(\varphi+tf)\leq\varphi+tf\leq\varphi+|t|\lVert
f\rVert_\infty$.

(3) From $\varphi+sf\leq(\varphi+tf)+|s-t|\lVert f\rVert_\infty$ and parts
(1) and (2) of \Cref{prop:ch23-formal-envelope-calculus}, which are proved
there for arbitrary obstacles with finite envelopes.

(4) Apply monotonicity of $E_\omega$ (\Cref{thm:extended-energy-basic}) and
the translation rule \eqref{eq:bfj-energy-translation} to the two
inequalities of (2):
\[
 E_\omega(\varphi)-|t|\lVert f\rVert_\infty V
 \leq E_\omega\bigl(P_\omega(\varphi+tf)\bigr)
 \leq E_\omega(\varphi)+|t|\lVert f\rVert_\infty V .
\]
Since $E_\omega(\varphi)>-\infty$, the middle term is finite.
\end{proof}

\begin{lemma}
\label{lem:ch27-e1-envelope-approximation}
Let $\varphi\in\mathcal E^1(X,\omega)$, $f\in C^0(X^{\An})$, and let
$(\varphi_\alpha)$ be a decreasing net of $\omega$-psh model functions
with $\varphi_\alpha\searrow\varphi$, as furnished by
\Cref{thm:bfj-decreasing-regularization}.  Then, for every $t\in\mathbf R$,
\[
 P_\omega(\varphi_\alpha+tf)\ \searrow\ P_\omega(\varphi+tf),
\]
and consequently
\[
 E_\omega\bigl(P_\omega(\varphi_\alpha+tf)\bigr)
 \longrightarrow E_\omega\bigl(P_\omega(\varphi+tf)\bigr),
\]
and
\[
 \MA_\omega\bigl(P_\omega(\varphi_\alpha+tf)\bigr)
 \longrightarrow
 \MA_\omega\bigl(P_\omega(\varphi+tf)\bigr)
\]
weakly.
\end{lemma}

\begin{proof}
Monotonicity of the envelope makes
$\alpha\mapsto P_\omega(\varphi_\alpha+tf)$ a decreasing net of
$\omega$-psh functions, each $\geq P_\omega(\varphi+tf)$.  Let
$\psi:=\inf_\alpha P_\omega(\varphi_\alpha+tf)$.  Since
$\psi\geq P_\omega(\varphi+tf)\geq\varphi-|t|\lVert f\rVert_\infty$, which
is not identically $-\infty$, \Cref{thm:ch21-decreasing-stability} gives
$\psi\in\PSH(X,\omega)$.  Moreover
\[
 \psi\ \leq\ \inf_\alpha\bigl(\varphi_\alpha+tf\bigr)=\varphi+tf ,
\]
so $\psi\leq P_\omega(\varphi+tf)$ by maximality
(\Cref{lem:ch27-e1-envelope-calculus}\,(1)).  Hence
$\psi=P_\omega(\varphi+tf)$.

All the functions involved lie in $\mathcal E^1(X,\omega)$, by
\Cref{lem:ch27-e1-envelope-calculus}\,(4) applied to $\varphi_\alpha$ and
to $\varphi$.  The convergence statements are therefore
\Cref{prop:e1-monotone-continuity}.
\end{proof}

\subsection{The derivative}
\label{sec:ch27-e1-derivative}

\begin{corollary}\label{cor:e1-projected-variation}
Assume orthogonality. If
$\varphi\in\mathcal E^1(X,\omega)$ and $f\in C^0(X^{\An})$, then
\[
 P_\omega(\varphi+tf)\in\mathcal E^1(X,\omega)
\]
for every $t\in\mathbf R$, and
\begin{equation}\label{eq:e1-projected-derivative}
 \left.\frac{\rd}{\rd t}\right|_{t=0}
 E_\omega(P_\omega(\varphi+tf))
 =\int f\,\MA_\omega(\varphi).
\end{equation}
\end{corollary}

\begin{proof}
Membership in $\mathcal E^1$ is
\Cref{lem:ch27-e1-envelope-calculus}\,(4).  Put
\[
 G(t):=E_\omega\bigl(P_\omega(\varphi+tf)\bigr),
 \qquad
 \gamma(s):=\int_{X^{\An}}f\,
     \MA_\omega\bigl(P_\omega(\varphi+sf)\bigr),
\]
and let $(\varphi_\alpha)$ be a decreasing net of $\omega$-psh model
functions with $\varphi_\alpha\searrow\varphi$.  Set
$G_\alpha(t):=E_\omega(P_\omega(\varphi_\alpha+tf))$.

\emph{Step 1: the approximants are differentiable.}
Each $\varphi_\alpha+tf$ is continuous, so
\Cref{thm:energy-envelope-differentiability} applies and gives
\[
 G_\alpha'(s)=\int f\,
   \MA_\omega\bigl(P_\omega(\varphi_\alpha+sf)\bigr),
 \qquad s\in\mathbf R .
\]
Moreover $G_\alpha$ is concave, by
\Cref{prop:ch27-F-elementary}\,(4), so $G_\alpha'$ is nonincreasing, and
$|G_\alpha'|\leq V\lVert f\rVert_\infty$ by the global Lipschitz bound
\Cref{prop:ch27-F-elementary}\,(3).

\emph{Step 2: pointwise convergence.}
By \Cref{lem:ch27-e1-envelope-approximation}, applied at the parameter
$s$, we get $G_\alpha(s)\to G(s)$ and $G_\alpha'(s)\to\gamma(s)$ for
every $s\in\mathbf R$.

\emph{Step 3: passing to the limit in the integral.}
Each $G_\alpha$ is concave and differentiable, hence absolutely continuous
on bounded intervals, so $G_\alpha(t)-G_\alpha(0)=\int_0^tG_\alpha'$.  The
net $(G_\alpha')$ consists of nonincreasing functions with values in
$[-V\lVert f\rVert_\infty,V\lVert f\rVert_\infty]$ and converges
pointwise to $\gamma$.  \Cref{lem:ch27-concave-nets}\,(3) therefore gives
$\int_0^tG_\alpha'\to\int_0^t\gamma$, and hence
\begin{equation}\label{eq:ch27-e1-integrated}
 G(t)-G(0)=\int_0^t\gamma(s)\,\rd s
 \qquad(t\in\mathbf R).
\end{equation}
For the chosen net, \Cref{lem:ch27-concave-nets}\,(3) supplies the
limit without an invalid application of dominated convergence to arbitrary
nets.  Alternatively, one may choose the potential-specific sequence of
\Cref{cor:psh-sequential-regularization} and use ordinary dominated
convergence.  That sequence does not make the model system countably cofinal.

\emph{Step 4: continuity of $\gamma$ at $0$.}
Note first that $P_\omega(\varphi)=\varphi$, so
$\gamma(0)=\int f\,\MA_\omega(\varphi)$.  Since the parameter space is
metric, it is enough to prove sequential continuity.  More precisely, it
suffices to show that every sequence $s_k\to0$ has, inside each of its
subsequences, a further subsequence along which
$\gamma(s_k)\to\gamma(0)$.  Pass to such a further subsequence so fast that
$\varepsilon_k:=|s_k|\lVert f\rVert_\infty$ is summable, and put
\[
 d_k:=\sum_{j\geq k}(\varepsilon_j+\varepsilon_{j+1}),
 \qquad
 w_k:=P_\omega(\varphi+s_kf)+d_k .
\]
By \Cref{lem:ch27-e1-envelope-calculus}\,(2), $w_k\geq
\varphi-\varepsilon_k+d_k\geq\varphi$; by
\eqref{eq:ch27-e1-contraction},
\[
 \begin{aligned}
 w_{k+1}
 &=P_\omega(\varphi+s_{k+1}f)+d_{k+1}\\
 &\leq P_\omega(\varphi+s_kf)
       +\varepsilon_k+\varepsilon_{k+1}+d_{k+1}
 =w_k .
 \end{aligned}
\]
These are extended-real order comparisons, also on the common pole set;
no difference of two infinite values is formed.  Thus $(w_k)$ is decreasing; and $w_k\leq\varphi+\varepsilon_k+d_k\to
\varphi$.  Hence $w_k\searrow\varphi$ in $\mathcal E^1$, and
\Cref{prop:e1-monotone-continuity} gives
$\MA_\omega(w_k)\to\MA_\omega(\varphi)$ weakly.  Since adding a constant
does not change the Monge--Amp\`ere measure,
$\MA_\omega(w_k)=\MA_\omega(P_\omega(\varphi+s_kf))$, so
$\gamma(s_k)\to\gamma(0)$ along this further subsequence.  The sequential
criterion just stated now gives $\gamma(s)\to\gamma(0)$ as $s\to0$.

\emph{Step 5: conclusion.}
By \eqref{eq:ch27-e1-integrated} and Step~4, $G$ is differentiable at
$0$ with $G'(0)=\gamma(0)=\int f\,\MA_\omega(\varphi)$, which is
\eqref{eq:e1-projected-derivative}.
\end{proof}

\begin{center}
\fbox{\begin{minipage}{0.91\textwidth}
\textbf{Warning (a moving extended-real obstacle).}
The function $\varphi+tf$ may equal $-\infty$, but no undefined
difference is formed in the proof. The envelope inequalities are monotone,
and energy is compared only after adding finite constants. Formula
\eqref{eq:e1-projected-derivative} is obtained from bounded approximants.
\end{minipage}}
\end{center}

\begin{remark}[The measure on the right is $\MA_\omega(\varphi)$, not
$\MA_\omega(P_\omega\varphi)$]
\label{rem:ch27-e1-derivative-measure}
In \eqref{eq:e1-projected-derivative} the measure is that of $\varphi$
itself, because $\varphi$ is already $\omega$-psh and hence
$P_\omega(\varphi)=\varphi$.  This is the special feature of the
finite-energy statement and the reason it is the one used in the variational
argument: the derivative of the projected variation at $t=0$ sees the
Monge--Amp\`ere measure of the \emph{candidate solution}, with no envelope
correction.  The general formula \eqref{eq:energy-envelope-derivative}
degenerates to this one when $f$ is replaced by an $\omega$-psh
obstacle.
\end{remark}

\begin{remark}[Where the hypotheses are used]
\label{rem:ch27-e1-hypotheses}
Orthogonality enters \Cref{cor:e1-projected-variation} only through
\Cref{thm:energy-envelope-differentiability} in Step~1, that is, only at
\emph{continuous} obstacles $\varphi_\alpha+sf$.  We never assert
orthogonality for the upper semicontinuous obstacle $\varphi+sf$, and the
proof does not need it.  This is why the finite-energy statement is a
corollary and not an independent theorem: it transports a statement about
continuous obstacles along a decreasing regularization, and decreasing
regularization is \textup{(D3)}, hence available.  Continuity of envelopes is
used only through that chain.
\end{remark}

\section{The Euler--Lagrange equation}
\label{sec:ch27-euler-lagrange}

\subsection{The variational functional}
\label{sec:ch27-variational-functional}

Let $\mu$ be a positive Radon measure of mass $V$. On the locus where the
integral is defined, set
\begin{equation}\label{eq:bfj-variational-functional}
                  F_\mu(\varphi)
                  :=E_\omega(\varphi)-\int\varphi\,\rd \mu.
\end{equation}
The equality of masses makes $F_\mu$ invariant under constants: by
\eqref{eq:bfj-energy-translation},
\[
 F_\mu(\varphi+c)=E_\omega(\varphi)+cV-\int\varphi\,\rd\mu-c\mu(X^{\An})
 =F_\mu(\varphi).
\]
\Cref{chap:bfj-finite-energy-measures}
will identify the measures for which this functional has the required
finite-energy coercivity and compactness.

Three conventions must be fixed before the functional is used.

\begin{enumerate}
\item \emph{Domain.}  $F_\mu(\varphi)$ is defined for
$\varphi\in\mathcal E^1(X,\omega)$ such that $\int\varphi\,\rd\mu$ is
finite.  The integral is always defined in $[-\infty,+\infty)$, because
$\varphi$ is upper semicontinuous and bounded above; finiteness is a
genuine hypothesis, and it is exactly what $\mu\in\mathcal M^1_V$
guarantees, by \eqref{eq:measure-fenchel-energy}.
\item \emph{Sign.}  $F_\mu$ is concave and is to be \emph{maximized};
\eqref{eq:appJ-energy-variation} and \eqref{eq:appJ-dual-energy} record the
sign convention, and \eqref{eq:dual-energy-abstract} the abstract shape.
Maximizing $F_\mu$ is the same problem as minimizing $-F_\mu$, but
the supergradient calculus of
\Cref{sec:ch27-superdifferential} is written for the concave convention and
should not be mixed with the convex one.
\item \emph{Normalization.}  Since $F_\mu$ is constant on the
$\mathbf R$-orbits $\varphi+\mathbf R$, one normalizes by
$\sup_{X^{\An}}\varphi=0$, which by
\Cref{thm:bfj-psh-compactness} places the competitors in a compact set.
\end{enumerate}

\subsection{Stationarity}
\label{sec:ch27-stationarity}

\begin{proposition}\label{prop:bfj-euler-lagrange}
Assume orthogonality. Suppose
$\varphi\in\mathcal E^1(X,\omega)$ maximizes $F_\mu$ on its domain,
with $\int\varphi\,\rd\mu$ finite. Then
\begin{equation}\label{eq:euler-lagrange-equation}
                           \MA_\omega(\varphi)=\mu.
\end{equation}
Conversely, if $\MA_\omega(\varphi)=\mu$, then $\varphi$ maximizes
$F_\mu$ wherever that functional is defined.
\end{proposition}

\begin{proof}
\emph{Necessity.}  Fix $f\in C^0(X^{\An})$ and define, for
$t\in\mathbf R$,
\[
 H(t):=E_\omega\bigl(P_\omega(\varphi+tf)\bigr)
       -\int(\varphi+tf)\,\rd \mu .
\]
By hypothesis $\int\varphi\,\rd\mu$ is finite, and $\int f\,\rd\mu$ is
finite because $f$ is bounded and $\mu$ has mass $V$.
The bounds
$\varphi-|t|\lVert f\rVert_\infty
 \leq P_\omega(\varphi+tf)
 \leq\varphi+|t|\lVert f\rVert_\infty$
also make the projected potential $\mu$-integrable.
Thus every projected variation is an admissible competitor and $H$ is
real valued, and it is the sum of the concave function
$t\mapsto E_\omega(P_\omega(\varphi+tf))$ and an affine function, hence
concave.

Since $P_\omega(\varphi+tf)\leq\varphi+tf$ and $\mu\geq0$,
\[
 -\int(\varphi+tf)\,\rd\mu\ \leq\
 -\int P_\omega(\varphi+tf)\,\rd\mu ,
\]
so
\[
 H(t)\ \leq\ F_\mu\bigl(P_\omega(\varphi+tf)\bigr)
       \ \leq\ F_\mu(\varphi)=H(0),
\]
the middle inequality by maximality of $\varphi$ and the last equality
because $P_\omega(\varphi)=\varphi$.  Thus $H$ attains its maximum at
the interior point $t=0$.

By \Cref{cor:e1-projected-variation}, $H$ is differentiable at $0$ with
\[
 H'(0)=\int f\,\MA_\omega(\varphi)-\int f\,\rd\mu .
\]
A differentiable function with an interior maximum has vanishing derivative
there, so
\[
 \int f\,\MA_\omega(\varphi)=\int f\,\rd\mu
 \qquad\text{for every }f\in C^0(X^{\An}).
\]
Both sides are positive Radon measures of the same mass $V$---the left by
\Cref{prop:e1-integrability-tail}, the right by hypothesis---so the Riesz
uniqueness statement \Cref{cor:riesz-uniqueness} gives
\eqref{eq:euler-lagrange-equation}.

\emph{Sufficiency.}  Suppose $\MA_\omega(\varphi)=\mu$ and let
$\psi\in\mathcal E^1(X,\omega)$ be such that $\int\psi\,\rd\mu$ is
finite.  The right-hand inequality of the two-sided sandwich
\eqref{eq:ch27-two-sided-sandwich} gives
\[
 E_\omega(\psi)-E_\omega(\varphi)
 \ \leq\ \int(\psi-\varphi)\,\MA_\omega(\varphi)
 =\int(\psi-\varphi)\,\rd \mu,
\]
all terms being finite by \Cref{lem:ch27-mutual-integrability}.  Rearranging
gives $F_\mu(\psi)\leq F_\mu(\varphi)$.
\end{proof}

\begin{remark}
\label{rem:ch27-euler-lagrange-not-existence}
The proposition is an Euler--Lagrange statement, not an existence theorem.
It does not assert that $F_\mu$ attains its supremum, that the maximizer is
continuous, or that every measure of mass $V$ has finite energy. These are
separate questions in \Crefrange{chap:bfj-finite-energy-measures}{chap:bfj-continuity}.
\end{remark}

\begin{corollary}
\label{cor:ch27-euler-lagrange-supergradient}
Assume orthogonality and let $\varphi\in\mathcal E^1(X,\omega)$ with
$\int\varphi\,\rd\mu$ finite.  Then $\varphi$ maximizes $F_\mu$ if and
only if $\mu$ is a supergradient of $E_\omega$ at $\varphi$, in the
sense that
\begin{equation}\label{eq:ch27-energy-supergradient}
 E_\omega(\psi)\ \leq\ E_\omega(\varphi)+\int(\psi-\varphi)\,\rd\mu
 \qquad\text{for every }\psi\in\mathcal E^1(X,\omega)
 \text{ with }\int\psi\,\rd\mu>-\infty,
\end{equation}
and this holds if and only if $\MA_\omega(\varphi)=\mu$.
\end{corollary}

\begin{proof}
Rearranging \eqref{eq:ch27-energy-supergradient} is exactly
$F_\mu(\psi)\leq F_\mu(\varphi)$ for all admissible $\psi$, which is
maximality.  If $\MA_\omega(\varphi)=\mu$, then
\eqref{eq:ch27-energy-supergradient} is the right-hand inequality of
\eqref{eq:ch27-two-sided-sandwich}.  Conversely, maximality gives
\eqref{eq:euler-lagrange-equation} by
\Cref{prop:bfj-euler-lagrange}.
\end{proof}

\begin{remark}[Two Legendre transforms, one equation]
\label{rem:ch27-two-legendre}
\Cref{cor:ch27-euler-lagrange-supergradient} identifies the Euler--Lagrange
equation with the equality case of the Fenchel inequality
\eqref{eq:fenchel-energy} for the pair $(E_\omega,E^\vee_\omega)$, by
\Cref{prop:concave-biconjugation}\,(3).  On the other hand
\Cref{prop:ch27-conjugate-identification}\,(3) identifies orthogonality at
$f$ with the equality case of the Fenchel inequality for the pair
$(F,F^\vee)=(E_\omega\circ P_\omega,E^\vee_\omega)$.  Both statements
concern the same dual functional $E^\vee_\omega$; they differ in the
primal variable, a potential in the first case and an obstacle in the
second.  The passage between them is $P_\omega$, and the assertion that
it is compatible with the duality is precisely
\Cref{cor:ch27-orthogonality-iff-differentiability}.  This is the cleanest
statement of what the variational method actually needs.
\end{remark}

\begin{remark}[What replaces uniqueness here]
\label{rem:ch27-uniqueness-pointer}
Neither \Cref{prop:bfj-euler-lagrange} nor
\Cref{cor:ch27-euler-lagrange-supergradient} asserts that the maximizer is
unique modulo constants.  That statement is the Calabi--Yau uniqueness
theorem, quoted with its source in
\Cref{cor:bfj-original-continuous-solution}; see also
\Cref{rem:ch27-equality-case} for the exact point at which the formal
argument stops.  What the two statements above do give is that the
\emph{set} of maximizers is convex and that all of them have the same
Monge--Amp\`ere measure, namely $\mu$.
\end{remark}

\section{Where envelope continuity is needed}
\label{sec:ch27-where-continuity}

\index{envelope!continuity, where used}
This section is an audit.  It lists every place in
\Cref{chap:bfj-ma,chap:bfj-capacity,chap:bfj-energy} and in the present
chapter at which \Cref{thm:bfj-envelope-regularity} is used, distinguishes
structural from local uses, and records what is known when its hypotheses
fail.  The hypothesis signature being invoked is fixed once, and it is the
one boxed in \Cref{sec:ch23-continuity}:

\begin{center}
\fbox{\begin{minipage}{0.91\textwidth}
\textbf{Hypothesis signature of \Cref{thm:bfj-envelope-regularity},
as invoked here.}
$K$ complete discretely valued, residue characteristic $0$; $X$ smooth
projective over $K$; $\omega$ a closed model $(1,1)$-form with
$\{\omega\}$ ample; the obstacle continuous.  Conclusion: $P_\omega(f)$
is a uniform limit of $\omega$-psh model functions, in particular
continuous, and $P^{\mathcal D}_\omega(f)=P_\omega(f)$ by
\Cref{cor:ch23-envelope-by-model-functions}.  Discreteness is used in the
model theory of \Cref{chap:bfj-model-forms}; smoothness and residue
characteristic zero for log resolutions and for cofinality of SNC models;
ampleness in \Cref{lem:ch23-envelope-finiteness} and in the rationalization
step.
\end{minipage}}
\end{center}

\subsection{The structural use, and why it dominates}
\label{sec:ch27-structural-use}

\begin{enumerate}
\item[\textup{(EC1)}]
\emph{This construction of the Monge--Amp\`ere operator uses it.}
\Cref{thm:bounded-ma-extension} constructs $\MA_\omega$ on bounded psh
functions as the unique monotone extension of the model intersection
measures, and the construction proceeds by choosing \emph{decreasing model
regularizations} of each entry.  Their existence is
\Cref{thm:bfj-decreasing-regularization}, which is deduced from
\Cref{thm:bfj-envelope-regularity} in
\Cref{cor:ch23-decreasing-regularization-pointer}.  Everything downstream
inherits this dependence: capacity \eqref{eq:capacity-definition}, the
maximum principle \eqref{eq:bounded-maximum-principle}, the energy of a
bounded potential (\Cref{prop:bounded-energy-identities}), the
nonpluripolar measure \eqref{eq:nonpluripolar-ma}, the class
$\mathcal E^1$, and hence every statement of this chapter.
\end{enumerate}

This one item swamps all the others, and it should be stated plainly: in the
original BFJ architecture, \emph{continuity of envelopes is a prerequisite
for the existence of the Monge--Amp\`ere operator itself}, not a convenience
in a later proof.  A reader who wishes to work over a field where
\Cref{thm:bfj-envelope-regularity} is unavailable must first replace
\textup{(EC1)}, and only then worry about the local uses below.  This is
exactly the order in which
\Cref{thm:gjkm-positive-characteristic-cy} proceeds in equal characteristic
$p$: envelope continuity by test ideals first, then cofinality of SNC
models by embedded resolution, then the whole calculus, and only then
orthogonality.

\subsection{The local uses}
\label{sec:ch27-local-uses}

Granting \textup{(EC1)}, the direct uses below fall into five groups.
In four groups the role is identical: an
object built from $P_\omega$ must be fed back into a statement whose
hypothesis is a \emph{continuous} function.

\begin{enumerate}
\item[\textup{(EC2)}]
\Cref{prop:ch27-F-elementary}\,(5)--(6): to know that $P_\omega(f)$ is
an admissible continuous argument of $F$ and an admissible competitor in the supremum defining $F(f)$, that is, that the
supremum is attained.
\item[\textup{(EC3)}]
\Cref{prop:ch27-supergradients-on-contact}: to substitute $g=P_\omega(f)$
into the supergradient inequality \eqref{eq:ch27-supergradient-inequality},
which is quantified over $g\in C^0(X^{\An})$.  Without it one cannot
conclude that a supergradient is carried by the contact set, and the
equivalence $(4)\Rightarrow(1)$ of
\Cref{prop:ch27-orthogonality-equivalences} fails to be available.
\item[\textup{(EC4)}]
\Cref{prop:ch27-orthogonality-uniform-limits}: to know that
$f_j-P_\omega(f_j)$ and $f-P_\omega(f)$ are continuous, so that the
joint convergence lemma \Cref{lem:joint-weak-uniform}, which tests weakly
convergent measures against uniformly convergent \emph{continuous}
functions, applies.
\item[\textup{(EC5)}]
\Cref{prop:ch27-conjugate-identification}\,(2): to approximate
$P_\omega(f)$ uniformly by $\omega$-psh \emph{model} functions, thereby
comparing the continuous-obstacle conjugate \eqref{eq:ch27-F-conjugate}
with the model-test conjugate \eqref{eq:dual-energy-measure}.  This is the
one use in which the full strength of
\Cref{thm:bfj-envelope-regularity}---uniform approximation by model
functions, not merely continuity---is needed.
\item[\textup{(EC6)}]
\Cref{prop:differentiability-iff-orthogonality}: to know that
$g=P_\omega(f)-f$ is a continuous direction, so that the hypothesis
\eqref{eq:energy-envelope-derivative} may be applied to it.  This is the
converse implication of
\Cref{cor:ch27-orthogonality-iff-differentiability}, and it is the only
place where envelope continuity is used to \emph{produce} orthogonality
rather than to manipulate it.
\end{enumerate}

Beyond \textup{(EC1)} there is one further, external use:

\begin{enumerate}
\item[\textup{(EC7)}]
\Cref{thm:orthogonality-provenance}\,(2), that is,
\Cref{thm:bggjkm-remove-algebraizability}: continuity of the envelope is an
input to its application in the orthogonality deduction, since the
derivative theorem must be applied at the continuous semipositive envelope.
The derivative theorem itself starts at a continuous semipositive metric
and does not assume envelope continuity for arbitrary obstacles.
This distinction is explicit in \cite[Theorems~5.4.3 and~6.3.2]{BGGJKM20}.
\end{enumerate}

\subsection{Where it is \emph{not} needed}
\label{sec:ch27-not-needed}

It is equally important to record the statements of this chapter whose
proofs use \emph{nothing} from
\Cref{thm:bfj-envelope-regularity} except \textup{(EC1)}.  Each of these
would survive verbatim in any setting in which a Monge--Amp\`ere operator
and an energy with the properties of
\Cref{chap:bfj-ma,chap:bfj-energy} are available by other means.

\begin{center}
\small
\begin{tabular}{p{.42\textwidth}|p{.50\textwidth}}
\toprule
\textbf{Statement}&\textbf{Only input beyond \textup{(EC1)}}\\
\midrule
\Cref{prop:ch27-energy-polynomial}
 &the cocycle and the Hodge index inequality\\
\Cref{lem:ch27-mutual-integrability}
 &the mass tail \eqref{eq:e1-mass-tail} and portmanteau\\
\Cref{thm:e1-energy-differentiability}
 &Lagrange interpolation and monotone measures\\
\Cref{cor:ch27-two-sided-sandwich}
 &concavity of a function of one variable\\
\Cref{thm:ch27-superdifferential}
 &Hahn--Banach and Riesz\\
\Cref{lem:ch27-uniform-ma}
 &the monotone extension \Cref{thm:bounded-ma-extension}\\
\Cref{thm:ch27-model-direction-differentiability}
 &orthogonality at two obstacles\\
\Cref{thm:energy-envelope-differentiability}
 &the model case plus density and contraction\\
\Cref{cor:e1-projected-variation}
 &decreasing regularization, that is \textup{(EC1)} again\\
\Cref{prop:bfj-euler-lagrange}
 &\Cref{cor:e1-projected-variation} and Riesz uniqueness\\
\Cref{thm:ch27-legendre-orthogonality}
 &convex analysis only; no non-Archimedean input at all\\
\bottomrule
\end{tabular}
\end{center}

The table is the practical content of the dependency diagram
\eqref{eq:ch27-dependency-diagram}: the analytic core of this chapter is
independent of the envelope theorem, and the two genuine couplings are
\textup{(EC1)}, which is upstream of everything, and the converse
implication \textup{(EC6)}.

\subsection{What is known in positive residue characteristic}
\label{sec:ch27-positive-characteristic}

\index{residue characteristic!positive, envelope and orthogonality}
We now record the status of the two theorems---continuity of envelopes and
orthogonality---outside residue characteristic zero.  Nothing in this
subsection is used in a proof; every entry is quoted with its source and its
hypotheses, in accordance with the standard fixed in
\Cref{sec:ch23-later}.

\paragraph{Curves over an arbitrary complete field.}
For ample line bundles on smooth projective geometrically connected
curves, semipositive metric-envelope continuity holds without discreteness
or residue-characteristic restrictions:
\cite[Theorem~1.1]{GJKM19} and \Cref{thm:ch23-gjkm-curves}.
This assertion concerns rational-model competitors.
Orthogonality on a curve is the complementary-slackness statement
\Cref{thm:ch23-model-orthogonality-curve} at model level, upgraded by
\Cref{prop:ch27-curve-orthogonality} and
\Cref{cor:ch27-curve-orthogonality-global} under the semistability
hypothesis discussed in \Cref{rem:ch27-curve-base-change}.  These are distinct scopes: the arbitrary-field envelope theorem is a
metric statement, while the orthogonality deduction just given has an
additional model-envelope comparison hypothesis.
At model level the graph Laplacian makes the orthogonality calculation
elementary, but
the mechanism that makes them elementary---retraction to a skeleton
preserves semipositivity---is false for $n\geq2$, by the toric
counterexample of Burgos Gil and Sombra recalled in
\Cref{rem:ch23-curve-specialness}.

\paragraph{Equal characteristic $p$.}
Let $K$ be complete discretely valued with
$\operatorname{char}K=p>0$, let $X$ be smooth projective of dimension
$n$, of geometric origin from a $d$-dimensional family over a perfect
field $k$ in the sense of \Cref{def:ch23-geometric-origin}, and assume
resolution of singularities over $k$ in dimension $d+n$.  Then
continuity of the envelope holds:
\Cref{thm:ch23-gjkm-positive-characteristic}, that is,
\cite[Theorem~8.2]{GJKM19}, with asymptotic test ideals replacing asymptotic
multiplier ideals throughout.  If in addition embedded resolution holds in
dimension $d+n$, so that projective SNC models are cofinal and the entire
calculus of \Cref{chap:bfj-psh,chap:bfj-ma,chap:bfj-energy} is available,
then orthogonality holds as well: this is
\cite[Theorem~9.2]{GJKM19}, whose proof uses the volume theorem
\cite[Theorem~6.3.2]{BGGJKM20}, and it is reconstructed in the proof of
\Cref{thm:gjkm-positive-characteristic-cy}.  The logical order there is
exactly \eqref{eq:ch27-dependency-diagram} with the multiplier-ideal input
replaced by the test-ideal input; in particular envelope continuity is
again \emph{upstream} of orthogonality and not the other way round.

For a smooth projective \emph{surface} of geometric origin from a
one-dimensional family over a perfect field of characteristic $p$, both
hypotheses are theorems---resolution and embedded resolution in dimension
three over a perfect field, by Cossart--Piltant as quoted in
\Cref{cor:ch23-gjkm-surfaces}---so continuity and orthogonality hold
unconditionally.  This is the unconditional non-toric higher-dimensional case supplied by
the sources quoted here; no exhaustive classification of other special
families is claimed.

\paragraph{Mixed characteristic.}
Here the situation is different in kind, because there is no Frobenius on
$K$-schemes and hence no test-ideal substitute of the equal-characteristic
type; see \Cref{rem:ch23-gjkm-hypotheses}.  What is available is
\Cref{thm:fgk-mixed-characteristic-cy}, that is, \cite{FGK25}: for $K$
complete discretely valued of mixed characteristic $(0,p)$ and $X$
smooth projective, \emph{assuming} resolution and embedded resolution for
projective $K^\circ$-models of $X$, the Monge--Amp\`ere equation is
solvable for a measure supported on the skeleton of a projective SNC model,
the envelope step being run with perturbation-friendly mixed-characteristic
test ideals.  Both hypotheses are open in general.
\Cref{ex:ch23-outside-scope} exhibits an explicit smooth projective
threefold over $\mathbf Q_p$ for which no theorem proved or cited in this
book gives an unconditional envelope-continuity theorem.  The
orthogonality routes cited here likewise do not apply unconditionally.

\paragraph{Summary table.}

\begin{center}
\small
\begin{tabular}{p{.24\textwidth}|p{.14\textwidth}|p{.27\textwidth}|p{.27\textwidth}}
\toprule
\textbf{Field and variety}&\textbf{$\dim X$}
 &\textbf{Envelope continuity}&\textbf{Orthogonality}\\
\midrule
cdv, res.\ char.\ $0$; $X$ smooth projective; class ample
 &any
 &\Cref{thm:bfj-envelope-regularity}
 &\Cref{thm:orthogonality-provenance}\,(2)\\
cdv, res.\ char.\ $0$; $X$ in addition algebraizable
 &any
 &\Cref{thm:bfj-envelope-regularity}
 &\Cref{thm:orthogonality-provenance}\,(1), by asymptotic orthogonality\\
any complete nontrivially valued $K$; ample $L$, curve geometrically connected
 &$1$
 &\Cref{thm:ch23-gjkm-curves}
 &\Cref{cor:ch27-curve-orthogonality-global} under its positive-class
   standing hypotheses and semistability;
   \Cref{rem:ch27-curve-base-change}\\
cdv, $\operatorname{char}K=p$, geometric origin, resolution in
 $d+n$
 &any
 &\Cref{thm:ch23-gjkm-positive-characteristic}
 &\cite[Theorem~9.2]{GJKM19}, with embedded resolution as well\\
cdv, $\operatorname{char}K=p$, surface of geometric origin from a
 $1$-dimensional family
 &$2$
 &\Cref{cor:ch23-gjkm-surfaces}, unconditional
 &unconditional, same source\\
cdv, mixed characteristic $(0,p)$, smooth projective $X$
 &any
 &\cite[Theorem~7.1]{FGK25}, with resolution for projective models
 &volume route with resolution and embedded resolution;
   \cite[Theorem~8.1]{FGK25}\\
smooth polarized toric, rational invariant model obstacle, standing BFJ field
 &any
 &rational polyhedral envelope
 &invariant-envelope \Cref{thm:ch27-toric-orthogonality}\\
\bottomrule
\end{tabular}
\end{center}

\begin{remark}[Reading the last row]
\label{rem:ch27-toric-row}
The toric row is not a counterexample to the pattern of the others; it is a
statement about a subcategory.  \Cref{thm:ch27-toric-orthogonality} computes
the \emph{invariant} envelope, and \Cref{rem:ch27-toric-scope} explains that
this is a hypothesis and not a reduction.  What the row does show is that
the difficulty of orthogonality is not a difficulty about fields or residue
characteristics: it is a difficulty about the absence of a Legendre duality
on the space of psh metrics, and it disappears exactly when such a duality is
present.
\end{remark}

\begin{warning}{Do not transport a continuity theorem across its dimension}
The curve theorem \Cref{thm:ch23-gjkm-curves} holds over every complete
nontrivially valued field in its stated ample metric scope, and it is
tempting to combine it with a fibration
argument to obtain continuity in higher dimension.  No such argument is
known, and the obstruction is identified in
\Cref{rem:ch23-curve-specialness}: retraction to a skeleton preserves
semipositivity on a curve and fails to do so already on a toric surface.
The same warning applies to orthogonality, whose curve proof
\Cref{thm:ch23-model-orthogonality-curve} is a finite linear program and has
no higher-dimensional analogue, by
\Cref{rem:ch23-model-slackness-fails}.
\end{warning}

\section{Curve and toric obstacle problems, carried to numbers}
\label{sec:ch27-worked}

Every theorem of this chapter is now verified on explicit data.  The first
computation differentiates the envelope energy along a curve skeleton and
watches the Monge--Amp\`ere measure move; the second verifies orthogonality
and the derivative formula in a toric case by Legendre duality; the third
solves the Euler--Lagrange equation for explicit measures on a graph.

\subsection{A directional derivative on a curve skeleton}
\label{sec:ch27-curve-numbers}

\begin{example}[The five-vertex chain: envelope, measure, and derivative]
\label{ex:ch27-curve-derivative}
We reuse verbatim the data of \Cref{ex:ch23-curve-envelope-numbers}.  Thus
$K$ is discretely valued with $v(\pi)=1$, $X=\mathbf P^1_K$, and
$\mathcal X$ is the regular model with reduced special fibre the chain
\[
 \mathcal X_s=E_0+E_1+E_2+E_3+E_4,
 \qquad b_i=1,\qquad (E_i\cdot E_{i+1})=1,
\]
so that the dual graph $\Gamma$ is the path
$e_0-e_1-e_2-e_3-e_4$ with all edge weights $w_{i,i+1}=1$.  We take
$L=\cO_{\mathbf P^1}(2)$ with the model extension whose residual degrees
are
\[
 (c_0,c_1,c_2,c_3,c_4)=(1,0,0,0,1),
 \qquad \sum_ic_i=2=\deg_XL=V,
\]
and $\omega:=c_1(\mathcal L,\lVert\cdot\rVert_{\mathcal L})$, which is
semipositive because all $c_i\geq0$.

\emph{The energy on the graph.}  For a model function with vertex vector
$v\in\mathbf R^5$, \eqref{eq:model-energy} in dimension one reads
\[
 E_\omega(v)=\int v\,\omega+\tfrac12\int v\,\ddc v
 =\sum_ic_iv_i-\tfrac12\sum_{i}w_{i,i+1}(v_{i+1}-v_i)^2 ,
\]
the second equality because
$\ddc\varphi_v=\sum_i(\Delta_\Gamma v)_i\delta_{e_i}$ by
\eqref{eq:ch20-graph-laplacian} and because
$\sum_iv_i(\Delta_\Gamma v)_i=-\sum_{i<j}w_{ij}(v_i-v_j)^2$.  Writing
$d_i:=v_{i+1}-v_i$ for $i=0,1,2,3$, this is
\begin{equation}\label{eq:ch27-chain-energy}
 E_\omega(v)=v_0+v_4-\tfrac12\bigl(d_0^2+d_1^2+d_2^2+d_3^2\bigr),
\end{equation}
and by \Cref{lem:ch23-curve-residual-degree} the Monge--Amp\`ere measure is
$\sum_iA_i(v)\delta_{e_i}$ with
\begin{equation}\label{eq:ch27-chain-A}
 A_0=1+d_0,\quad A_1=d_1-d_0,\quad A_2=d_2-d_1,\quad A_3=d_3-d_2,
 \quad A_4=1-d_3,
\end{equation}
of total mass $2$.  Semipositivity is $A_i\geq0$ for all $i$, that is,
\begin{equation}\label{eq:ch27-chain-slope-window}
 -1\ \leq\ d_0\leq d_1\leq d_2\leq d_3\ \leq\ 1 :
\end{equation}
convexity along the path, with the two boundary slopes confined to
$[-1,1]$ by the residual degrees $c_0=c_4=1$.

\emph{The obstacle and its envelope.}  Take
\[
 u=(0,-\tfrac32,-3,-\tfrac32,0),
\]
the piecewise affine ``V'' with vertex value $-3$ at $e_2$.  As computed
in \Cref{ex:ch23-curve-envelope-numbers},
\[
 P:=P_\omega(u)=(-1,-2,-3,-2,-1),
 \qquad
 \MA_\omega(P)=2\,\delta_{e_2},
\]
and by \Cref{rem:ch23-model-dependence} this model-level envelope is the
global one.  From \eqref{eq:ch27-chain-energy},
\[
 E_\omega(P)=(-1)+(-1)-\tfrac12(1+1+1+1)=-4 .
\]
On the skeleton contact occurs only at $e_2$, and $\int(u-P)\MA_\omega(P)=0$, which is
\eqref{eq:orthogonality} verified in numbers.

\emph{First direction: a bump at the contact point.}  Let
$g^{(2)}:=(0,0,1,0,0)$, a model function.  The perturbed obstacle is
$u+tg^{(2)}=(0,-\tfrac32,-3+t,-\tfrac32,0)$.  For $t\leq\tfrac12$ we
claim
\[
 P_\omega\bigl(u+tg^{(2)}\bigr)=P+t\,(1,1,1,1,1).
\]
Feasibility: the vector $P+t\mathbf 1$ satisfies the obstacle constraints
$-1+t\leq0$, $-2+t\leq-\tfrac32$, $-3+t\leq-3+t$, and their mirror
images, and it has the same slope vector $d=(-1,-1,1,1)$ as $P$, hence
the same $A$-vector $(0,0,2,0,0)\geq0$.  Maximality: for any competitor
$v$, convexity at $e_1$ gives $2v_1\leq v_0+v_2$, while $A_0\geq0$
gives $v_0\leq v_1+1$; combining with $v_2\leq-3+t$,
\[
 2v_1\ \leq\ (v_1+1)+(-3+t),
 \qquad\text{so}\qquad v_1\leq t-2,\quad v_0\leq t-1,
\]
and symmetrically $v_3\leq t-2$, $v_4\leq t-1$.  Hence
\[
 E_\omega\bigl(P_\omega(u+tg^{(2)})\bigr)
 =E_\omega(P)+tV=-4+2t
 \qquad\bigl(t\leq\tfrac12\bigr),
\]
by the translation rule \eqref{eq:bfj-energy-translation}, so the derivative
at $t=0$ is $2$.  On the other side,
\[
 \int g^{(2)}\,\MA_\omega(P)=1\cdot2=2 .
\]
The two agree, verifying
\eqref{eq:ch27-model-direction-derivative}.  Note that the measure did not
move at all here: raising the obstacle at the unique contact point simply
translates the envelope.

\emph{Second direction: a bump away from the contact point.}  Let
$g^{(0)}:=(1,0,0,0,0)$, so that $\int g^{(0)}\MA_\omega(P)=A_0(P)=0$ and
the predicted derivative at $t=0$ is $0$.  Here the answer is genuinely
piecewise, and it is worth computing on the whole line.  The obstacle is
$u+tg^{(0)}=(t,-\tfrac32,-3,-\tfrac32,0)$.  Using
\eqref{eq:ch27-chain-slope-window} and maximizing coordinatewise:
\[
 P_\omega\bigl(u+tg^{(0)}\bigr)
 =\begin{cases}
   (t,\ t+1,\ t+2,\ t+3,\ t+4), & t\leq-5,\\[2pt]
   \bigl(t,\ \tfrac{t-3}2,\ -3,\ -2,\ -1\bigr), & -5\leq t\leq-1,\\[2pt]
   (-1,-2,-3,-2,-1), & t\geq-1 .
  \end{cases}
\]
Indeed, on the middle range the binding constraints are $v_0\leq t$,
$v_2\leq-3$ and $d_0\leq d_1$; maximizing $v_1=t+d_0$ subject to
$d_0\leq d_1$ and $t+d_0+d_1\leq-3$ forces
$d_0=d_1=\tfrac{-3-t}2$, which lies in $[-1,1]$ exactly for
$-5\leq t\leq-1$, and then $d_2=d_3=1$ is optimal.  On the left range all
slopes are pinned at $1$; on the right range the obstacle at $e_0$ is
slack and the envelope is the unperturbed $P$.

The corresponding measures, from \eqref{eq:ch27-chain-A}, are
\[
 \MA_\omega\bigl(P_\omega(u+tg^{(0)})\bigr)
 =\begin{cases}
   2\,\delta_{e_0}, & t\leq-5,\\[2pt]
   \tfrac{-1-t}2\,\delta_{e_0}+\tfrac{5+t}2\,\delta_{e_2},
     & -5\leq t\leq-1,\\[2pt]
   2\,\delta_{e_2}, & t\geq-1,
  \end{cases}
\]
each of total mass $2$, the mass migrating continuously from $e_2$ to
$e_0$ as $t$ decreases from $-1$ to $-5$.  The energies, from
\eqref{eq:ch27-chain-energy}, are
\[
 E_\omega\bigl(P_\omega(u+tg^{(0)})\bigr)
 =\begin{cases}
   2t+2, & t\leq-5,\\[2pt]
   t-2-\tfrac{(3+t)^2}4, & -5\leq t\leq-1,\\[2pt]
   -4, & t\geq-1 .
  \end{cases}
\]
The three expressions agree at $t=-5$ (value $-8$) and at $t=-1$
(value $-4$), and differentiating gives
\[
 \frac{\rd}{\rd t}
 E_\omega\bigl(P_\omega(u+tg^{(0)})\bigr)
 =\begin{cases}
   2, & t\leq-5,\\[2pt]
   \tfrac{-1-t}2, & -5\leq t\leq-1,\\[2pt]
   0, & t\geq-1,
  \end{cases}
\]
while, on the other side,
\[
 \int g^{(0)}\,\MA_\omega\bigl(P_\omega(u+tg^{(0)})\bigr)
 =A_0\bigl(P_\omega(u+tg^{(0)})\bigr)
 =\begin{cases}
   2,& t\leq-5,\\[2pt]
   \tfrac{-1-t}2,& -5\leq t\leq-1,\\[2pt]
   0,& t\geq-1 .
  \end{cases}
\]  The two functions of $t$
coincide identically:
\eqref{eq:ch27-model-direction-derivative} holds not only at $t=0$ but at
every parameter, as the theorem asserts.

\emph{Four checks against the general theory.}
\begin{enumerate}
\item The derivative is nonincreasing in $t$---it equals $2$, then
$\tfrac{-1-t}2$ which decreases from $2$ to $0$, then $0$---in
accordance with concavity of $t\mapsto F(u+tg^{(0)})$
(\Cref{prop:ch27-F-elementary}\,(4)).
\item Its absolute value never exceeds
$V\lVert g^{(0)}\rVert_\infty=2$, and the bound is attained exactly on the
range $t\leq-5$ where the obstacle at $e_0$ is so low that the envelope
is pinned to it; this is the Lipschitz estimate
\Cref{prop:ch27-F-elementary}\,(3) being sharp.
\item The derivative is continuous, as
\Cref{thm:ch27-model-direction-differentiability} asserts, even though the
envelope itself is only piecewise affine in $t$ and the measure changes
support at $t=-1$ and $t=-5$.
\item Orthogonality holds at every $t$.  For instance at $t=-3$ the
envelope is $(-3,-3,-3,-2,-1)$, the obstacle is
$(-3,-\tfrac32,-3,-\tfrac32,0)$, the skeletal contact set is $\{e_0,e_2\}$, and
$\MA_\omega=\delta_{e_0}+\delta_{e_2}$ charges only those two vertices,
so $\int(u+tg^{(0)}-P_\omega(\cdot))\MA_\omega(\cdot)=0$.  This is
\Cref{prop:ch27-curve-orthogonality} in numbers.
\end{enumerate}
\end{example}

\begin{remark}[What the example shows about the two-sided estimate]
\label{rem:ch27-curve-example-lesson}
Follow the proof of
\Cref{thm:ch27-model-direction-differentiability} on the second direction at
$t_0=-1$, the parameter at which the measure begins to move.  The upper
estimate \eqref{eq:ch27-upper-estimate} uses $\mu_{t_0}=2\delta_{e_2}$ and
gives $\int g^{(0)}\rd\mu_{t_0}=0$; the lower estimate
\eqref{eq:ch27-lower-estimate} uses $\mu_t$ for nearby $t$, and gives
$\tfrac{-1-t}2\to0$.  The squeeze is exactly the continuity of
$t\mapsto\mu_t$ in the weak topology, which here is visible as a
continuous transfer of mass between two atoms.  Had we possessed only the
upper estimate, we would have obtained $D^+\leq0\leq D^-$ and no
conclusion, in accordance with
\Cref{ex:ch27-corner-without-orthogonality}.
\end{remark}

\subsection{A toric verification by Legendre duality}
\label{sec:ch27-toric-numbers}

\begin{example}[Orthogonality and the derivative on $\mathbf P^1$, by
Legendre transform]
\label{ex:ch27-toric-legendre}
Let $K$ be as in the standing hypotheses, $X=\mathbf P^1_K$
with its standard torus---a smooth projective toric variety---and
$L=\cO(2)$, so that
$N_{\mathbf R}=M_{\mathbf R}=\mathbf R$, $V=(L^1)=2$, and in the
notation of \Cref{sec:ch27-toric-transport} we may take the body
\[
 Q=[0,2]\subset M_{\mathbf R},
 \qquad
 h_Q(x)=\sup_{m\in[0,2]}mx=\max\{0,2x\} .
\]
Take as reference the invariant semipositive model metric with weight
$w_0:=h_Q$, which is rational piecewise affine and convex, and put
$\omega=c_1(L,\lVert\cdot\rVert_0)$; then
$w_0^*=0$ on $Q$ and $E_\omega(0)=0$.

\emph{The obstacle.}  Let
\[
 g(x):=\max\{0,1-|x|\},
 \qquad
 f:=-g,
\]
both invariant model functions (piecewise affine with integer slopes and
compact support), and set $w_f:=w_0+f=h_Q-g$, explicitly
\[
 w_f(x)=
 \begin{cases}
  0,&x\leq-1,\\
  -1-x,&-1\leq x\leq0,\\
  3x-1,&0\leq x\leq1,\\
  2x,&x\geq1 .
 \end{cases}
\]
Its slopes are $0,-1,3,2$ in that order, so $w_f$ is not convex, and
$\lVert w_f-h_Q\rVert_\infty=1$.

\emph{The Legendre transform.}  For $m\in[0,2]$, computing the supremum of
$mx-w_f(x)$ branch by branch gives the four candidate values
\[
 -m\ \ (\text{at }x=-1),
 \qquad 1\ \ (\text{at }x=0),
 \qquad 1\ \ (\text{at }x=0),
 \qquad m-2\ \ (\text{at }x=1),
\]
whence
\[
 w_f^*(m)=\max\{-m,\,1,\,m-2\}=1
 \qquad\text{for all }m\in[0,2],
\]
attained \emph{only} at $x=0$.  Therefore, by
\Cref{thm:ch27-legendre-orthogonality}\,(1),
\[
 u:=\operatorname{conv}(w_f)=w_f^{**}
   =\sup_{m\in[0,2]}(mx-1)=h_Q(x)-1,
\]
and the contact set is
$C=\{u=w_f\}=\{0\}$, since $u<w_f$ strictly elsewhere (for instance
$u(-1)=-1<0=w_f(-1)$ and $u(1)=1<2=w_f(1)$).

\emph{Orthogonality.}  The convex function $u$ has a single slope jump,
from $0$ to $2$, at $x=0$, so
\[
 \MA_{\mathbf R}(u)=2\,\delta_0,
 \qquad
 \MA_{\mathbf R}(u)(\mathbf R)=2=\operatorname{vol}_M(Q),
\]
in accordance with \eqref{eq:appD-real-ma-mass}; the measure is carried by
the contact set $\{0\}$, and
\[
 \int_{\mathbf R}(w_f-u)\,\MA_{\mathbf R}(u)
 =\bigl(w_f(0)-u(0)\bigr)\cdot2=(-1-(-1))\cdot2=0,
\]
which is \eqref{eq:ch27-convex-orthogonality}.  Transported by
\Cref{thm:ch27-toric-orthogonality}, and using
$\trop_*\MA_\omega(P^{\mathrm{inv}}_\omega(f))=1!\,\MA_{\mathbf R}(u)$, this
is \eqref{eq:ch27-toric-orthogonality} in numbers:
$P^{\mathrm{inv}}_\omega(f)=-1$, $\MA_\omega(P^{\mathrm{inv}}_\omega(f))=
2\delta_{x_0}$ with $x_0$ the Gauss point, and the contact integral
vanishes.

\emph{The derivative, in closed form.}  Move the obstacle in the direction
$g$, so that $w_{f+tg}=h_Q+(t-1)g$.  Exactly the same branchwise
computation gives, for every $t\in\mathbf R$,
\[
 (w_{f+tg})^*(m)=\max\{-m,\ 1-t,\ m-2\},
 \qquad m\in[0,2].
\]
By \Cref{cor:ch27-toric-energy} with the reference $w_0=h_Q$, for which
$w_0^*=0$ on $Q$, the envelope energy is
\[
 \Lambda(t):=E_\omega\bigl(P^{\mathrm{inv}}_\omega(f+tg)\bigr)
 =-n!\int_Q(w_{f+tg})^*\,\rd\operatorname{Leb}_M
 =-\int_0^2(w_{f+tg})^*(m)\,\rd m ,
\]
the second equality because $n=1$.
Since $\max\{-m,m-2\}=-m$ on $[0,1]$ and $=m-2$ on $[1,2]$, with
minimum value $-1$ at $m=1$, the plateau $1-t$ is the maximum exactly
on the interval $[t-1,3-t]\cap[0,2]$.  Hence
\[
 \Lambda(t)=
 \begin{cases}
  2t-2, & t\leq1,\\[2pt]
  2(t-1)-(t-1)^2, & 1\leq t\leq2,\\[2pt]
  1, & t\geq2 .
 \end{cases}
\]
The three expressions agree at $t=1$ (value $0$) and at $t=2$ (value
$1$), and
\[
 \Lambda'(t)=
 \begin{cases}
  2,& t\leq1,\\
  4-2t,& 1\leq t\leq2,\\
  0,& t\geq2 .
 \end{cases}
\]

\emph{The derivative, from the measure.}  For $1\leq t\leq2$ the
maximizer of $mx-w_{f+tg}(x)$ is $x=-1$ for $m\in(0,t-1)$, $x=0$ for
$m\in(t-1,3-t)$, and $x=1$ for $m\in(3-t,2)$.  By
\Cref{thm:ch27-legendre-orthogonality}\,(3) the Monge--Amp\`ere measure of
the envelope $u_t$ is therefore the pushforward of Lebesgue measure on
$[0,2]$,
\[
 \MA_{\mathbf R}(u_t)
 =(t-1)\,\delta_{-1}+(4-2t)\,\delta_0+(t-1)\,\delta_1,
\]
of total mass $2$ for every $t\in[1,2]$.  Consequently
\[
 \int g\,\MA_{\mathbf R}(u_t)
 =g(-1)(t-1)+g(0)(4-2t)+g(1)(t-1)
 =0+(4-2t)+0=4-2t,
\]
which is exactly $\Lambda'(t)$.  At $t=0$ this reads
$\Lambda'(0)=2=\int g\,\MA_{\mathbf R}(u)=g(0)\cdot2$, which is
\eqref{eq:ch27-convex-derivative}, and, after the toric dictionary,
\eqref{eq:energy-envelope-derivative}.

\emph{Three normalization checks.}  First, $\Lambda(0)=-2$, while
$P^{\mathrm{inv}}_\omega(f)=-1$ as a relative weight, so
$E_\omega(-1)=E_\omega(0)-1\cdot V=-2$ by
\eqref{eq:bfj-energy-translation}: the two agree, confirming the mass
$V=2$.  Second, $|\Lambda'|\leq V\lVert g\rVert_\infty=2$ throughout,
with equality on $t\leq1$.  Third, for $1<t<2$ the perturbed weight $w_{f+tg}=h_Q+(t-1)g$ has
slopes $0,\ t-1,\ 3-t,\ 2$ in that order, and this sequence is
nondecreasing precisely for $1\leq t\leq2$.  So on that whole range
$w_{f+tg}$ is already convex, $u_t=w_{f+tg}$, and the contact set is all
of $N_{\mathbf R}$; the three displayed equalities
$w_{f+tg}(-1)=0=u_t(-1)$, $w_{f+tg}(0)=t-1=u_t(0)$,
$w_{f+tg}(1)=2=u_t(1)$ are three instances of that, not a proof that the
contact set is $\{-1,0,1\}$.  Orthogonality is therefore automatic here,
$f-P$ vanishing identically.  What the computation does exhibit is the
genuinely three-atom measure
$\MA_{\mathbf R}(u_t)=(t-1)\delta_{-1}+(4-2t)\delta_0+(t-1)\delta_1$, whose
atoms are the three slope jumps $t-1$, $4-2t$, $t-1$ at $-1,0,1$, of
total mass $2$.  The envelope is a \emph{proper} envelope only outside
$[1,2]$: for $t<1$ the contact set is $\{0\}$, and for $t>2$ the
slope sequence again fails to be monotone.
\end{example}

\begin{example}[A toric contact set]
\label{ex:ch27-toric-contact-prototype}
For an invariant obstacle, use the invariant envelope of \eqref{eq:ch27-invariant-envelope}.  Then $P^{\mathrm{inv}}_\omega(f)$ becomes the convex envelope with
the recession behavior prescribed by the polarization polytope. On an open
region where the envelope lies strictly below the obstacle it is affine in
the one-dimensional model, and its Alexandrov measure vanishes there. In
higher dimension the Alexandrov measure of the convex envelope is supported
on the contact set in the corresponding obstacle theorem. Consequently
\[
       \int(f-P^{\mathrm{inv}}_\omega(f))\MA_\omega(P^{\mathrm{inv}}_\omega(f))=0.
\]
This is the real convex prototype for orthogonality. It explains the theorem
but does not replace the global model-volume proof.
\end{example}

The last example is now a theorem, in two forms:
\Cref{thm:ch27-legendre-orthogonality} on the convex side and
\Cref{thm:ch27-toric-orthogonality} after transport, with the scope
restrictions of \Cref{rem:ch27-toric-scope}.

\begin{example}[The same computation in two variables]
\label{ex:ch27-toric-two-variables}
Take $X=\mathbf P^1_K\times\mathbf P^1_K$ with
$L=\cO(2,2)$, so that $n=2$ and one may take $Q=[0,2]^2$ with
$h_Q(x_1,x_2)=\max\{0,2x_1\}+\max\{0,2x_2\}$.  Consider the separable
obstacle
\[
 w_f(x_1,x_2)=h_Q(x_1,x_2)-g(x_1)-g(x_2),
 \qquad g(s)=\max\{0,1-|s|\}
\]
as in \Cref{ex:ch27-toric-legendre}.  Conjugation turns a separable sum into
a separable sum, so
$w_f^*(m_1,m_2)=1+1=2$ on $Q$, the supremum being attained only at
$(0,0)$; hence
\[
 \begin{aligned}
 u&=\operatorname{conv}(w_f)=h_Q-2,
 &C&=\{(0,0)\},\\
 \MA_{\mathbf R}(u)
   &=\operatorname{vol}_M(Q)\,\delta_{(0,0)}
   =4\,\delta_{(0,0)},
 \end{aligned}
\]
the last equality by \Cref{ex:appD-polyhedral-atom}, since
$\partial u(0,0)=[0,2]^2$.  Orthogonality is immediate:
$w_f(0,0)=-2=u(0,0)$, so the whole mass sits on the contact set.

The normalizations now carry the factor $n!=2$, and it is worth watching
it.  The total mass is
$V=n!\operatorname{vol}_M(Q)=2\cdot4=8$ by
\eqref{eq:appJ-toric-degree}, whereas the Alexandrov mass is
$\operatorname{vol}_M(Q)=4$; the two differ exactly by $n!$, which is
the test recorded in the warning at the end of
\Cref{sec:appD-real-ma}.  By \Cref{cor:ch27-toric-energy} with reference
$w_0=h_Q$,
\[
 \Lambda(0)=E_\omega\bigl(P^{\mathrm{inv}}_\omega(f)\bigr)
 =-n!\int_Qw_f^*\,\rd\operatorname{Leb}_M=-2\cdot2\cdot4=-16 ,
\]
which agrees with the direct computation
$E_\omega(-2)=E_\omega(0)-2V=-16$, the envelope having relative weight
$u-h_Q=-2$.  For the direction
$g\oplus g:(x_1,x_2)\mapsto g(x_1)+g(x_2)$ one finds
\[
 \Lambda'(0)=n!\int(g\oplus g)\,\MA_{\mathbf R}(u)
 =2\bigl(g(0)+g(0)\bigr)\cdot4=16,
\]
which saturates the Lipschitz bound
$V\lVert g\oplus g\rVert_\infty=8\cdot2=16$, as it must, because the
direction attains its maximum precisely on the contact set.  Compare
\Cref{ex:ch23-toric-orthogonality}, where the same two-variable envelope is
computed on the domain $(-2,2)^2$ with the unreflected obstacle.
\end{example}

\subsection{Euler--Lagrange verifications}
\label{sec:ch27-euler-lagrange-numbers}

\begin{example}[Solving the two-vertex equation]
\label{ex:ch27-two-vertex-euler-lagrange}
In \Cref{ex:two-vertex-ma}, let the target be
$\mu=q_0\delta_{x_0}+q_1\delta_{x_1}$, with
$q_i\geq0$ and $q_0+q_1=m_0+m_1$. For vertex values $a,b$,
\[
 \MA_\omega(\varphi)=\mu
 \quad\Longleftrightarrow\quad
 b-a=q_0-m_0=m_1-q_1.
\]
Thus the solution is unique modulo constants. Directly differentiating
\[
 F_\mu(a,b)
 =m_0a+m_1b-\frac12(a-b)^2-q_0a-q_1b
\]
gives the same linear equation. This is
\Cref{prop:bfj-euler-lagrange} reduced to graph calculus.
\end{example}

\begin{example}[The Euler--Lagrange equation on the five-vertex chain]
\label{ex:ch27-euler-lagrange-chain}
Keep the data of \Cref{ex:ch27-curve-derivative}: the path $\Gamma$ with
vertices $e_0,\dots,e_4$, unit edge weights, residual degrees
$c=(1,0,0,0,1)$ and $V=2$.  Let
\[
 \mu=\sum_{i=0}^4q_i\delta_{e_i},
 \qquad q_i\geq0,\qquad \sum_iq_i=2 .
\]
We solve $\MA_\omega(\varphi)=\mu$ and verify
\Cref{prop:bfj-euler-lagrange} completely.

\emph{Reduction to slopes.}  With $d_j=v_{j+1}-v_j$ as before and
$S_j:=\sum_{i>j}(c_i-q_i)$, one has, for any $v\in\mathbf R^5$,
\[
 \sum_i(c_i-q_i)v_i=\sum_{j=0}^3S_jd_j,
\]
because $\sum_i(c_i-q_i)=2-2=0$ kills the constant term when $v_i$ is
expanded as $v_0+\sum_{j<i}d_j$.  Extend the graph formula for $F_\mu$
as a quadratic function of every vertex vector; on the psh vectors it is
the actual variational functional.  Hence
\begin{equation}\label{eq:ch27-chain-variational}
 F_\mu(v)=E_\omega(v)-\int v\,\rd\mu
 =\sum_{j=0}^3\Bigl(S_jd_j-\tfrac12d_j^2\Bigr),
\end{equation}
using \eqref{eq:ch27-chain-energy}.  This is a sum of four independent
concave parabolas: $F_\mu$ is maximized over all of $\mathbf R^5$
precisely when
\[
 d_j=S_j\quad(j=0,1,2,3),
 \qquad
 \max F_\mu=\tfrac12\sum_{j=0}^3S_j^2 ,
\]
the maximizer being unique modulo the additive constant $v_0$, as it must
be by translation invariance of $F_\mu$.

\emph{The maximizer solves the equation.}  Substituting $d_j=S_j$ into
\eqref{eq:ch27-chain-A} and using $S_j-S_{j-1}=q_j-c_j$,
\[
 A_0=1+S_0=q_0,\qquad
 A_i=S_i-S_{i-1}=q_i\ \ (i=1,2,3),\qquad
 A_4=1-S_3=q_4 ,
\]
so $\MA_\omega(v)=\mu$ exactly.  In particular $A_i=q_i\geq0$, so the
unconstrained maximizer is automatically $\omega$-psh: on this graph the
unconstrained optimum satisfies the positivity constraints.  Some of them
are equalities when $q_i=0$.  This is the
finite-dimensional shadow of the statement that a maximizer of $F_\mu$
solves $\MA_\omega(\varphi)=\mu$ with no envelope correction
(\Cref{rem:ch27-e1-derivative-measure}).

\emph{Two explicit measures.}
For $\mu=2\delta_{e_2}$ one has $c-q=(1,0,-2,0,1)$, hence
$S=(-1,-1,1,1)$, so
\[
 v=(a+2,\ a+1,\ a,\ a+1,\ a+2),
 \qquad
 \max F_\mu=\tfrac12(1+1+1+1)=2 .
\]
Taking $a=-3$ recovers the envelope
$P=(-1,-2,-3,-2,-1)$ of \Cref{ex:ch27-curve-derivative}, whose
Monge--Amp\`ere measure was computed there to be $2\delta_{e_2}$; and
indeed $F_\mu(P)=E_\omega(P)-\int P\,\rd\mu=-4-2(-3)=2$, as predicted.

For $\mu=\delta_{e_1}+\delta_{e_3}$ one has $c-q=(1,-1,0,-1,1)$, hence
$S=(-1,0,0,1)$, so
\[
 \begin{aligned}
 v&=(a+1,\ a,\ a,\ a,\ a+1),\\
 \MA_\omega(v)&=\delta_{e_1}+\delta_{e_3},\\
 \max F_\mu&=\tfrac12(1+0+0+1)=1 .
 \end{aligned}
\]
One checks directly: $d=(-1,0,0,1)$ gives
$A=(1-1,\,0+1,\,0,\,1,\,1-1)=(0,1,0,1,0)$, of total mass $2$; and
$F_\mu(v)=\bigl[(a+1)+(a+1)-\tfrac12(1+0+0+1)\bigr]-\bigl[a+a\bigr]=1$.

\emph{The converse direction.}  Both maximizers satisfy the supergradient
inequality \eqref{eq:ch27-energy-supergradient} directly: for any
$v'\in\mathbf R^5$, \eqref{eq:ch27-chain-variational} gives
\[
 F_\mu(v')-F_\mu(v)
 =-\tfrac12\sum_{j}\bigl(d_j'-S_j\bigr)^2\ \leq\ 0,
\]
which is \Cref{cor:ch27-euler-lagrange-supergradient} with the defect
expressed as a Dirichlet energy.  The quantity
$\tfrac12\sum_j(d_j'-S_j)^2$ is precisely the directed functional
\eqref{eq:directed-dirichlet-functional}
$J_\omega(v,v')$ of \Cref{chap:bfj-finite-energy-measures} in this
finite-dimensional model, and its vanishing forces $v'-v$ to be constant.
\end{example}

\begin{example}[A measure that no model function on $\mathcal X$ can
solve]
\label{ex:ch27-measure-off-vertices}
Keep $\Gamma$ as above and let $x$ be the midpoint of the edge
$[e_1,e_2]$, so its path coordinate is $r=3/2$, and put $\mu=2\delta_x$.
No model function determined on $\mathcal X$ has measure $\mu$, since
\eqref{eq:ch23-curve-model-measure} puts every such measure on the original
vertices.

Blowing up the node corresponding to this edge introduces an exceptional
component of multiplicity two and subdivides the modulus-length-one edge
at its midpoint.  The resulting model is regular with SNC special-fibre
support, but the fibre is not reduced.  Consequently the unit-length graph
formula cannot be reused without its length and multiplicity conventions.
Here the solution can be written directly on the subdivided path:
\[
 v(r)=-\min\{r,3/2\}-\min\{4-r,5/2\},
 \qquad 0\leq r\leq4.
\]
Its slope is $-1$ before $x$ and $1$ after $x$.  Pull it back by the
skeletal retraction, so it is constant on attached branches.
The outgoing slopes cancel the two reference atoms at $e_0,e_4$ and
produce the atom $2\delta_x$ at the breakpoint, by
\eqref{eq:ch25-curve-ma}.  Thus it is a psh model function on the
refined model and solves the equation.  The energy uses modulus lengths,
as in \Cref{prop:ch26-graph-energy}; its graph integral is invariant under
this subdivision.

The equation concerns $X^{\An}$ rather than a fixed model, so models must
be allowed to grow with the data.  The support hypothesis in
\Cref{thm:bfj-supported-cy} asks for one dual complex carrying $\mu$;
it does not require the support to consist of its original vertices.
\end{example}

\begin{example}[Failure of the naive argument, visible on the chain]
\label{ex:ch27-naive-failure-chain}
Return to \Cref{ex:ch27-curve-derivative} at $t=0$, with obstacle
$u=(0,-\tfrac32,-3,-\tfrac32,0)$ and envelope
$P=(-1,-2,-3,-2,-1)$.  The heuristic criticized in
\Cref{rem:model-slackness-heuristic} would say: at $e_0$ the envelope lies
strictly below the obstacle ($-1<0$), so one should be able to raise it
there, and hence $A_0(P)$ must vanish.  Here that conclusion happens to be
\emph{true}, $A_0(P)=0$; but the heuristic gets the mechanism backwards.
Raising $v_0$ alone is blocked by the constraint $A_0\geq0$ at $e_0$
itself, which is tight: $A_0(P)=1+(P_1-P_0)=0$, and replacing $P_0$ by
$P_0+\varepsilon$ gives $A_0=-\varepsilon<0$, while it \emph{relaxes}
convexity at $e_1$, $A_1$ rising to $+\varepsilon$.  What actually
proves $A_0(P)=0$ is Step~4 of the proof of
\Cref{thm:ch23-model-orthogonality-curve}, which raises the coordinate and
checks that \emph{all five} constraints survive---possible only because each
$A_i$ is affine in $v$ with nonnegative off-diagonal coefficients.  In
dimension $n\geq2$ the analogous check fails at the first step, as
\Cref{rem:ch23-model-slackness-fails} explains.
\end{example}

\section{Chapter dictionary}
\label{sec:ch27-dictionary}

\subsection{From the objects of this chapter to their use}
\label{sec:ch27-dict-forward}

\begin{center}
\small
\begin{tabular}{p{.28\textwidth}|p{.30\textwidth}|p{.34\textwidth}}
\toprule
\textbf{Analytic object}&\textbf{Variational meaning}
 &\textbf{Established at}\\
\midrule
$\MA_\omega(\varphi)$ for
 $\varphi\in\mathcal E^1$
 &supergradient of $E_\omega$ at $\varphi$
 &\Cref{cor:ch27-two-sided-sandwich},
  \Cref{cor:ch27-euler-lagrange-supergradient}\\
one-sided derivative of $E_\omega$ along a psh segment
 &integration of the increment against the endpoint measure
 &\Cref{thm:e1-energy-differentiability},
  \Cref{cor:ch27-far-endpoint}\\
concavity of $E_\omega$ on a segment
 &the polynomial \eqref{eq:ch27-energy-derivative-binomial} with
  nonincreasing derivative
 &\Cref{prop:ch27-energy-polynomial}\\
$F=E_\omega\circ P_\omega$
 &concave, monotone, $V$-Lipschitz, translation equivariant
 &\Cref{prop:ch27-F-elementary}\\
$\partial F(f)$
 &compact convex set of measures of mass $V$ on the contact set
 &\Cref{thm:ch27-superdifferential},
  \Cref{prop:ch27-supergradients-on-contact}\\
orthogonality at $f$
 &$\MA_\omega(P_\omega f)\in\partial F(f)$; Fenchel equality for
  $(F,E^\vee_\omega)$
 &\Cref{prop:ch27-orthogonality-equivalences}\\
Gateaux differentiability of $F$
 &$\partial F(f)$ is a singleton
 &\Cref{thm:ch27-superdifferential}\,(4),
  \Cref{cor:ch27-singleton}\\
contact set $C_f$
 &support of every supergradient
 &\eqref{eq:ch27-contact-set},
  \Cref{prop:ch27-supergradients-on-contact}\\
convex envelope, Legendre transform
 &the maximizer $\nabla u^*(m)$ lies in $C_f$
 &\Cref{thm:ch27-legendre-orthogonality}\,(2)\\
Alexandrov measure as a pushforward
 &the derivative of the conjugate integral
 &\Cref{thm:ch27-legendre-orthogonality}\,(3), (4)\\
complementary slackness $A_i(P)(u_i-P_i)=0$
 &orthogonality on a curve model
 &\Cref{thm:ch23-model-orthogonality-curve},
  \Cref{prop:ch27-curve-orthogonality}\\
$F_\mu(\varphi)=E_\omega(\varphi)-\int\varphi\,\rd\mu$
 &the variational functional; stationarity is
  $\MA_\omega(\varphi)=\mu$
 &\eqref{eq:bfj-variational-functional},
  \Cref{prop:bfj-euler-lagrange}\\
$E^\vee_\omega(\mu)$
 &concave conjugate of $F$; Fenchel equality at supergradients
 &\Cref{prop:ch27-conjugate-identification}\\
\bottomrule
\end{tabular}
\end{center}

\subsection{The four statements that must never be merged}
\label{sec:ch27-dict-four}

\begin{center}
\small
\begin{tabular}{p{.22\textwidth}|p{.36\textwidth}|p{.34\textwidth}}
\toprule
\textbf{Name}&\textbf{Statement}&\textbf{Status here}\\
\midrule
Energy differentiability
 &$\tfrac{\rd}{\rd t}E_\omega((1-t)\varphi+t\psi)\big|_{0+}
    =\int(\psi-\varphi)\MA_\omega(\varphi)$
 &proved: \Cref{thm:e1-energy-differentiability}\\
Envelope differentiability
 &$\tfrac{\rd}{\rd t}E_\omega(P_\omega(f+tg))\big|_{0}
    =\int g\,\MA_\omega(P_\omega f)$
 &proved from orthogonality:
  \Cref{thm:energy-envelope-differentiability}\\
Orthogonality
 &$\int(f-P_\omega f)\MA_\omega(P_\omega f)=0$
 &imported for $n\geq2$:
  \Cref{thm:orthogonality-provenance}\\
Envelope continuity
 &$P_\omega f\in C^0$, uniform limit of model psh functions
 &proved in \Cref{chap:bfj-envelopes}:
  \Cref{thm:bfj-envelope-regularity}\\
\bottomrule
\end{tabular}
\end{center}

Given the energy calculus, the first needs no orthogonality input.
The second and the third are
equivalent (\Cref{cor:ch27-orthogonality-iff-differentiability}).  The
fourth is upstream of all of them, through
\textup{(EC1)} of \Cref{sec:ch27-structural-use}, and it appears again as a
hypothesis in the historical proofs of the third.

\subsection{Numbers to remember}
\label{sec:ch27-dict-numbers}

\begin{center}
\small
\begin{tabular}{p{.36\textwidth}|p{.26\textwidth}|p{.30\textwidth}}
\toprule
\textbf{Quantity}&\textbf{Value}&\textbf{Where}\\
\midrule
Beta weight in the energy cocycle
 &$\binom nk\!\int_0^1\!(1-s)^{n-k}s^k\rd s=\tfrac1{n+1}$
 &Step~4 of \Cref{prop:ch27-energy-polynomial}\\
Lipschitz constant of $E_\omega\circ P_\omega$
 &$V$
 &\Cref{prop:ch27-F-elementary}\,(3)\\
mass of every element of $\partial F(f)$
 &$V$
 &\Cref{thm:ch27-superdifferential}\,(2)\\
mutual integrability bound
 &$\int\psi\,\MA_\omega(\varphi)\geq
    E_\omega(\psi)+(n+1)E_\omega(\varphi)$
 &\eqref{eq:ch27-mutual-integrability}\\
envelope on the five-vertex chain
 &$(-1,-2,-3,-2,-1)$, energy $-4$
 &\Cref{ex:ch27-curve-derivative}\\
its Monge--Amp\`ere measure
 &$2\delta_{e_2}$, mass $2=\deg L$
 &\Cref{ex:ch27-curve-derivative}\\
derivative in the direction $(0,0,1,0,0)$
 &$2$
 &\Cref{ex:ch27-curve-derivative}\\
derivative in the direction $(1,0,0,0,0)$
 &$0$ for $t\geq-1$, $\tfrac{-1-t}2$ for
  $-5\leq t\leq-1$, $2$ for $t\leq-5$
 &\Cref{ex:ch27-curve-derivative}\\
toric conjugate of the model obstacle
 &$w_f^*\equiv1$ on $[0,2]$
 &\Cref{ex:ch27-toric-legendre}\\
toric envelope energy
 &$\Lambda(t)=2t-2$ for $t\leq1$
 &\Cref{ex:ch27-toric-legendre}\\
toric measure at $t\in[1,2]$
 &$(t-1)\delta_{-1}+(4-2t)\delta_0+(t-1)\delta_1$
 &\Cref{ex:ch27-toric-legendre}\\
maximum of $F_\mu$ on the chain
 &$\tfrac12\sum_{j}S_j^2$, with $S_j=\sum_{i>j}(c_i-q_i)$
 &\eqref{eq:ch27-chain-variational}\\
$\max F_{2\delta_{e_2}}$, $\max F_{\delta_{e_1}+\delta_{e_3}}$
 &$2$ and $1$
 &\Cref{ex:ch27-euler-lagrange-chain}\\
\bottomrule
\end{tabular}
\end{center}

\section*{Arithmetic heights and equidistribution}
\markright{Arithmetic heights and equidistribution}
\addcontentsline{toc}{section}{Arithmetic heights and equidistribution}

The variational equation also has an arithmetic incarnation: a
directional derivative determines the distribution of algebraic points
whose heights approach an essential minimum. The moving functional is
then a minimum of heights, rather than the energy of a local envelope.
Balla\"y--Sombra's adelic approximation theorem
\cite[Theorems~6.9 and~6.17]{BallaySombra26Published} makes this
connection without requiring equality in Zhang's lower bound.
The comparison is useful here because it exhibits another reason to
control a first-order error \emph{relative to a geometric scale}.

\paragraph{A separate arithmetic setting.}
In this discussion, let $F$ be a number field and $Y/F$ a normal
integral projective variety of dimension $d\geq1$. These hypotheses
replace the standing local BFJ hypotheses for this discussion only.
At a place $v$ use the completed algebraic closure $\mathbf C_v$,
its analytification $Y_v^{\An}$, and the weight
\[
 n_v=\frac{[F_v:\mathbf Q_v]}{[F:\mathbf Q]}.
\]
Use the usual absolute values and the product formula; the weights
at the Archimedean places sum to one.
A continuous adelic $\mathbf R$-Cartier divisor
$\overline D=(D,(g_v)_v)$ has Green functions
$g_v=-\log\|s_D\|_v$, with a common model outside finitely many
places. In particular, $g_v$ is a minus-log-norm function, not a
minus-log-squared-norm function. It is continuous after adding the
logarithm of a local defining equation.

Assume the geometric divisor $D$ is big. Write
\[
 \mu=\mu^{\mathrm{ess}}(\overline D)
   =\sup_{Z\subsetneq Y}\ \inf_{x\in(Y\setminus Z)(\overline F)}
                                      h_{\overline D}(x),
\]
where $Z$ ranges over proper closed subsets. A sequence $(x_j)$ is
generic if it eventually avoids each such subset, and small if
$h_{\overline D}(x_j)\to\mu$.
For a semipositive adelic divisor $\overline Q$ put
$\mu^{\mathrm{abs}}(\overline Q)=\inf_{x\in Y(\overline F)}
h_{\overline Q}(x)$. Here semipositivity means continuous local psh
metrics and a geometrically nef divisor. Adelic neffness additionally
requires $\mu^{\mathrm{abs}}\geq0$.
The Archimedean unit divisor $[\infty]$ has underlying divisor zero,
Green function one at every Archimedean place, and zero at the other
places; adding $c[\infty]$ adds $c$ to all heights.

\paragraph{The approximation scale.}
For big geometric $\mathbf R$-divisors $P,A$ define
\[
 r(P;A)=\sup\{\lambda\in\mathbf R:P-\lambda A\text{ is big}\}.
\]
This inradius is finite and positive. A semipositive approximation
of $\overline D$ is a normal modification $\pi:Y'\to Y$ and a
semipositive $\overline Q$ on $Y'$ with $Q$ big and
$\pi^*\overline D-\overline Q$ adelically pseudo-effective.
Arithmetically big means positive arithmetic volume of small sections.
An adelic divisor is pseudo-effective if its sum with every
arithmetically big adelic divisor is arithmetically big.
This is an arithmetic condition; it does not mean pointwise
domination of every local Green function.

Take a sequence $(\pi_m,\overline Q_m)$ of such approximations,
with finite absolute minima, and set
\[
 \delta_m=\mu-\mu^{\mathrm{abs}}(\overline Q_m),\qquad
 r_m=r(Q_m;\pi_m^*D).
\]
The hypothesis of the cited Theorem~6.9 is
\[
                         \frac{\delta_m}{r_m}\longrightarrow0.
\]
The target metric need not itself be semipositive. The conclusion
is Gateaux differentiability of the essential minimum in every
continuous adelic direction $\overline E$, with
\[
 \left.\frac{\rd}{\rd t}\right|_{t=0}
       \mu^{\mathrm{ess}}(\overline D+t\overline E)
 =\lim_{m\to\infty}
 \frac{(\overline Q_m^{\,d}\cdot\pi_m^*\overline E)
       -d\mu(Q_m^{d-1}\cdot\pi_m^*E)}{(Q_m^d)}.
\]
The first intersection is arithmetic and has $d+1$ factors;
the second is geometric and has $d$ factors. The denominator is
the approximant's geometric mass $(Q_m^d)$, which can tend to
zero. Replacing it by $\operatorname{vol}(D)$ changes the statement.

At every place $v$, the probability measures
\[
 \nu_{m,v}=(\pi_m)_*
       \frac{c_1(\overline Q_{m,v})^{\wedge d}}{(Q_m^d)}
\]
have a weak limit $\nu_v$, independent of the chosen approximation
sequence satisfying the criterion. For every generic small sequence,
its Galois-orbit probability measures converge to $\nu_v$.
For a zero-geometric-divisor metric direction
$\overline E=\overline0^{\,f}$ supported at $v$, the formula becomes
\[
 \mathrm D_{\overline0^{\,f}}\mu^{\mathrm{ess}}(\overline D)
                         =n_v\int_{Y_v^{\An}}f\,\rd\nu_v.
\]
The same weight $n_v$ multiplies the orbit average in the point
height. Thus the derivative recovers the probability measure after
division by $n_v$; the arithmetic weight is not part of its mass.

\paragraph{Why an inradius enters the proof.}
After the height normalization
$\widetilde{\overline Q}_m=\overline Q_m-
\mu^{\mathrm{abs}}(\overline Q_m)[\infty]$, the approximants are
adelically nef. Arithmetic Siu and Hodge estimates bound the
one-sided slope error, after a fixed comparison of geometric
reference divisors, by
\[
                  \frac{\delta_m}{\lambda}
                     +C\frac{\lambda}{r_m},\qquad
                          0<\lambda<c r_m.
\]
Here $C,c>0$ depend only on the fixed data.
For $\delta_m>0$, take
$\lambda_m=r_m\sqrt{\delta_m/r_m}$; if $\delta_m=0$, take
$\lambda_m=r_m/\sqrt m$. These choices eventually lie in the stated interval, and both errors
tend to zero. Applying the
slope comparison to opposite directions identifies the derivative.
This is the arithmetic counterpart of the two-sided slope squeeze
used earlier in this chapter. Convergence of $\delta_m$ alone does
not control the squeeze when the geometric scale also shrinks.

This paragraph describes the proof mechanism, not a new proof of
the imported arithmetic inputs. A nef comparison divisor must stay
nef during the normalization: a downward constant shift can destroy
that property. A simultaneous upward constant shift of the target
and all approximants preserves their difference, $\delta_m$, the
geometric inradii, and the normalized nef approximants.
The arithmetic Siu inequality used here is a different statement
from the homogeneous Siu decomposition discussed in
\Cref{chap:bj-homogeneity-bdivisors}; it supplies no repair of that problem.

\paragraph{A quadratic roof with strict Zhang inequality.}
For $Y=\mathbf P^1_{\mathbf Q}$ and $D=\mathcal O(1)$ the moment
polytope is $[0,1]$, with lattice length one. Choose the continuous
semipositive toric metric whose Archimedean roof is
\[
                 \vartheta_\infty(x)=1-(x-\tfrac12)^2,
                        \qquad 0\leq x\leq1,
\]
and whose finite-place roofs are zero. The existence of the metric
is the concave Legendre metric--roof dictionary of
\cite[Section~8.1]{BallaySombra26Published}.
The global roof is the displayed quadratic. The toric height and
intersection formulas give
\[
 \mu^{\mathrm{ess}}(\overline D)=1,\qquad
 \mu^{\mathrm{abs}}(\overline D)=\frac34,\qquad
 \frac{(\overline D^{\,2})}{2(D)}
       =\int_0^1\vartheta_\infty(x)\,\rd x
       =\frac{11}{12}<1.
\]
Indeed, the maximum is at $1/2$, the minimum is at the endpoints,
and the integral of $(x-1/2)^2$ is $1/12$.
Thus Zhang's lower bound is strict; the constant approximation
$\overline Q=\overline D$ does not give the vanishing-defect
criterion.

For $t=1-\varepsilon$ with $0<\varepsilon<1/4$, the upper-level set is
\[
 \Gamma_t=\{\vartheta_\infty\geq t\}
     =[\tfrac12-\sqrt\varepsilon,\tfrac12+\sqrt\varepsilon].
\]
Its width, and its convex inradius relative to $[0,1]$, are both
$2\sqrt\varepsilon$. Consequently
\[
              \frac{1-t}{r(\Gamma_t;[0,1])}
                            =\frac{\sqrt\varepsilon}{2}\longrightarrow0.
\]
This is the \emph{wide-roof} criterion of
\cite[Theorem~8.7 and Proposition~A.3]{BallaySombra26Published}.
The local sup-gradients at $x=1/2$ are all zero, and hence balanced.
The limiting measure at infinity is probability Haar measure on
the unit circle. At each finite place it is the Dirac probability
measure at the unit Gauss point, characterized on Laurent
polynomials by
\[
                    \left|\sum_j a_jz^j\right|=\max_j|a_j|_v.
\]
This example connects an explicit convex level-set calculation to
small-point equidistribution even when Zhang equality fails.

\paragraph{A smooth metric whose small points select a circle.}
The small-point limit can differ from the curvature measure of the
given metric even when that metric is smooth and positive.
Take $Y=\mathbf P^1_{\mathbf Q}$, $D=[\infty]$, the Fubini--Study
metric at infinity, and the canonical metric of
$(\mathbf P^1_{\mathbf Z},[\infty])$ at every finite place.
In the affine coordinate $z$, the Green functions are
\[
 g_\infty(z)=\frac12\log(1+|z|^2),\qquad
 g_v(z)=\log\max\{1,|z|_v\}\quad(v\text{ finite}).
\]
The factor $1/2$ records the ordinary norm used in the height.
The corresponding roof is
\[
 \vartheta_\infty(x)
    =-\frac12\bigl(x\log x+(1-x)\log(1-x)\bigr),
       \qquad 0\leq x\leq1,
\]
with $0\log0=0$, and the finite-place roofs vanish.
It has unique maximum $\frac12\log2$ at $x=1/2$.
The toric height formula gives
\[
 \mu^{\mathrm{ess}}(\overline D)=\frac12\log2,
 \qquad
 \frac{(\overline D^{\,2})}{2(D)}
        =\int_0^1\vartheta_\infty(x)\,\rd x
        =\frac14<\frac12\log2.
\]
One can check the essential minimum directly on the torus.
For $r>0$,
$\frac12\log(1+r^2)-\frac12\log r\geq\frac12\log2$,
while
$\log\max\{1,r\}-\frac12\log r\geq0$.
Summing over all places of a field of definition and using the
weighted product formula gives the lower bound for every nonzero
algebraic $z$. Roots of unity attain it and form a generic sequence.
The two boundary points have height zero; their heights do not
determine the essential minimum.

For every generic small sequence the limiting measures are
unit-circle probability Haar measure at infinity and unit-Gauss
Dirac measures at finite places. This is the standard example in
\cite[Theorem~6.4 and Example~6.5]{BGPRS19Published}; the canonical
finite-place models are part of its hypotheses.
The original Archimedean curvature measure, however, is
\[
 c_1\bigl(\mathcal O(1),\|\cdot\|_\infty\bigr)
   =\ddc(2g_\infty)
   =\frac{1}{\pi}\frac{\rd x\,\rd y}{(1+|z|^2)^2},
                  \qquad z=x+\rmi y.
\]
Here the complex operator and squared-norm convention are those
of \Cref{sec:appJ-complex-ddc}.
This probability measure gives the unit circle mass zero and
the unit disc mass $1/2$, whereas the small-point limit gives
the circle mass one. Thus the measures supplied by the arithmetic
variation must be identified from the limiting functional;
smooth positivity of the initial metric does not identify them.

\paragraph{A translated abelian polarization with Zhang equality.}
An abelian example makes the height normalization equally explicit.
Let $B/\mathbf Q$ be an abelian variety of dimension $g\geq1$, and
let $\overline L$ be an ample symmetric line bundle with its
rigidified canonical adelic metric. For $q\in B(\mathbf Q)$, write
$h_q=\widehat h_L(q)$, let $t_q$ denote translation by $q$, and put
\[
 \overline T=t_q^*\overline L-
                 \pi^*(\overline L|_q),\qquad
 \overline R=\overline T-\overline L,
\]
where $\pi:B\to\operatorname{Spec}\mathbf Q$ is the structure map.
The constant metrized line $\overline L|_q$ fixes the rigidification;
its height contribution is $h_q$. Thus
\[
 h_{\overline T}(x)=\widehat h_L(x+q)-h_q.
\]
The canonical height is nonnegative, and a generic sequence of
torsion points translated by $-q$ attains the lower bound. Hence
\[
 \mu^{\mathrm{abs}}(\overline T)
       =\mu^{\mathrm{ess}}(\overline T)=-h_q.
\]

Geometrically, $R$ is algebraically trivial and $(T^g)=(L^g)$.
Theorem~2.5 of \cite{CL00Published} gives
\[
 (\overline R^{\,2}\cdot\overline L^{\,g-1})
       =-\frac2g(L^g)h_q,
\]
and all other mixed arithmetic terms of total degree $g+1$ vanish.
Expanding the full power therefore gives
\[
 (\overline T^{\,g+1})
   =\binom{g+1}{2}
      (\overline R^{\,2}\cdot\overline L^{\,g-1})
   =-(g+1)(L^g)h_q.
\]
The binomial coefficient belongs to the full power. Dividing by
$(g+1)(T^g)$ gives $-h_q$, so Zhang's lower bound is an equality
in this example, including when the normalized height is negative.

The canonical metric on $R$ is flat at infinity, and constant
metric shifts preserve curvature. Consequently
$c_1(\overline T_\infty)=c_1(\overline L_\infty)$, and its normalized
top power is probability Haar measure on $B(\mathbf C)$.
For the constant approximation $\overline Q=\overline T$, the
inradius is $r(T;T)=1$ and the height defect is zero. The preceding
approximation criterion therefore recovers the original normalized
curvature measures as the small-point limits. The two examples
show why the normalized arithmetic intersection and the essential
minimum must be compared before identifying those measures.

\paragraph{A curve over a higher-dimensional geometric base.}
The dimension of the base enters the comparison with successive minima.
Let $C$ be a smooth projective curve of genus $g>1$ over an
algebraically closed field $k$, let $J$ be its Jacobian, and assume
$\rho(C\times J)=3$. Choose a degree-one line bundle $\alpha$ with
$(2g-2)\alpha=\omega_C$, and let $\theta$ be the symmetric principal
theta polarization associated with this choice. Write
$i_\alpha(x)=x-\alpha$ and
\[
 \alpha_1=p_1^*\alpha,\qquad \theta_2=p_2^*\theta,\qquad
 f_{m,n}(x,y)=m i_\alpha(x)+ny.
\]
Use the addition convention for the Poincare class $Q$, so that
\[
 f_{m,n}^*\theta=g m^2\alpha_1+n^2\theta_2+mnQ
 \qquad(m,n\in\mathbf Z).
\]
The numerical classes $\alpha_1,\theta_2,Q$ form a basis.
Guo--Yuan's cone calculation gives
\[
 \operatorname{Psef}(C\times J)=\operatorname{Nef}(C\times J)
 =\{a\alpha_1+b\theta_2+cQ:a,b\geq0,\ ab\geq gc^2\};
\]
see \cite[Proposition~2.2 and Theorem~2.4]{GY26RelativeMinimaV2}.

Put $K=k(J)$, $\pi=p_2$, and
$\mathcal L=f_{1,1}^*\theta=g\alpha_1+\theta_2+Q$.
This line bundle is nef on the entire model, and its restriction
to the generic curve $C_K$ has degree $g$.
With base polarization $\theta$, an algebraic point $x$ has height
\[
 h^\theta_{\mathcal L}(x)
   =\frac{(\overline x\cdot\mathcal L\cdot\theta_2^{g-1})}
           {[K(x):K]},
 \qquad
 h^\theta_{\mathcal L}(C_K)
   =\frac{(\mathcal L^2\cdot\theta_2^{g-1})}{2g},
\]
where $\overline x$ is its closure in $C\times J$.
The base polarization remains in both formulas.

The intersections can be checked with the same pullback identity.
Since $f_{m,n}$ takes values in the $g$-dimensional variety $J$,
$(f_{m,n}^*\theta)^{g+1}=0$.
Using $\alpha_1^2=0$ and
$(\alpha_1\cdot\theta_2^g)=(\theta^g)=g!$, the coefficient
of $m^2n^{2g}$ gives
\[
 (g+1)g\,g!+\binom{g+1}{2}(Q^2\cdot\theta_2^{g-1})=0,
 \qquad (Q^2\cdot\theta_2^{g-1})=-2g!.
\]
The restrictions of $Q$ to the two kinds of fibre are algebraically
trivial, so the other cross terms with one $Q$ vanish. Consequently
\[
 ((a\alpha_1+b\theta_2+cQ)\cdot\mathcal L\cdot\theta_2^{g-1})
                      =(a+gb-2c)g!,
 \qquad
 h^\theta_{\mathcal L}(C_K)=(g-1)(g-1)!.
\]

The effective divisor $\overline x$ has numerical class
$a\alpha_1+b\theta_2+cQ$ with $a=[K(x):K]>0$.
Set $s=b/a$ and $t=c/a$. The cone inequality gives $s\geq gt^2$,
and therefore
\[
 \frac{h^\theta_{\mathcal L}(x)}{g!}
       =1+gs-2t
       \geq1+g^2t^2-2t
       \geq1-\frac1{g^2}.
\]
This bound is attained by points of arbitrarily large degree.
Indeed, a smooth connected divisor $H_n$ in $|n\theta|$ for
sufficiently large $n$ pulls back under $f_{g,1}$ to an integral
divisor $E_n$. The automorphism
$(x,y)\mapsto(x,g i_\alpha(x)+y)$ identifies $E_n$ with
$C\times H_n$. Its generic closed point has degree $g^3n$,
and its numerical class is
$n(g^3\alpha_1+\theta_2+gQ)$, which attains the displayed bound.
These points form a generic sequence on the curve.
Thus \cite[Theorem~3.2]{GY26RelativeMinimaV2} gives
\[
 \mu^{\mathrm{ess}}(h^\theta_{\mathcal L})
   =\mu^{\mathrm{abs}}(h^\theta_{\mathcal L})
   =\left(g-\frac1g\right)(g-1)!
   >h^\theta_{\mathcal L}(C_K).
\]
For $g=2$ the two minima are $3/2$, while the curve height is $1$.
Hence the model-height inequality
$h(C_K)\geq\tfrac12(\mu^{\mathrm{ess}}+\mu^{\mathrm{abs}})$
fails over this higher-dimensional base, even for a nef model
line bundle with ample generic restriction.
The upper bound $\mu^{\mathrm{ess}}\geq h(C_K)$ is consistent
with these values. Transferring a two-sided height comparison
to function fields therefore requires the base and its polarization
to be specified.

\paragraph{The remaining directions.}
Returning to the number-field setting, the sufficient intrinsic condition
of \cite[Theorem~6.17]{BallaySombra26Published} is
\[
 \liminf_{t\uparrow\mu,\ t<\mu}
       \frac{\mu-t}{\rho(\overline D-t[\infty])}=0,
\]
where $\rho$ is the supremum of geometric inradii of adelic nef
approximants. The source's Question~6.21 asks whether
differentiability, or even all-place equidistribution, forces this
condition. Proposition~7.17 proves a converse under an additional
dominating semipositive metric satisfying Zhang equality;
Theorem~8.10 gives the semipositive toric equivalence. The general
necessity question remains open in the source.

There is a separate issue for logarithmically singular tests.
Weak convergence alone controls continuous tests. For an effective
divisor $E$, \cite[Theorem~7.13]{BallaySombra26Published} additionally
requires the directional derivative in $\overline E$ to equal the
sum of its limiting local Green integrals. Corollary~7.15 gives a
cycle-height criterion when all approximating local measures
eventually stabilize; removing that stabilization is Remark~7.16.
Neither issue is resolved by the roof computation above.

\section*{Notes and further reading}
\markright{Notes and further reading}

\enlargethispage{4\baselineskip}

\paragraph{The original theory.}
Differentiability along $\mathcal E^1$ segments is
\cite[Proposition~6.12]{BFJ15}. Orthogonality, differentiability of
$E_\omega\circ P_\omega$, and finite-energy projected variations are
\cite[Appendix~A, Theorem~7.2, and Corollary~7.3]{BFJ15}. The original
orthogonality theorem assumes algebraizability. Its removal in residue
characteristic zero is \cite[Theorem~6.3.3]{BGGJKM20}.  Continuity of
envelopes, used throughout and audited in
\Cref{sec:ch27-where-continuity}, is \cite[Theorem~8.3]{BFJ16}, proved in
\Cref{chap:bfj-envelopes}.  The variational functional
\eqref{eq:bfj-variational-functional} and the Euler--Lagrange statement
\Cref{prop:bfj-euler-lagrange} are \cite[\S7]{BFJ15}; the existence theory
that completes them is \cite[\S8]{BFJ15} and occupies
\Cref{chap:bfj-finite-energy-measures}.

\paragraph{The complex source of the method.}
The architecture of this chapter---concavity of the energy, projection by an
envelope, orthogonality, differentiability of $E\circ P$, and the
resulting Euler--Lagrange equation---is that of the variational approach to
complex Monge--Amp\`ere equations of Berman--Boucksom--Guedj--Zeriahi
\cite{BBGZ13}, itself resting on the equilibrium-energy differentiability of
Berman--Boucksom \cite{BB10}.  In the complex case orthogonality is
available by a Bedford--Taylor argument on the contact set, and it is not
the bottleneck; the whole difficulty of the non-Archimedean theory is that
no such local argument exists, which is why
\Cref{rem:model-slackness-heuristic} and
\Cref{rem:ch23-model-slackness-fails} are stated at such length.  Finite
energy in big classes, the setting in which the complex statements are
usually quoted, is \cite{BEGZ10}; the corresponding Monge--Amp\`ere and
stability theory is \cite{BBEGZ16}.

\paragraph{The presentational choices made here.}
Five things are organized differently from the source.

First, the superdifferential of $E_\omega\circ P_\omega$ is introduced
before orthogonality, in \Cref{sec:ch27-superdifferential}, and the soft
facts---nonemptiness, compactness and mass $V$---use standard functional
analysis.  The contact-set conclusion additionally uses envelope continuity.  This isolates precisely what
orthogonality adds, namely the identification of one supergradient, and it
makes \Cref{ex:ch27-corner-without-orthogonality} available as a warning
about one-sided estimates.

Second, the two-sided sandwich \eqref{eq:ch27-two-sided-sandwich} is stated
for an arbitrary pair of finite-energy potentials, with no ordering
hypothesis, consistently with the unconditional form now proved in
\Cref{prop:ch26-e1-cocycle}.  Both inequalities are needed in
\Cref{thm:ch27-model-direction-differentiability}, applied to a pair that is
not comparable, and deriving them from
\Cref{lem:ch27-concave-quotients} rather than from monotonicity of mixed
measures is what makes that possible.

Third, \Cref{thm:e1-energy-differentiability} is proved through the exact
polynomial structure of \Cref{prop:ch27-energy-polynomial} and Lagrange
interpolation, rather than by an approximation estimate.  The point is that
a decreasing limit of concave functions need not have convergent
derivatives; a decreasing limit of polynomials of bounded degree does, and
the energy along a bounded psh segment is such a polynomial.  The
identification of the limit then rests on the two elementary measure
statements \eqref{eq:ch27-diagonal-limit} and
\eqref{eq:ch27-mixed-limit}, both proved from the mass tail
\eqref{eq:e1-mass-tail} and the setwise monotonicity in
\Cref{def:nonpluripolar-ma}.

Fourth, differentiability along model directions and along continuous
directions are separated, and the second is deduced from the first by an
argument that uses no orthogonality
(\Cref{rem:ch27-two-proofs}).  The direct argument would prove both at once;
the separation records that the passage from model to continuous data is
soft, which is useful when orthogonality arrives first for model obstacles
from an intersection-theoretic argument.

Fifth, the convex case is stated as a theorem with a complete proof,
\Cref{thm:ch27-legendre-orthogonality}, in which orthogonality and the
derivative formula are two readings of the same Danskin computation
(\Cref{rem:ch27-one-computation}).  \Cref{chap:bfj-envelopes} proves the
convex orthogonality statement by the nondifferentiability of the conjugate,
\Cref{thm:ch23-convex-envelope-orthogonality}; the two arguments are the
same duality seen from opposite sides, and both are recorded because the
present one also yields \eqref{eq:ch27-convex-derivative}.

\paragraph{Nets rather than sequences.}
\Cref{lem:ch27-concave-nets}\,(3) replaces dominated convergence by
monotonicity in the integration variable.  It is used twice, but for two
different reasons.  In \Cref{thm:energy-envelope-differentiability} the
model data are chosen as a sequence, so dominated convergence would work as
well; using the same lemma keeps the limiting argument uniform.  In
\Cref{cor:e1-projected-variation}, the chosen regularization is a net whose
index set need not be countably cofinal, so the lemma justifies that proof.
One could instead choose a sequence for the single potential and use
dominated convergence, as in \cite[Corollary~7.3]{BFJ15}; this is not a
countably cofinal choice of models.  The relevant counterexamples to
sequential shortcuts are \Cref{ex:no-cofinal-sequence} and
\Cref{ex:decreasing-net-not-sequential}; the general policy is that of
\Cref{rem:net-vs-sequence-checklist}.

\paragraph{Asymptotic orthogonality.}
The mechanism sketched in \Cref{sec:ch27-orthogonality-mechanism} rests on
approximate Zariski decompositions.  Its absolute counterparts are
differentiability of the volume function \cite{BFJ09}, restricted volumes
and their asymptotic orthogonality \cite{ELMNP09}, Okounkov bodies
\cite{LM09}, and the movable-cone duality \cite{BDPP13}; the relative
statement is \cite[Appendix~A]{BFJ15}, where algebraizability is used
exactly once, to spread a model over a projective curve.  Non-Archimedean
volumes of metrized nef line bundles, the functional whose
differentiability replaces this argument in the later route, are developed
in \cite{BGM21}.  In \cite{BGGJKM20}, with an appendix by Lazarsfeld, the
differentiability theorem is Theorem~5.4.3, the passage from a continuous
semipositive envelope to metric orthogonality is Theorem~6.3.2, and the
residue-characteristic-zero BFJ form is Theorem~6.3.3.

\paragraph{Positive residue characteristic.}
The continuity theorems quoted in \Cref{sec:ch27-positive-characteristic}
are Gubler--Jell--K\"unnemann--Martin \cite{GJKM19}: curves over an
arbitrary complete nontrivially valued field with an ample line bundle,
\cite[Theorem~1.1]{GJKM19}; equal characteristic $p$ with geometric origin and
resolution, \cite[Theorem~8.2]{GJKM19}; the unconditional surface case,
\cite[Corollary~8.3]{GJKM19}.  Orthogonality in equal characteristic $p$
is \cite[Theorem~9.2]{GJKM19}, which uses the metric orthogonality theorem
\cite[Theorem~6.3.2]{BGGJKM20}; the full chain is reconstructed in
\Cref{thm:gjkm-positive-characteristic-cy}.  Mixed characteristic is
\cite{FGK25}, recorded as
\Cref{thm:fgk-mixed-characteristic-cy}.  None of these removes the
dependence of the whole calculus on continuity of envelopes; they replace
its proof.

\paragraph{Toric and convex sources.}
The toric dictionary used in \Cref{sec:ch27-toric-transport} is that of
\Cref{chap:cld-examples}, following Burgos Gil--Philippon--Sombra
\cite{BGPS14} for the arithmetic toric formalism and Burgos
Gil--Gubler--Jell--K\"unnemann \cite{BGJKK25} for the tropical toric
pluripotential theory; the comparison of the real and non-Archimedean
Monge--Amp\`ere operators is \cite{Vil20}.  The convex analysis is standard
and is collected in \Cref{app:convexity}; the classical reference is
\cite{Roc70}.  The identity \eqref{eq:ch27-toric-energy}, that the energy is
affine in the roof variable, is the toric shadow of the fact that the
Legendre transform linearizes the Monge--Amp\`ere energy; it is used again
in \Cref{chap:bj-energy} in the trivially valued theory, where the roof
variable becomes a concave function on an Okounkov body.

\paragraph{Measures of finite energy.}
The dual functional $E^\vee_\omega$ of
\Cref{prop:ch27-conjugate-identification} is the object completed in
\cite{BJ25FiniteEnergy}; its use here is confined to identifying the
concave conjugate of $E_\omega\circ P_\omega$, and no completion theorem
is imported.  \Cref{chap:bfj-finite-energy-measures} states exactly what is
and is not proved about that completion.

\paragraph{What is new here, and what is not.}
No theorem of this chapter is new.
\Cref{thm:ch27-superdifferential} and
\Cref{prop:ch27-supergradients-on-contact} are the standard convex-analytic
facts about concave Lipschitz functionals, written out because the
dependency audit of \Cref{sec:ch27-where-continuity} needs to know exactly
which of them use \Cref{thm:bfj-envelope-regularity}.
\Cref{thm:ch27-legendre-orthogonality} is classical convex analysis in the
form used by \cite{BGPS14} and \cite{BGJKK25}.
\Cref{prop:ch27-curve-shadow} is the linear-programming shadow of
\Cref{thm:ch27-model-direction-differentiability} and is stated only to
exhibit the mechanism.  The computations of \Cref{sec:ch27-worked} are
carried out in full because they are the test suite against which the sign
and mass conventions of \Cref{app:normalizations} are checked.
